\part{The Spaces Between}
% After \part{The Spaces Between}
\begin{quote}
  \small
  \textit{``The Ways Between do not measure miles. They measure choices. A wrong turn might lead to yesterday. A right turn might lead to a version of yourself who made a different decision -- and is still alive to regret it. The mist does not hide landmarks. It hides destinations.''}
  \noindent --- Dusana of the Raven Road, at a crossroads that did not exist on any map
  \end{quote}
% ------------------------------
% CHAPTER: UBRAL
% ------------------------------
\chapter{Ubral}
\emph{The Stone Between Spears.}

\noindent
{\Large
\textit{The stone here remembers every oath sworn upon it and every blood spilled in its shadow. To lead is to understand that the land itself is the true sovereign, and we are merely its stewards.}

\vspace{0.5em}

I am a Hearth-Aunt. Hospitality's a sacred thing in these hills---break bread at my table and you're family till the sun sets twice. But cross that trust and the very stones will remember your name.

\vspace{0.5em}

Ubral is a land of upland mists, stone cairns, and clan oaths held tighter than iron. Caught between the Aelerian mountains and the lowland courts of Viterra, its hills are scarred by old raids and crowned by watch-fires that speak faster than any rider. Here, law is written not in charters but in cattle, wergild, and songs; a guest's word can buy more than a sword, but one broken vow can spark a feud lasting generations. But here's the secret the Lady of Tor doesn't want you to know: the cairns are not just markers. They are witnesses. Every stone in every cairn has a name. If you touch the right stone, you can hear the dead voting on the disputes of the living. The dead have long memories. They vote with the wind.

\vspace{0.5em}

Viterra has tried to conquer Ubral a dozen times; each attempt ended with supply lines cut by reivers, passes sealed by Aeler outcast engineers, and lowland knights lost in mist. The Utaran Imperium itself broke its leg on these hills and limped away. Now the clans watch the lowlands with patient eyes, and the exiled Aeler of Khaz-Vurim weigh every traveler's honesty on iron scales. A guest here is kin---until they prove otherwise. And kin here will die for you---or kill you over a stolen cow.

\vspace{0.5em}

Never refuse a cup from a hearth-aunt. Hospitality is a contract sealed in salt and bread. Refusal is an insult that blood must answer.

\vspace{1em}

\noindent\textit{--- Hearth-Aunt Brigid, who has never drawn a blade and has never needed to}
}

\vspace{1em}

\begin{almanac}
\textbf{What do people eat for breakfast?} Oatcakes and skyr---the oats are grown in highland clearings, the skyr is cultured in stone crocks buried in the peat. The peat is ancient. The peat remembers the old kings. The old kings are buried beneath it.

\textbf{What does a child learn before they can walk?} To count cattle. Ubral children are taught the eight breeds of the hills: Black, Red, White, Brindled, Dun, Gray, Spotted, and the ninth---the breed that has no color because it belongs to the cairn-folk. The cairn-folk are the dead. The dead do not pay taxes.

\textbf{What is the first thing you notice about a stranger?} Their guest-token. Every Ubral traveler carries a token---a stone, a bone, a braid of hair---given by their last host. A stranger whose token is warm has been recently welcomed. A stranger whose token is cold has been traveling long. A stranger with no token is a spy. Spies are given a token---a rope. The rope is for their neck.

\textbf{What is the most common crime?} Cattle raiding---but raiding is not a crime, it is a sport. The cattle are the wealth. The wealth is the clan's. A stolen cow is a stolen child. The child must be returned. The return is a negotiation. The negotiation is a feud. The feud is a song.

\textbf{What does no one talk about but everyone knows?} The Lady of Tor is not the ruler. The cairns rule. The cairns vote at each full moon. Their votes are carried by the wind. The wind speaks to the hearth-aunts. The hearth-aunts speak to the clan. The clan obeys. The Lady is the voice. The cairns are the vote.

\textbf{Where do people go when they die?} To the cairns. The dead are not buried. Their bones are added to the cairns. The cairns grow taller each year. The tallest cairn is the oldest. The oldest cairn has no name. The oldest cairn remembers the first clan.

\textbf{How do you apologize for a serious wrong?} You offer a wergild---a payment in cattle, calculated by the wronged party's estimate of their loss. The cattle are driven to the cairn. The cairn is the court. If the cattle graze peacefully, you are forgiven. If they scatter, you are not. If they refuse to enter the cairn's shadow, your debt is blood.

\textbf{What is the most important festival?} The Gathering of the Clans (autumn). The clans meet at the Great Cairn. The cattle are counted. The songs are sung. The feuds are settled. The dead are named. The wind carries the names to the hills. The hills remember.

\textbf{What is the secret that could destroy a powerful person?} The Great Cairn is not a cairn. It is a prison---the tomb of a Viterran king who tried to conquer Ubral and failed. He was buried alive, with his crown and his sword. His ghost votes with the cairns. He votes for war. The cairns have not yet agreed.
\end{almanac}

\newpage
\section{Ubral --- ``The Stone Between Spears''}
\label{chap:ubral}

\subsection*{Quick Reference}
\textbf{Tagline:} \textit{The lowlands write their laws in ink. We carve ours in stone. The ink fades.}

\begin{quote}
  \textbf{Lady of Tor (Elite):}  
  ``The stone here remembers every oath sworn upon it and every blood spilled in its shadow. To lead is to understand that the land itself is the true sovereign, and we are merely its stewards. The lowland kings write their claims on parchment. We write ours in cairns.''
\end{quote}

\begin{quote}
  \textbf{Hearth-Aunt (Commoner):}  
  ``Hospitality's a sacred thing in these hills -- break bread at my table and you're family till the sun sets twice. But cross that trust and the very stones will remember your name. The Viterran tax men learned that. Their bones are under the sheepwalk now.''
\end{quote}

\begin{tcolorbox}[colback=black!3,colframe=blue!60!black,title={Theme \& Atmosphere}]
Ubral is a land of knife‑edged ridges, mist‑veiled glens, and stone cairns that mark the bones of old kings. Caught between the Aelerian mountains to the north and the lowland courts of Viterra to the south, its hills are scarred by centuries of raiding and reprisal. The clans speak a tongue thick with swallowed consonants and older grudges. They measure wealth in cattle, honour in wergild, and time in the generations since a Viterran banner last burned in the high passes. Exiled Aeler — dwarves who could not abide the Deep Vaults' suffocating tradition — have found refuge here, sharing their cunning with forge and toll‑road, and their outcast bitterness with any clan willing to listen. The lowland lords of Viterra claim suzerainty over Ubral; the clans call those claims ``ink on wet stone.'' Every pass holds a watch‑fire, every ford a memory of a battle, every hearth a knife waiting for a guest who stays too long.

\vspace{2mm}
\textbf{What you notice first:} The smell of peat smoke and wet granite. The constant lowing of cattle on distant hills. The way every threshold has a knife hidden nearby---for guests, not against them.

\textbf{The one rule that kills newcomers:} Never refuse a cup from a hearth‑aunt. Hospitality is a contract sealed in salt and bread. Refusal is an insult that blood must answer.
\end{tcolorbox}

\subsection*{Regional Mood: Cairn, Tor, and Oath}

Here, the past is never past. Every cairn holds a voice, every burn scar on a hillside remembers a raid. Viterra has tried to conquer Ubral a dozen times; each attempt ended with supply lines cut by reivers, passes sealed by Aeler outcast engineers, and lowland knights lost in mist. The Utaran Imperium itself broke its leg on these hills and limped away. Now the clans watch the lowlands with patient eyes, and the exiled Aeler of Khaz‑Vurim weigh every traveler's honesty on iron scales. A guest here is kin---until they prove otherwise. And kin here will die for you---or kill you over a stolen cow. The Lady of Tor rules by consent of the Council of Cairns, and the Council speaks for the dead as much as the living. Every nine years, a moot is called at Grey Tor to settle the great feud — a dispute between the two oldest clans that no one can remember the origin of, but everyone remembers the blood.

\textbf{Starting Location:} A Warden's Cairn at the fog line. Spears planted like grave markers, their shafts wrapped in faded clan ribbons. Hoofbeats echo in the valley below. The Hearth‑Aunt offers a cup of small beer and a piece of oatcake. ``The Wolf Road wakes,'' she says. ``Viterran tax men are asking which passes you used. The clans are counting your steps. Speak true, or the cairn will add your name to its tally. And mind the Aeler outcast in the smithy — he's been here forty years and still forgets to duck under doorframes.''

\subsection*{The Clans of Ubral (A Highland Primer)}
\label{subsec:ubral-clans}
\begin{tcolorbox}[colback=black!3,colframe=blue!60!black,title={The Great Houses of the Stone}]
\small
\begin{tabular}{|p{0.18\textwidth}|p{0.18\textwidth}|p{0.22\textwidth}|p{0.22\textwidth}|}
\hline
\textbf{Clan} & \textbf{Homeland} & \textbf{Feuds (Ancient)} & \textbf{Relationship with Viterra} \\
\hline
Mac Tern & Grey Tor, upper Glen Tern & O'Creagh (blood feud) & Defiant; has never paid a lowland tax \\
\hline
O'Creagh & Black Broom Bog, eastern fells & Mac Tern, Mac Gowan & Wary; accepts bribes but no garrisons \\
\hline
Mac Gowan & Khaz-Vurim Steps, western passes & O'Creagh & Pragmatic; trades iron for lowland silver \\
\hline
O'Ruarc & Bride's Causey, southern marches & None active & Collaborator; factor in Valora is a cousin \\
\hline
Mac Taggart & Three-Fires Ridge, northern heights & Mac Tern (old) & Neutral; sells beacon service to both sides \\
\hline
O'Cahan & Pass of Ashes, far west & None & Isolationist; speaks to lowlanders only through Aeler outcasts \\
\hline
\end{tabular}
\smallskip

\textit{The Mac Tern – O'Creagh feud has outlasted every lowland king. Each has a cairn for every generation lost. No wergild has ever settled it; no marriage has ever bridged it. The Lady of Tor keeps the peace by reminding both that a Viterran judge would hang them all.}
\end{tcolorbox}

\subsection*{Faction Clocks (Campaign Engines)}
\label{subsec:ubral-clocks}
\begin{itemize}
\item \textbf{The Viterran Yoke [6]} – Lowland crown's creeping control. Advances when a chief accepts a Viterran title, the party enforces lowland law, or a garrison is established. Retreats when a tax collector is sent home empty (or headless) or a claim is exposed as fraud. On full: a Queen's Justiciar arrives with a circuit court; lowland law becomes default.
\item \textbf{Tor's Hold [6]} – Lady Morwen's personal authority. Advances when she wins a moot, brokers peace, or gains an ally. Retreats when she loses a feud or a client clan defects. On 0: she is deposed; civil war follows.
\item \textbf{Ancestral Grudge [8]} – The dead's patience. Advances when oaths are broken, wergild unpaid, or cairns defiled. On full: a cairn‑wight (see \textit{Witnessed Prey}) rises and demands blood.
\item \textbf{The Forge Ledger [6]} – Khaz‑Vurim's economic leverage. Advances when they secure toll contracts, sell iron, or gain a clan's debt. On full: they can close a pass to enforce terms.
\item \textbf{The Deep Road [8]} – Aeler outcasts' secret ambition. Advances as they survey, bribe, or blast. On full: construction begins on a dwarf under‑way that would destroy several clan cairns.
\item \textbf{Mac Tern's Revenge [6]} – Advances when Mac Tern suffers an insult or loses cattle. On full: they launch a devastating raid.
\item \textbf{O'Creagh's Patience [6]} – Retreats when O'Creagh is wronged. When it hits 0, they strike without warning.
\end{itemize}

\subsection*{Wergild Price Table (Quick Reference)}
\begin{tcolorbox}[colback=black!3,colframe=black!40!white,title={Blood Prices (in cattle, unless noted)}]
\small
\begin{tabular}{|l|l|}
\hline
Insult (spoken) & 1 cow \\
\hline
Stolen sheep & 3 cows \\
\hline
Broken betrothal & 5 cows + a horse \\
\hline
Wounding & 7–15 cows (by severity) \\
\hline
Killing a clansman & 21 cows + a public apology \\
\hline
Killing a chief & 42 cows + a hostage \\
\hline
\end{tabular}
\end{tcolorbox}

\subsection*{The Hollow's Cairn (Legendary Site)}
Hidden in the high mist, a ninth stone circle with no clan markings. The dead do not vote here; they \textit{listen}. Once per arc, a character may climb the cairn and ask one question of the Hollow. The answer is always a price.

\paragraph{Mechanics:}
Spirit + Resolve (DV 6). Success: the GM answers truthfully (one sentence). Failure: the character gains the Condition \textit{Marked by the Hollow} (the GM may spend 2 SB to have a Debt‑Walker appear).

\subsection*{Starting Hooks}
\begin{itemize}
\item \textbf{The Wolf Road wakes:} Viterran tax men are asking which passes you used. The clans are counting your steps.
\item \textbf{The Bride‑Price Default:} A Mac Tern groom cannot pay the wergild for an O'Creagh bride. The bride has fled to the Pass of Ashes. Retrieve her – or pay her price yourself.
\item \textbf{The Toll‑Master's Ledger:} Khaz‑Vurim has doubled tolls on the Steps. A clan chief suspects they are skimming to fund the Deep Road. Infiltrate the counting‑house and copy the ledger.
\item \textbf{The Cairn that Weeps:} A shepherd reports that a cairn on Three‑Fires Ridge is weeping water that tastes of salt. The dead are agitated. Find out why before the clan blames a rival.
\end{itemize}

\subsection*{Faction Turn Cheat Sheet}
Between sessions, roll 1d6 for each major faction. Write one sentence of fallout and adjust the relevant clock.

\begin{center}
\begin{tabular}{|c|c|c|}
\hline
\textbf{Roll} & \textbf{Progress} & \textbf{Clock Change} \\
\hline
1–2 & Setback & –1 segment (min 0) \\
3–4 & Stalled & No change \\
5–6 & Advance & +1 segment \\
\hline
\end{tabular}
\end{center}

\paragraph*{Example:}
Mac Tern raided O'Creagh's cattle (\textbf{Mac Tern's Revenge} +1). The Viterran factor raised tariffs (\textbf{The Viterran Yoke} +1). Khaz‑Vurim surveyed a new pass (\textbf{The Deep Road} +1).

\subsection*{Prominent NPCs of Ubral}
\label{sec:ubral-npcs}
\begin{longtable}{|p{0.2\textwidth}|p{0.25\textwidth}|p{0.3\textwidth}|p{0.15\textwidth}|}
\hline
\textbf{Name} & \textbf{Role/Title} & \textbf{Motivation} & \textbf{Rumoured Quirk} \\
\hline
Lady Morwen & Clan‑chief of the Grey Tor & Keep the old oaths and resist Viterran encroachment & She lost a hand to a Viterran knight; her iron hook has a built‑in bell. \\
\hline
Hearth‑Aunt Brigid & Keeper of guest‑right (commoner exemplar) & Protect her household and settle feuds before blood is spilled & She has never drawn a blade. She has never needed to. \\
\hline
Councillor Gethin & Elder of the Council of Cairns & Preserve the wergild system and the memory of the dead & He speaks once a year; his sentences take a full season to deliver. \\
\hline
Oath‑Singer Lilias & Verse‑weaver and feud mediator & Turn vengeance into cattle, not corpses & Her tongue is forked—she can praise and curse in the same breath. \\
\hline
Vent‑Prior Orla & Aeler outcast leader of Khaz‑Vurim & Build a new under‑way through Ubral, linking mountains to lowlands & She carries a bell that rings once for each miner she left behind. \\
\hline
Sister Niamh & Healer of the Cairn-Keepers of the Grey Palm (Temple of Light) & Run a free hospice at Grey Tor without angering the old gods & She asks no payment, but she writes every patient's name on a slate that never fills. \\
\hline
\end{longtable}

\subsection*{Names of Ubral}
\label{sec:ubral-names}
\begin{longtable}{|p{0.25\textwidth}|p{0.25\textwidth}|p{0.2\textwidth}|p{0.2\textwidth}|}
\hline
\textbf{First Names (Male)} & \textbf{First Names (Female)} & \textbf{Surnames} & \textbf{Epithets} \\
\hline
Gethin & Morwen & ap Cairn & \textit{the Grey} \\
\hline
Bran & Brigid & ferch Tor & \textit{Oath‑Sworn} \\
\hline
Drystan & Lilias & mac Tern & \textit{the Stone‑Foot} \\
\hline
Eilian & Niamh & O'Ruarc & \textit{Hollow‑Tongue} \\
\hline
Padraig & Siorcha & Mac Gowan & \textit{the Unbroken} \\
\hline
Tegid & Ailbhe & O'Creagh & \textit{Ash‑Finger} \\
\hline
\end{longtable}

\subsection*{Spades --- Places}
\label{sec:ubral-places}
\begin{enumerate}
\setcounter{enumi}{1}
\item \textbf{Sheepwalk Ledge} --- Goat path with room for one honest lie at a time. `[THRESHOLD]`
  \hfill \textit{The lie must be small---a false weight, a borrowed name. Tell it twice and the ledge crumbles. A Viterran spy learned that the hard way.}
\item \textbf{Warden's Cairn} --- Windy tor; signal‑fire basket and dry cache; eyes of the hills. `[WATCH]`
  \hfill \textit{The cache holds a single arrow. Light it with the fire, and the message is war. The arrow is fletched with eagle feathers — a clan's highest oath.}
\item \textbf{Wergild Ford} --- Flat stones, deep pools, table rock for counting silver. `[LAW]` `[TRADE]`
  \hfill \textit{The table rock is smooth as glass. Blood washes off; lies do not. The water remembers every price ever paid here.}
\item \textbf{Droppers' Bridge} --- Stone span rigged to fall; pins already loosened. `[TRAP]`
  \hfill \textit{The pins are greased with mutton fat. The bridge knows who paid for the grease. Its fall has ended three feuds and started two.}
\item \textbf{Scree-Ladder} --- Climbing to a notch; red rags mark safe steps through danger. `[TRAVEL]`
  \hfill \textit{The rags are old clan banners. Step where they wave, not where they lie flat. One rag is white — that marks a traitor's grave.}
\item \textbf{Moot Hollow} --- Ring of standing stones; voices carry and won't quite stop. `[SOCIAL]` `[OATH]`
  \hfill \textit{The stones have ears. Speak a false oath here, and every cairn will know. The dead vote by how the mist moves.}
\item \textbf{Reiver's Gate} --- Between boulders; cart‑ruts vanish into heather's memory. `[STEALTH]`
  \hfill \textit{The gate has no lintel. Above it, ravens nest in a skull. The ravens are the clan's spies.}
\item \textbf{Khaz-Vurim Steps} --- Dwarf‑cut switchbacks with iron mile studs; toll and toil. `[TRAVEL]` `[AUTHORITY]`
  \hfill \textit{The studs are cold iron. Touch them, and the mountain knows your weight. Aeler outcasts maintain the toll — they accept stories as often as coin.}
\item \textbf{Grey Tor Fort} --- Earthen rampart, timber crown, smoky cook pits. `[COMMUNITY]`
  \hfill \textit{The cook pits have been burning for three hundred years. The ash is deep enough to bury a secret — and more than one Viterran scout.}
\item[J] \textbf{Black Broom Bog} --- Stepping‑logs; wrong one drinks you into silence. `[DEATH]` `[FEAR]`
  \hfill \textit{The logs are painted with warding runes. One is unpainted — that one is the quick way. The bog exhales names of the drowned.}
\item[Q] \textbf{Bride's Causey} --- Raised road to a valley kirk; ribbons hang like warnings. `[SOCIAL]`
  \hfill \textit{The ribbons are from wedding gowns — cut short when the groom fell in a feud. The kirk is neutral ground, but the path is not.}
\item[K] \textbf{Three-Fires Ridge} --- Watchposts that see Viterra and Vhasia both; news travels fast. `[WATCH]`
  \hfill \textit{The fires are layered: birch for peace, pine for raid, oak for invasion. A fourth fire — yew — means the Hollow walks.}
\item[A] \textbf{Pass of Ashes} --- When snow closes here, the upland becomes an island. `[WEATHER]` `[DEATH]`
  \hfill \textit{The ash is from old sacrifices. The pass remembers every offering. The last sacrifice was a Viterran herald. His voice is still in the wind.}
\end{enumerate}

\subsection*{Hearts --- People \& Factions}
\label{sec:ubral-people}
\begin{enumerate}
\setcounter{enumi}{1}
\item \textbf{Hearth-Aunt} --- Holds guest‑cup and house's temper; hospitality as power.
  \hfill \textit{She has never drawn a blade. She has never needed to. Her hospitality has ended more feuds than swords.}
\item \textbf{Hill Guide} --- Thorn‑staff and ten quiet shortcuts; the land as a map.
  \hfill \textit{He taps each stone before stepping. The stones talk to him. He learned the tongue from an Aeler outcast who went deaf and learned to listen with his feet.}
\item \textbf{Feud-Broker} --- Knows the weight of a life in cattle and coin; peace as trade.
  \hfill \textit{His ledger is a row of tally sticks. Each notch is a death---or a birth. He has notched both sides of the same feud for forty years.}
\item \textbf{Reiver Band} --- Light on tack, heavy on nerve, laughing in the rain.
  \hfill \textit{They wear no colors. Their horses answer to whistles, not names. They raid Viterran supply trains and leave heather sprigs as signatures.}
\item \textbf{Fire Warden} --- Braziers speak faster than riders; warning as weapon.
  \hfill \textit{She can light a signal with a single word. The word is a dead man's name. The dead man was a Viterran knight who burned her village.}
\item \textbf{Wergild Counter} --- Keeps tallies, ends grudges, starts others; law as math.
  \hfill \textit{His abacus beads are sheep knuckles. He clicks them in his sleep. His grandmother taught him; she learned from an Aeler vent‑prior.}
\item \textbf{Dwarf Warden (Aeler Outcast)} --- Khaz‑Vurim road; toll first, friendship later.
  \hfill \textit{Her hammer has a chip where she cracked a liar's skull. She never fixed it. She came to the surface for the sky — and stayed for the feuds.}
\item \textbf{Oath-Singer} --- Verses bind hands and open gates; word as iron.
  \hfill \textit{Her tongue is forked---she can praise and curse in the same breath. Her songs are older than the Viterran crown.}
\item \textbf{Lowland Factor} --- Buying iron blooms, selling trouble; coin as catalyst.
  \hfill \textit{He carries two ledgers. One is for the clans. One is for the tax man. The third ledger is in his head; it lists every bribe he has ever taken.}
\item[J] \textbf{Bride-Carrier} --- Peaceweaver walking bloodlines with a knife and a smile.
  \hfill \textit{Her knife has a wooden blade. The threat is real. She has walked between warring clans nine times. The ninth was a Viterran lord who wanted a wife. He got a treaty instead.}
\item[Q] \textbf{Lady of Tor} --- Clan‑chief in cloak and mail; her nod is winter or spring.
  \hfill \textit{She lost a hand to a Viterran knight. Her iron hook has a built‑in bell. The bell rings when she lies. It has never rung.}
\item[K] \textbf{Council of Cairns} --- Greybeards and granite wills; they do not hurry.
  \hfill \textit{They have sentenced men to exile. The sentences take a year to deliver. The exile always knows before they speak.}
\item[A] \textbf{Stone-Speaker} --- Aeler outcast envoy whose word moves roads and rates.
  \hfill \textit{He carries a tuning fork. Strike it, and the mountain answers. He has not seen his homeland in sixty years. He calls the mist his ceiling now.}
\end{enumerate}

\subsection*{Clubs --- Complications/Threats}
\label{sec:ubral-complications}
\begin{enumerate}
\setcounter{enumi}{1}
\item \textbf{Upland Mist} --- Hear horns but not edges; navigation by sound and instinct. `[WEATHER]` `[STEALTH]`
  \hfill \textit{The horns are calling a muster. Which clan? You will know when they arrive. The mist smells of wet wool and old iron.}
\item \textbf{Feud Rekindled} --- Cousin spits on guest‑law; knives wake from sleep. `[SOCIAL]` `[COMBAT]`
  \hfill \textit{The spit lands at your feet. The insult is now yours to answer. The cousin's father was killed in a cattle raid thirty years ago. He has been waiting.}
\item \textbf{Bridge Dropped} --- Pursuers fall... and the route with them; passage becomes peril. `[TRAP]` `[TRAVEL]`
  \hfill \textit{The bridge was cut from the far side. Someone expected you. The cut ropes smell of dwarf‑made pitch — a signature of the Aeler outcasts.}
\item \textbf{Black-Rent Demand} --- ``Privateering on land''; pay or be ``escorted.'' `[AUTHORITY]` `[DEBT]`
  \hfill \textit{The escort has a noose. Choose your words carefully. The man collecting is a Viterran factor with a royal warrant. The warrant is a forgery. No one will admit it.}
\item \textbf{Wergild Breach} --- Silver short by a head; tempers long by a spear's length. `[LAW]` `[DEBT]`
  \hfill \textit{The bereaved clan has a new widow. She is sharpening a blade. The missing silver was paid in lowland coin — debased, clipped, and worthless. The insult is not the shortfall; it is the contempt.}
\item \textbf{Snow-Squall} --- Seals the notch; tents turn to coffins if you dally. `[WEATHER]` `[TRAVEL]`
  \hfill \textit{The snow is tinged grey---ash from the Pass. Someone burned an offering. The offering was a Viterran tax ledger. The gods accepted.}
\item \textbf{Dwarf Toll Hike} --- At the Steps; papers right, purses wrong; law costs coin. `[AUTHORITY]` `[DEBT]`
  \hfill \textit{The warden quotes a rate from a book that wasn't there yesterday. The book was written by an Aeler outcast who lost his sense of humour but not his arithmetic.}
\item \textbf{Cattle Scatter} --- Bells ringing downslope; the cover story runs with the herd. `[STEALTH]` `[DISASTER]`
  \hfill \textit{The bells are yours. The herd is not. The real owner is a clan chief with a long memory and three sons with sharp swords.}
\item \textbf{False Alarm} --- Watch‑fire beacons ridge to ridge; levies seize the road. `[WATCH]` `[AUTHORITY]`
  \hfill \textit{The fire was lit by a child. Or a spy. The levy doesn't care. The spy is a Viterran agent. The child is his nephew.}
\item[J] \textbf{Bride-Theft} --- A wedding becomes a war‑party; you are caught between both. `[SOCIAL]` `[COMBAT]`
  \hfill \textit{The bride's veil is on the ground. The groom's own kin drew first blood. The bride was promised to a neighbouring clan. She chose her own groom. Now three clans want her head — and her hand.}
\item[Q] \textbf{Royal Incursion} --- Neighbor's ``lawful'' arrests in the uplands; flags and chains. `[AUTHORITY]` `[DEBT]`
  \hfill \textit{The flags are Viterran. The chains are new. The prisoners are children. The crime is not paying a tax that was never law. The clans are calling the muster.}
\item[K] \textbf{Clan Muster} --- Horns call men from steadings to spears; all traffic stops. `[WATCH]` `[COMBAT]`
  \hfill \textit{The muster stone is black with old blood. The new blood will be red. The last muster was forty years ago, to drive out a Viterran garrison. The garrison's bones are under the sheepwalk.}
\item[A] \textbf{Hill-Fall} --- Rain liquefies a slope; trail, proof, and bodies slide together. `[WEATHER]` `[DEATH]`
  \hfill \textit{The mud is warm. Something underneath is still breathing. The slide exposed a mass grave — clan warriors who fell in a forgotten feud. Their ghosts are counting the wergild still owed.}
\end{enumerate}

\subsection*{Diamonds --- Rewards/Leverage}
\label{sec:ubral-rewards}
\begin{enumerate}
\setcounter{enumi}{1}
\item \textbf{Guest-Token} --- One hearth owes food, bed, and steel at dawn. `[SOCIAL]`
  \hfill \textit{The token is a knucklebone from a feast. The host's name is carved on it. Break it, and the obligation ends. The host will not thank you.}
\item \textbf{Guide's Braid} --- Lawful passage on named sheepwalks (once). `[TRAVEL]`
  \hfill \textit{The braid is woven from horsehair and a promise. Unravel it to return. The promise was made to a hill guide who died last winter. His ghost still checks the knots.}
\item \textbf{Ford Remission} --- Cross Wergild Ford free for a season; water as ally. `[TRAVEL]` `[DEBT]`
  \hfill \textit{The remission is a white stone. Skip it across the ford to pay. The stone was taken from a Viterran courthouse. The law has not found it yet.}
\item \textbf{Feud-Charter} --- Two clans sheath blades until the next harvest. `[SOCIAL]` `[OATH]`
  \hfill \textit{The charter is a broken spear. Each clan keeps a half. The spear was used to kill a Viterran tax collector. His family still wants it back.}
\item \textbf{Bloom Allotment} --- Claim on a week's iron from a hill bloomery. `[TRADE]`
  \hfill \textit{The bloomery is a woman's name. She decides how much you get. She is an Aeler outcast's granddaughter. Her iron is harder than any lowlander's steel.}
\item \textbf{Watch-Code} --- Today's beacon order from Three‑Fires Ridge. `[WATCH]` `[TRAVEL]`
  \hfill \textit{The code is a sequence of smoke puffs. Read it, and the hills will watch for you. The code changes weekly. The fire wardens burn the old codes — the smoke carries the memory.}
\item \textbf{Pass-Ring} --- Dwarf road priority for one train of carts. `[TRAVEL]` `[AUTHORITY]`
  \hfill \textit{The ring is iron. Wear it on your off hand. The Aeler outcasts who made it stamped their names inside. One name is scratched out — the smith was killed in a feud. The ring still works.}
\item \textbf{Bride-Escrow} --- Hold the purse; both sides must humor you. `[SOCIAL]` `[DEBT]`
  \hfill \textit{The purse is a leather bag. The coins inside are from a dead man's hoard. The dead man was a Viterran lord who tried to tax the clans. His ghost asks for his money back.}
\item \textbf{Shelter Writ} --- Grey Tor opens its gates during storm or pursuit. `[SOCIAL]`
  \hfill \textit{The writ is a raven feather. Burn it to claim shelter. The raven was a clan spy. His memory still watches the gates.}
\item[J] \textbf{Oath-Bracelet} --- One binding vow ends cleanly, witnessed. `[OATH]` `[MAGIC]`
  \hfill \textit{The bracelet is silver. The clasp is a wolf's head. Bite it to break the oath. The wolf's head was a seal of a Viterran king. The oath was sworn under duress. The bracelet remembers.}
\item[Q] \textbf{Council Audience} --- The Cairns hear you alone; precedent sticks like stone. `[SOCIAL]` `[LAW]`
  \hfill \textit{The audience is held at midnight. The stones glow when they agree. The last audience exiled a clan chief. He is still walking the passes, looking for a cairn that will forgive him.}
\item[K] \textbf{Road Commission} --- Collect tolls on a stretch of pass (for now). `[AUTHORITY]` `[TRADE]`
  \hfill \textit{The commission is a wooden staff. Plant it at the toll gate. The staff was carved from a gallows tree. The last toll‑keeper was hanged. The staff still creaks.}
\item[A] \textbf{Stone-Clause} --- Temporary exception to dwarf toll or law, sealed. `[AUTHORITY]` `[OATH]`
  \hfill \textit{The clause is a single rune. Carve it on the toll post, and the mountain will honor it — once. The rune was taught to the clans by an Aeler outcast who wanted to see a Viterran knight pay double.}
\end{enumerate}

\subsection*{Quick use notes}
\label{sec:ubral-quick-use}
\begin{itemize}
\item Draw until you have all four suits: Spade = place, Heart = actor, Club = pressure, Diamond = leverage. Highest rank sets the main Clock (2--5 $\rightarrow$ 4, 6--10 $\rightarrow$ 6, J/Q/K $\rightarrow$ 8, A $\rightarrow$ 10).
\item Diamonds are codified outcomes (oaths/rights/tokens) that change position rather than call for a roll.
\item If any A appears, echo \textbf{stone and spear} motifs---cairns that watch, voices that echo, and oaths that cut deeper than steel.
\end{itemize}

\subsection*{Additional Features (Cultural Mechanics)}
\label{sec:ubral-mechanics}
\begin{itemize}
\item \textbf{Guest-Law:} Breaking hospitality stains a name worse than breaking a sword. Guest‑right tokens guarantee shelter, but abuse invites a blood feud. A Viterran lord who broke guest‑law at Grey Tor had his name struck from the Council's cairn. His own clan disowned him.
\item \textbf{Wergild System:} Every injury has a price in cattle or coin. Wergild counters keep tallies that outlive the injured; debts pass to kin. The lowland courts have no equivalent; they imprison the poor and hang the desperate. The clans consider that barbaric.
\item \textbf{Dwarf Toll Roads:} Khaz‑Vurim gates charge for passage, but their roads are sure and swift. Pass‑rings grant priority, but dwarf law is strict. The Aeler outcasts who run the tolls are exiles — they fled the Deep Vaults because they could not abide the counting. Now they count for the clans, and they are gentle with the weary.
\item \textbf{Local Patrons (Syncretic Names):}
  \begin{itemize}
    \item \textbf{The Bone Keeper} — Varnek Karn. *The accountant of ancestral wergild; the witness to every blood price.*
    \item \textbf{The Iron Crown} — Mykkiel. *The law of the lowlands, cold and distant; the clans honour it by breaking it.*
    \item \textbf{The Grey Mare} — The Traveler. *The rider who never tires; the path that leads home through the mist.*
    \item \textbf{The Forge Father (Aeler outcast cult)} — The Sealed Gate. *The fire that shapes iron; the exile who builds a new hearth.*
  \end{itemize}
\end{itemize}

\subsection*{Lore Echoes (Cross‑Expansion Hooks)}
\begin{itemize}
  \item \textbf{The Bitter Root:} Daughters of the Cord run a hospice near the Pass of Ashes, tending to the wounded of both sides of the Great Feud. A renegade Sister has hidden a Hollowed in a cairn, trying to discharge a blood‑debt by feeding the Hollow memories of Viterran tax collectors. The clans do not know whether to thank her or burn her.
  \item \textbf{Malachai, the Chained Angel:} A silver coin that never spends appears in the palm of every third corpse pulled from a hill‑fall. Take it, and the dead will give you their last breath — but the coin will taste of peat smoke and regret.
  \item \textbf{The Grey Wanderer's Grimoire:} Ubral psions are called ``Stone‑Listeners.'' They press their foreheads to cairns and hear the names of everyone who has touched that stone. The Lady of Tor employs three; they have not spoken aloud in a decade.
  \item \textbf{The Threshold Compendium:} The Pass of Ashes is a Class II Anchor — the spirits of Viterran knights who froze to death in a blizzard still march there, looking for a way home. Saikou Ira left a silver mirror at the summit. The mirror shows the knights their own corpses. They have not advanced since.
  \item \textbf{The Ninth:} In any Ubral scene, the ninth stone in a cairn is always warm. Touch it, and you will remember an oath you never swore. Those who touch the ninth stone three times are said to become a cairn themselves — silent, still, and waiting for the next hand. The Aeler outcasts say the ninth stone is a door to the Deep Vaults they fled. They do not touch it.
\end{itemize}

\subsection*{Ubral Superstitions (burn for +1 die once)}
\begin{itemize}
  \item \textbf{Bread to the Cairn:} Crumble a crust at the base of a standing stone. Ignore the first \emph{Upland Mist} penalty. `[SUPERSTITION]`
  \item \textbf{Knot the Guest-Cup:} Tie a single knot in the handle of a drinking cup. Cancel \emph{Feud Rekindled} or take --1 die on your next hospitality roll (your choice). `[SUPERSTITION]`
  \item \textbf{Name to the Ford:} Whisper a dead enemy's name into the Wergild Ford. Gain +1 Effect when collecting from their kin this scene, but mark \emph{Obligation} to the Lady of Tor. `[SUPERSTITION]`
\end{itemize}

\subsection*{Thematic SB Spend Table}
\label{sec:ubral-sb}

\paragraph{Minor Complications (1 SB)}
\begin{itemize}
\item \textbf{Exposure:} Your actions draw unwanted attention from \textbf{clan guards or reiver scouts}.
\item \textbf{Noise:} Sounds of your actions alert nearby \textbf{watch‑fires or hill guides}.
\item \textbf{Trace:} Evidence of your passage marks your route for \textbf{trackers or feud‑kin}.
\item \textbf{Delay:} A brief but meaningful setback costs you \textbf{time or favorable weather}.
\item \textbf{Supply Strain:} Mark +1 segment on a relevant \textbf{resource clock}.
\end{itemize}

\paragraph{Moderate Setbacks (2 SB)}
\begin{itemize}
\item \textbf{Alarm Raised:} \textbf{Hill guide or clan chief} becomes aware and begins responding.
\item \textbf{Position Lost:} You lose advantageous ground/cover/stealth due to \textbf{mist or bridge collapse}.
\item \textbf{Foe Appears:} A \textbf{reiver band or feud‑kin} arrives on scene.
\item \textbf{Gear Trouble:} A piece of equipment becomes \textbf{Compromised/Neglected}.
\item \textbf{Lock/Barrier:} A simple obstacle now requires a test to overcome.
\end{itemize}

\paragraph{Serious Trouble (3 SB)}
\begin{itemize}
\item \textbf{Reinforcements:} Additional \textbf{clan warriors, reivers, or dwarf wardens} arrive.
\item \textbf{Key Gear Breaks:} A crucial tool/weapon becomes temporarily unusable.
\item \bf{Major Twist:} The situation fundamentally changes — \textbf{feud declared/bridge falls/snow closes pass}.
\item \textbf{Rail Tick:} Advance a relevant campaign/front clock by 1 segment.
\item \textbf{Condition Applied:} Mark \textbf{Fatigue 1/Harm 1/Condition} appropriate to fiction.
\end{itemize}

\paragraph{Major Turns (4+ SB)}
\begin{itemize}
\item \textbf{Trap Springs:} A prepared danger activates with full effect.
\item \textbf{Authority Arrival:} \textbf{Lady of Tor, Council of Cairns, or Stone-Speaker} intervenes.
\item \textbf{Scene Shift:} The environment changes dramatically — \textbf{mist thickens/bridge drops/hill falls}.
\item \textbf{Patron Omen:} Divine/arcane forces take notice — \textbf{omen appears/blessing lost/curse manifests}.
\item \textbf{Narrative Pivot:} The story takes an unexpected turn that reframes objectives.
\end{itemize}

\paragraph{Region-Specific SB Options}
\begin{itemize}
\item \textbf{Ubral (Stone \& Oath):} Cairns shift position, voices echo from empty air, oaths bind mid‑speech.
\item \textbf{Ubral (Hill Law):} Paths reroute without warning, tolls increase mid‑journey, guest‑rights expire.
\item \textbf{Ubral (Feud Culture):} Insults are overheard, weapons appear in sheaths, kin arrive unsummoned.
\end{itemize}

\subsection*{Pursuit Ladder (Near / Mid / Far)}
\begin{itemize}
\item \textbf{Advance/Withdraw (Wits + Navigation):} On a hit, shift one band. On a strong hit, also impose \emph{Bog Mist}.
\item \textbf{Cut the Bridge (Presence + Command):} On a hit, freeze range for one exchange; if a \textbf{Pass-Ring} is in your possession, +1 die.
\item \textbf{Weather the Squall (Resolve + Survival):} On a hit in \emph{Snow-Squall}, you pick who finds shelter; on a miss, the GM does.
\end{itemize}

\subsection*{Boss Seeds (Stone \& Spear)}
\begin{itemize}
\item \textbf{The Wolf-Road Captain (reiver‑king on a piebald mare)} --- \emph{Phases}: Ride‑Through $\rightarrow$ Torch the Fold $\rightarrow$ Laugh in Rain. \emph{Cracks}: present a \textbf{Watch-Code} or \textbf{Guide's Braid} to reduce his effect; strike a truce with a season's cattle.
\item \textbf{The Bride-Thief of Black Broom (mask of bog‑rush)} --- \emph{Phases}: Switch the Bride $\rightarrow$ Bog Embrace $\rightarrow$ Nettle Oath. \emph{Cracks}: play a \textbf{Bride-Charter} to freeze blades; a \textbf{Well-Blessing} nulls one bog move.
\item \textbf{Stone Auditor of Khaz-Vurim (Aeler outcast toll‑magistrate)} --- \emph{Phases}: Double the Toll $\rightarrow$ Weigh the Lie $\rightarrow$ Seal the Switchback. \emph{Cracks}: present two of \{\textbf{Pass-Ring}, \textbf{Stamped Load-List}, \textbf{Named Sponsor}\} to downgrade his authority.
\item \textbf{The Oath-Eater Wight (cairn‑bound revenant)} --- \emph{Phases}: Take the Word $\rightarrow$ Cold of Cairn $\rightarrow$ Bind Witness. \emph{Cracks}: \textbf{Oath-Bracelet} or a skald's stanza spoken over bread and blood cancels one move; repay an ancestral wergild.
\item \textbf{Grey-Tor Hedge-Witch (heather crown, ash staff)} --- \emph{Phases}: Weather‑Turn $\rightarrow$ Hag‑Stone Sight $\rightarrow$ Bind the Gate. \emph{Cracks}: gift a \textbf{Bloom Allotment} or a winter's salt; a \textbf{Council Audience} compels parley.
\end{itemize}

\subsection*{Legendary Sites \& Trials}
\begin{itemize}
\item \textbf{Sheepwalk Ledge Trial} --- Speak one honest lie to pass. On success, set Position one step safer for the next border scene; on failure, start \emph{Kin Arrive}.
\item \textbf{Wergild Ford Weighing} --- Cast silver on the table‑rock. If paid in full, clear 1 \emph{Feud-Smoke}; if short, \emph{Wergild Breach} enters play.
\item \textbf{Pass of Ashes Vigil} --- Hold a fire through the night storm. Success grants \textbf{Shelter Writ}; a miss summons the \emph{Oath-Eater Wight} at weak strength.
\end{itemize}

\subsection*{Extensions (Plug-in)}
\label{sec:ubral-extensions}
\begin{itemize}
  \item \textbf{Oath Dial (Honor \(\leftrightarrow\) Cunning):} Track how the party navigates Ubral norms.
  \begin{itemize}
    \item \emph{Honor High:} +1 Position when invoking guest‑right, keeping truce, or paying wergild promptly; first ambush/raid action each session starts one step worse.
    \item \emph{Cunning High:} +1 Effect on raids, evasions, and border tricks; first formal parley or oath‑bound request each session starts one step worse.
    \item Center the dial by restoring stolen cattle, hosting an enemy well, or submitting to a cairn ruling.
  \end{itemize}

  \item \textbf{Feud Math (fast settlement):} Lay three markers: \textbf{Blood}, \textbf{Beasts}, \textbf{Blame}.
  \begin{enumerate}
    \item Make a single \emph{Parley/Command/Study} roll per side. On a hit, move one marker to your side; great Effect moves two.
    \item Win any two to settle: \emph{Blood} = oath/hostage, \emph{Beasts} = wergild in cattle/coin, \emph{Blame} = public apology/verse.
    \item Miss: mark \emph{Feud-Smoke} [4]. At 4, the losing side names a lawful target (bridge, herd, guide).
  \end{enumerate}

  \item \textbf{Guest-Right Procedure:}
  \begin{itemize}
    \item Present a \textbf{Guest-Token} (or bread and salt). Host must choose: \emph{Shelter}, \emph{Escort at Dawn}, or \emph{Neutral Meal}.
    \item Breaking it creates \emph{Blood-Oath} [6] against the violator; any Ubral ally may spend 1 \Diamond{} to tick it twice.
    \item Abiding it clears 1 tick from \emph{Feud-Smoke} and grants +1d to the next social action in that steading.
  \end{itemize}

  \item \textbf{Watchfire Net (ridge signals):}
  \begin{itemize}
    \item Two clocks: \emph{Beat the Beacons} vs \emph{Ridge Aflame} (size by highest rank in your seed).
    \item Choose a lane each exchange: \emph{Low Fold} (slow, stealth), \emph{Sheepwalk} (tricky, quick), \emph{Road} (fast, obvious).
    \item Advantage: Low Fold\(\,\triangleright\,\)Road (conceal), Road\(\,\triangleright\,\)Sheepwalk (speed), Sheepwalk\(\,\triangleright\,\)Low Fold (shortcuts).
    \item \emph{False Alarm (1 SB):} light or douse a misleading brazier; shift one segment between the two clocks.
  \end{itemize}

  \item \textbf{Reiver Ride (mounted chase):}
  \begin{itemize}
    \item Track \emph{Remounts} (3): each push for speed trades 1 remount for +1d or +1 Effect. At 0, gain \emph{Blown Mounts} condition.
    \item \emph{Hag-Stone Cut}: pass a known gap (Sheepwalk Ledge, Droppers' Bridge) to force pursuers to roll or fall behind.
  \end{itemize}

  \item \textbf{Khaz-Vurim Protocol (Aeler Outcast Rules):}
  \begin{itemize}
    \item Present any two: \textbf{Pass-Ring}, \textbf{Stamped Load-List}, \textbf{Named Sponsor}. Two = safe passage; one = pay toll + scrutiny; none = delay and search.
    \item \emph{Stone-Clause} may override one dwarf law once; afterwards start \emph{Ledger Notice} [3] with the wardens (prices rise, patience falls).
    \item \emph{Tally Truth}: you may substitute exact weights/measures for charm or coin to improve Position by one step with dwarves.
    \item Aeler outcast smiths accept stories as payment for small repairs — the story must be about a betrayal. They are collecting evidence for a case they will never bring.
  \end{itemize}

  \item \textbf{Bride-Peace Interlude:}
  \begin{itemize}
    \item Play a \textbf{Bride-Charter} to suspend violence at a scene. Everyone sheaths---or earns \emph{Song Shame} [4]. When it fills, a skald memorializes the insult (ongoing -1d with that clan).
  \end{itemize}

  \item \textbf{Start-of-Scene Omens (1d6):}
  \begin{enumerate}
    \item \emph{Clear Ridge}: first signal test gains +1d.
    \item \emph{Lowing Herd}: +1d to tracking or herding; -1d to stealth near cattle.
    \item \emph{Wet Stone}: climbs start one step worse; mending gear is easier (free repair tick).
    \item \emph{Old Verse Remembered}: first oath/plea gains +1d if spoken in rhyme.
    \item \emph{Border Drums}: Viterra patrols near; papers safer, raiding harsher.
    \item \emph{Dwarf Toll Day}: prices firm; any \Diamond{} tied to roads counts double value at the Steps.
  \end{enumerate}

  \item \textbf{Plug‑in SB Conversions (Ubral Flavor):}
  \begin{itemize}
    \item 1 SB \(\rightarrow\) \emph{Mud and Heather}: footing worsens; next physical action --1d or mark \emph{Winded}.
    \item 2 SB \(\rightarrow\) \emph{Kin Arrive}: add a rival cousin/ally who changes the ask or splits the pot.
    \item 3+ SB \(\rightarrow\) \emph{Cairn Decrees}: a standing stone ``speaks'' (witnesses recall the law); swap one held \Diamond{} for a lesser right until you make amends.
  \end{itemize}

  \item \textbf{Reputation Echoes:}
  \begin{itemize}
    \item \emph{Guest-Faithful}: once/session downgrade a social backlash inside a steading; next time you refuse hospitality, mark \emph{Feud-Smoke} +1.
    \item \emph{Reiver-Favored}: +1 Effect to seize cattle/contraband on the move; city gates start one step worse Position.
    \item \emph{Stone-Friend}: +1d with dwarf wardens on roads; any false measure triggers \emph{Ledger Notice} +1.
  \end{itemize}
\end{itemize}

\begin{tcolorbox}[colback=black!3,colframe=black!40!white,title={Patronage \& Power}]
In Ubral, power flows through clan bonds, guest‑right, and the careful balance of honor and iron. Clan chiefs maintain authority through wergild, moots, and the respect of their people, while the Aeler outcasts of Khaz‑Vurim wield influence through their control of mountain passes and forges. Law is personal and passed down through generations, with each oath and debt creating a web of obligation that can last centuries. The Viterran crown claims suzerainty; the clans call that claim ``the ink of thieves.''

\textbf{For the GM:}  
Patronage in Ubral revolves around hospitality, protection, and the resolution of feuds. Rewards often take the form of tokens, charters, or safe passage that can be leveraged into greater influence. To emphasize this, use the faction clocks, let rival clans issue conflicting protections, and remember the Aeler outcasts — they are exiles, not tourists. They will help a clan for a story, but they will never forget a slight.

In Ubral, your word is your bond, and your bond determines whether you walk in peace or peril. The lowlands write their laws in ink; Ubral carves them in stone. The ink fades. The stone remembers.
\end{tcolorbox}
\newpage

% ------------------------------
% CHAPTER: SILKSTRAND
% ------------------------------
\chapter{Silkstrand}
\emph{City of Bridges \& Dyewater.}

\noindent
{\Large
\textit{In this city, every bridge is a contract and every current carries coin. The bargains struck in our shadows never quite fade from memory, and the waters remember every handshake.}

\vspace{0.5em}

I am a Bobbin-Runner. I can cross this city faster than any grown-up, and I know which planks creak and which merchants lie. The bridges remember every step, and the canals carry more secrets than fish.

\vspace{0.5em}

Silkstrand is a city strung across canals and arches, a place where every bridge is a contract and every current carries rumor. The silk looms never sleep, the Exchange never closes, and the Archivolt's notaries seal fates with a drop of wax. Merchants rise and fall with the tide of the market, bravos carve reputations on the planks of Three-Queens Bridge, and the city's underworld runs on whispers and borrowed keys. But here's the secret the Matron doesn't want you to know: the canals are not just water. They are memory. Every bribe, every broken promise, every cargo that was never declared---the canals remember. And sometimes, at low tide, the water shows you the faces of the merchants who drowned rather than pay their debts.

\vspace{0.5em}

The Matron rules from her palazzo stairs with velvet and iron, but the city itself belongs to silk, silver, and shadow. The Exchange Volatility Timer ticks with every major cargo seizure, every trade route cut, every merchant default. When it fills, a panic hits---the Exchange closes, and all outstanding contracts are voided unless witnessed by the Matron herself. The Matron has voided nine contracts this year. The ninth was her own.

\vspace{0.5em}

Never name a rival while crossing a bridge at night. The water remembers, and the curse you send may come back wrapped around your own throat.

\vspace{1em}

\noindent\textit{--- Bobbin-Runner, who ties a knot in her apron string for every secret she carries}
}

\vspace{1em}

\begin{almanac}
\textbf{What do people eat for breakfast?} Rice porridge with dried fish---the rice is from the Oshiiran barges, the fish is from the Zakov fleets. Both are smuggled. Both are delicious. Both are paid for in favors.

\textbf{What does a child learn before they can walk?} To count bridges. Silkstrand children are taught the eight bridges of the city: Merchant's, Weaver's, Dyer's, Cooper's, Pilot's, Beggar's, Saint's, and the ninth---the bridge that is only there at low tide. The ninth bridge leads to the Sunken Quarter. The Sunken Quarter is where debts go to drown.

\textbf{What is the first thing you notice about a stranger?} Their ink. Every Silkstrand merchant has a seal-ink stain on their right thumb. A stranger whose ink is fresh is new to the trade. A stranger whose ink is faded has been retired. A stranger with no ink is not a merchant. Not a merchant is not a person.

\textbf{What is the most common crime?} Contract fraud---but fraud is not a crime, it is a negotiation. The contracts are written in disappearing ink. The ink's formula is a secret. The secret is held by the Archivolt notaries. The notaries sell the antidote. The antidote is the fee.

\textbf{What does no one talk about but everyone knows?} The Matron is not the ruler. The canals rule. The canals set the prices. The canals decide who lives and who drowns. The canals are not water. The canals are the city's memory. The memory is the Matron's. The Matron is the canals' voice.

\textbf{Where do people go when they die?} To the canals. The dead are not buried. Their bodies are wrapped in silk and sunk in the deep channels. The channels are the city's ledger. The ledger is never balanced.

\textbf{How do you apologize for a serious wrong?} You offer a bridge---a toll paid to the Bridge-Lords. The toll is calculated by the wronged party's estimate of their loss. The toll is paid in silver. The silver is melted and poured into the canal. The canal remembers the weight. The weight is your apology.

\textbf{What is the most important festival?} The Dyeing of the Waters (spring). The dyers empty their vats into the canals. The water turns red, blue, green, gold. The colors are the year's fashions. The fashions are set by the Matron. The Matron's favorite color is the one that sells best. The best color is the one she dyed her own hair.

\textbf{What is the secret that could destroy a powerful person?} The canals are not natural. They were dug by a Sidhi engineer who was paid in promises. The promises were broken. The engineer cursed the city. The curse is that every ninth drop of water is a tear. The tears belong to the engineer. She is still crying. She has been crying for nine centuries.
\end{almanac}

\newpage
\section{Silkstrand --- ``City of Bridges \& Dyewater''}
\label{chap:silkstrand}

\begin{quote}
  \textbf{The Matron (Elite):}  
  ``In this city, every bridge is a ledger and every current carries coin. The dye that stains our waters never washes out, just as the bargains struck in our shadows never quite fade from memory.''
\end{quote}

\begin{quote}
  \textbf{Bobbin-Runner (Commoner):}  
  ``I can cross this city faster than any grown‑up, and I know which planks creak and which merchants lie. The bridges remember every step, and the canals carry more secrets than fish.''
\end{quote}

\begin{tcolorbox}[colback=black!3,colframe=blue!60!black,title={Theme \& Atmosphere}]
Silkstrand is a city strung across canals and arches, a place where every bridge is a ledger line and every current carries rumor. The dyes that stain its waters never wash out, and neither do the bargains struck in shadow. Merchants rise and fall with the Exchange, bravos carve reputations on the planks of Three‑Queens Bridge, and curses cling to the very cloth that leaves the looms. The Matron rules from her palazzo stairs with velvet and iron, but the city itself belongs to silk and tide.

\vspace{2mm}
\textbf{What you notice first:} The clash of dye‑stained water against ancient stone. The constant rumble of cart wheels on bridge planks. The way every handshake lingers a heartbeat too long---as if both parties are checking for threads.

\textbf{The one rule that kills newcomers:} Never spit into a canal while naming a rival. The water remembers, and the curse you send may come back wrapped around your own throat.
\end{tcolorbox}

\subsection*{Regional Mood: Bridges and Dyewater}

Here, beauty is bought with blood and color, and both wash downstream. Silkstrand is not lawless---it is *layered*. Above the tide lies the world of guild charters, notarial seals, and the Matron's velvet justice. Below the bridges, in the weeping arches and the midnight boat lanes, a quieter economy runs on whispers, borrowed keys, and the careful omission of names. The city's bravos duel for reputation on Three‑Queens Bridge while factors bribe clerks in the Archivolt. Every alliance is a knot that can be pulled loose. Every fortune is a thread that can be cut.

\textbf{Starting Location:} Three‑Queens Bridge, dripping with dyewater, its stones painted in a thousand forgotten shades. A Guildmistress with dyed hands offers a dye‑permit. "The Redwater boils," she says. "And something in the canals remembers your scent. Bring me three bolts of untaxed crimson before the Exchange closes, or the Archivolt will find your name in a dead notary's ledger."

\paragraph*{(Bridge/Canal/Mill)} Mulberry garths and wormhouses outside North Gate with steam and sweet rot; Filature hall where whispers travel faster than steam.

\section*{Spades --- Places}
\label{sec:silkstrand-places}
\begin{enumerate}
\setcounter{enumi}{1}
\item \textbf{Mulberry Garths} --- Wormhouses outside North Gate; steam, sweet rot, desperate pickers.
  \hfill \textit{The pickers" hands are stained black. They say the color never fades, even after death.}
\item \textbf{Filature Hall} --- Cocoon‑boil chambers; whispers travel faster than steam.
  \hfill \textit{The steam carries voices from last season. Some say they"re warnings.}
\item \textbf{Redwater Dyeworks} --- Along stain‑canal; brick stained forever with forgotten dyes.
  \hfill \textit{The dye‑runners wear lead masks. The ones who don"t die young---or wrong.}
\item \textbf{Spindle Tower} --- Creaking windlass‑lifts and posted rates; heights and hazards.
  \hfill \textit{The lift operator has a bell. He rings it once for every passenger who didn"t pay.}
\item \textbf{Three-Queens Bridge} --- Stacked market stalls; cells under arches, reputations made and broken.
  \hfill \textit{The cells are for debtors. The walls are washed with canal water so no one forgets the smell.}
\item \textbf{Salt Gate Quay} --- Rope booms, chalk tallies, tired eyes; customs and coin.
  \hfill \textit{The chalk is mixed with ground bone. A false tally writes itself in a dead hand.}
\item \textbf{Silk Exchange} --- Floor with chalk circles and clappers for opening bids.
  \hfill \textit{The clappers are whale ribs. The bids echo, and sometimes the echo answers.}
\item \textbf{Ropewalk Sheds} --- Arrow‑straight; bruisers hired by the yard, muscle for hire.
  \hfill \textit{The rope‑walkers never look back. They say a rope remembers who pulled it last.}
\item \textbf{Old Arsenal} --- Free Company's barracks now; arms and ambitions stored.
  \hfill \textit{The arsenal's bell rings at odd hours. No one admits ringing it.}
\item[J] \textbf{The Archivolt} --- Arcaded street of notaries, seals, and quiet knives.
  \hfill \textit{The notaries write in lemon juice. Heat reveals the truth---or a facsimile.}
\item[Q] \textbf{Basilica of Azerin} --- Weaver‑saint; confraternity rooms hum with vows.
  \hfill \textit{The saint's icon has one eye. The other was traded for the secret of red dye.}
\item[K] \textbf{Palazzo della Matrona} --- Ruling seat with private river stairs; velvet and iron.
  \hfill \textit{The Matrona's barge is rowed by debtors who have repaid in years, not coin.}
\item[A] \textbf{Flood-Stairs} --- Bronze flood marks, bell rope ready when the Strand runs wild.
  \hfill \textit{The flood marks go up to nine. The ninth step is underwater---always.}
\end{enumerate}

\paragraph*{(Guild/Factor/Crown)} Bobbin‑runner child with feet sure on parapets; foreign factor seeking warehouse and friend at customs.

\section*{Hearts --- People \& Factions}
\label{sec:silkstrand-people}
\begin{enumerate}
\setcounter{enumi}{1}
\item \textbf{Bobbin-Runner} --- Child with sure feet on parapets, rumors in pocket.
  \hfill \textit{She ties a knot in her apron string for every secret she carries. The apron is nearly solid knots.}
\item \textbf{Mulberry Steward} --- Counting leaves; desperate for pickers, quick with coin.
  \hfill \textit{His left hand is missing two fingers---lost to a cocoon‑rot curse. He calls it a bargain.}
\item \textbf{Foreign Factor} --- Seeking warehouse and friend at customs; connections for sale.
  \hfill \textit{He carries three ledgers. One for the Matron, one for the Archivolt, one for himself.}
\item \textbf{Guildmistress} --- Dyers' hands stained, permits tighter than purse strings.
  \hfill \textit{She can tell your guild by the stain on your cuff. She never forgets a color.}
\item \textbf{Bridge Bailiff} --- Rents stalls, sells gossip by the breath; power in position.
  \hfill \textit{He rings a small bell before each announcement. The tone tells you if it's truth or theater.}
\item \textbf{Archivolt Notary} --- "Fixes" missing recitals---for a donation; law as service.
  \hfill \textit{His inkwell contains a single drop of his own blood. He renews it each spring.}
\item \textbf{Watch Captain} --- Condotta to three lords; passwords change with the wind.
  \hfill \textit{He wears a mask on feast days. No one has seen his face in seven years.}
\item \textbf{Spinner-Matron} --- Wormhouses; temper like hot copper, secrets in the steam.
  \hfill \textit{She names each cocoon after a dead rival. The silk is said to carry their complaints.}
\item \textbf{Exchange Caller} --- A clap can still a thousand voices; auction master.
  \hfill \textit{He lost his voice to a curse. Now he claps in a language only the bids understand.}
\item[J] \textbf{Night-Boat Smuggler} --- "Ravel," silent oar, louder favors; shadows for coin.
  \hfill \textit{His boat has no name. The canal would remember it, and memory is a kind of tax.}
\item[Q] \textbf{The Matron} --- Patient, velvet, iron; the city's thread pulled through her hands.
  \hfill \textit{She has never raised her voice. The city learned to fear her silences.}
\item[K] \textbf{Lame King's Envoy} --- Velvet boots; claims certain alleys "protected."
  \hfill \textit{His cane is hollow. Inside is a letter of marque that expired a decade ago.}
\item[A] \textbf{Saint of Warps} --- If real: sees curses braided in cloth; holy unraveler.
  \hfill \textit{Her eyes are milk‑white. She sees not the cloth, but the threads of obligation within.}
\end{enumerate}

\paragraph*{(Flood/Interdict/Riot)} Flood siren with gates lowering---move crates or kiss them goodbye; quarantine flag at Redwater sealing dyers' row.

\section*{Clubs --- Complications/Threats}
\label{sec:silkstrand-complications}
\begin{enumerate}
\setcounter{enumi}{1}
\item \textbf{Flood Siren} --- Gates lowering; move crates or kiss them goodbye to tide.
  \hfill \textit{The gate‑keeper's palm is marked with a bell. Cross it with silver to hurry.}
\item \textbf{Quarantine Flag} --- Redwater dyers' row sealed; your cargo implicated.
  \hfill \textit{The flag is dyed with lamp‑black and something else---something that smells of graves.}
\item \textbf{Loom Strike} --- Over "bad cocoons"; streets fill with idle frames.
  \hfill \textit{The striking weavers wear red threads around their wrists. They will not speak first.}
\item \textbf{Seal Counterfeit} --- Discovered at Archivolt; all contracts frozen.
  \hfill \textit{The forged seal is warm. The real one is cold. Touch both to know which is which.}
\item \textbf{Bridge Riot} --- A dropped stall blocks lanes; tempers boil faster than vats.
  \hfill \textit{The riot started over a missed bell‑ring. Or a stolen bolt. Or a look.}
\item \textbf{Condottieri Flip} --- Colors change; watchwords shift, escorts vanish.
  \hfill \textit{The captain's new banner is stitched with thread from his old captain's shroud.}
\item \textbf{Blackwood Panic} --- Witch‑posts hammered on quay doors overnight.
  \hfill \textit{The posts are from the Direwood edge. They sweat sap that smells of old fear.}
\item \textbf{Fungus Blight} --- Silk‑fungus; wormhouses burn, refugees flood piazze.
  \hfill \textit{The fungus glows faintly in the dark. It spells names---sometimes yours.}
\item \textbf{Salt-Tax Doubled} --- At Salt Gate; boats stack three deep, bribes flow.
  \hfill \textit{The tax collector wears a mask. Beneath it, his face is a salt‑caked skull.}
\item[J] \textbf{Duel Challenge} --- Booked on Three‑Queens Bridge; you"re named as seconds.
  \hfill \textit{The challenge is written on silk. Burn it to refuse---and lose face forever.}
\item[Q] \textbf{Exchange Corner} --- Rivals hoard raw thread; prices go vertical.
  \hfill \textit{The cornering merchant has a ledger bound in human skin. The names inside are his competitors.}
\item[K] \textbf{Saint's Procession} --- Closes half the wards; ignore at peril, faith walks the streets.
  \hfill \textit{The icon weeps dyewater. The faithful collect it in silver cups.}
\item[A] \textbf{Curse Awakened} --- In canals: no route reaches the address you seek.
  \hfill \textit{The water is still. Your reflection is gone. You are still here---for now.}
\end{enumerate}

\paragraph*{(Permit/Seat/Escort)} Bridge token for peak‑hour cart crossing free; dye‑permit chit to process "questionable" color this week.

\section*{Diamonds --- Rewards/Leverage}
\label{sec:silkstrand-rewards}
\begin{enumerate}
\setcounter{enumi}{1}
\item \textbf{Bridge Token} --- One peak‑hour cart crossing free; traffic as currency.
  \hfill \textit{The token is a bronze coin with a hole in the center. The hole is the toll.}
\item \textbf{Dye-Permit} --- Process "questionable" color this week; law bends for coin.
  \hfill \textit{The permit is written on dyed vellum. The color changes with the tide.}
\item \textbf{Warehouse Seal} --- From Matron's office; rivals legally shut out.
  \hfill \textit{The seal is a brass bell. Ring it once to claim the space.}
\item \textbf{Exchange Pass} --- Day pass to trade without sponsor; freedom in paper.
  \hfill \textit{The pass is stamped with a spider. The spider's legs are the signatures of three guilds.}
\item \textbf{Watergate Priority} --- Skip one flood closure when it matters; tide waits for none.
  \hfill \textit{The priority token is a dried eel skin. Show it, and the gate‑keeper spits but stands aside.}
\item \textbf{Notarial Indulgence} --- Retrofit a missing recital; past deals stand.
  \hfill \textit{The indulgence is written in lemon juice. Heat it to see the terms.}
\item \textbf{Wormhouse Allotment} --- Claim a share of the next hatch; future in cocoons.
  \hfill \textit{The allotment is a silk‑worm egg. Keep it warm, or it will hatch into a debt.}
\item \textbf{Ropewalk Credit} --- Hire crews before you have coin; muscle on credit.
  \hfill \textit{The credit is a knotted rope. Cut a knot for each favor called.}
\item \textbf{Arsenal Key} --- One night's issue for retinue; arms in trusted hands.
  \hfill \textit{The key is a finger bone. The lock is a gauntleted fist.}
\item[J] \textbf{Condotta Rider} --- City watch escorts through any ward, once.
  \hfill \textit{The rider's horse has iron shoes. The sparks tell him who to watch.}
\item[Q] \textbf{Matron's Audience} --- Private whisper becomes policy; velvet commands.
  \hfill \textit{The Matron's ear is a bell. Speak into it, and the city listens---once.}
\item[K] \textbf{Tax-Farm Share} --- On Three‑Queens stalls for a season; profit in peddlers.
  \hfill \textit{The share is a tally stick. Break it in half to claim your portion.}
\item[A] \textbf{Golden Thread} --- Temporary charter to set tolls on a named canal.
  \hfill \textit{The thread is real gold. It will not break---but it will stretch.}
\end{enumerate}

\section*{Quick use notes}
\label{sec:silkstrand-quick-use}
\begin{itemize}
\item Draw until you have all four suits: Spade = place, Heart = actor, Club = pressure, Diamond = leverage. Highest rank sets the main Clock (2--5 $\rightarrow$ 4, 6--10 $\rightarrow$ 6, J/Q/K $\rightarrow$ 8, A $\rightarrow$ 10).
\item Diamonds are codified outcomes (permits/seats/escorts) that change position rather than call for a roll.
\item If any A appears, echo \textbf{silk and water} motifs---threads that bind, currents that remember, bargains that stain canal beds.
\end{itemize}

\subsection*{Additional Features}
\begin{itemize}
\item \textbf{Duel Etiquette:} Duels fought on bridges draw crowds eager to wager on blood. Refusing stains a name worse than losing.
\item \textbf{Dyewater Omens:} Canals run strange colors at dawn; red foretells riots, blue brings calm, black means curses awakened.
\item \textbf{Matron's Seal:} Any contract bearing her seal is law---but the Matron is known to rescind seals without warning, leaving debts tangled.
\end{itemize}

\subsection*{Lore Echoes (Cross‑Expansion Hooks)}
\begin{itemize}
  \item \textbf{The Bitter Root:} Daughters of the Cord run a hospice near the Salt Gate. They ask no payment---but they cut threads. A renegade Sister hides among the bobbin‑runners, trying to untie her own cord.
  \item \textbf{The Threshold Compendium:} The Archivolt notaries keep a secret ledger of unquiet debts---Class II Anchors bound to unpaid contracts. Saikou Ira would have a seizure here.
  \item \textbf{Malachai, the Chained Angel:} A cursed bolt of silk surfaces on the Exchange every decade. Anyone who wears it prospers---until the next full moon. The bolt is due to appear tonight.
  \item \textbf{The Grey Wanderer's Grimoire:} Silkstrand psions are called "Weft‑Readers." They see the threads of obligation. The Matron employs three---and has never let them meet each other.
  \item \textbf{The Ninth:} In any Silkstrand scene, the ninth bridge crossed in a day leads somewhere unexpected. The wise stop at eight and take a boat.
\end{itemize}

\subsection*{Silkstrand Superstitions (burn for +1 die once)}
\begin{itemize}
  \item \textbf{Bread to the Canal:} Toss a crust into the water before crossing a new bridge. Ignore the first \emph{Bridge Riot} penalty.
  \item \textbf{Knot the Thread:} Tie a single knot in a scrap of silk. Cancel \emph{Seal Counterfeit} or take --1 die on your next negotiation (your choice).
  \item \textbf{Name to the Dyewater:} Whisper a debtor's name into the canal at dawn. Gain +1 Effect when collecting from them this scene, but mark \emph{Obligation} to the Archivolt.
\end{itemize}

\begin{tcolorbox}[colback=black!3,colframe=black!40!white,title={Patronage \& Power}]
In Silkstrand, power flows through commerce, connections, and the careful weaving of influence. The Matron maintains control through a web of guilds, factors, and condottieri, each bound by contracts and customs that shift like the tides. Merchants rise and fall with the Exchange, while the city's bravos and smugglers operate in the spaces between law and profit.

\textbf{For the GM:}  
Patronage in Silkstrand revolves around access---to bridges, to warehouses, to the Exchange floor itself. Rewards often take the form of permits, seals, or safe passage that can be leveraged into greater influence. To emphasize this:
\begin{itemize}
\item Tie rewards to visible symbols (seals, tokens, contracts) that can be challenged, stolen, or voided.
\item Let rival guilds issue conflicting charters, forcing players to choose whose favor matters more.
\item Use the bridges, markets, and Exchange as arenas for social contests, where reputation and connections determine success.
\end{itemize}
In Silkstrand, your word is your weave, and your weave determines whether you rise or drown.
\end{tcolorbox}

% Index entries (placeholder, keep existing style)
\index{Silkstrand|see{City of Bridges and Dyewater}}
\index{Theme generator}
\index{Bridges and Dyewater generator}
\index{Places!Silkstrand}
\index{People!Silkstrand}
\index{Complications!Silkstrand}
\index{Rewards!Silkstrand}
\index{The Matron}
\index{Three-Queens Bridge}
\index{Archivolt}
\index{Redwater Dyeworks}
\index{Silk Exchange}
\index{Flood-Stairs}
\index{Saint of Warps}
\index{Bobbin-runner}
\index{Guildmistress}
\index{Bridge bailiff}
\index{Night-boat smuggler}
\index{Exchange caller}
\index{Weft-Curse}

\section*{Thematic SB Spend Table}
\label{sec:silkstrand-sb}

\subsection*{Minor Complications (1 SB)}
\begin{itemize}
\item \textbf{Exposure:} Your actions draw unwanted attention from \textbf{city watch or rival merchants}.
\item \textbf{Noise:} Sounds of your actions alert nearby \textbf{crowds or bridge traffic}.
\item \textbf{Trace:} Evidence of your passage marks your route for \textbf{trackers or customs}.
\item \textbf{Delay:} A brief but meaningful setback costs you \textbf{time or favorable tide}.
\item \textbf{Supply Strain:} Mark +1 segment on a relevant \textbf{resource clock}.
\end{itemize}

\subsection*{Moderate Setbacks (2 SB)}
\begin{itemize}
\item \textbf{Alarm Raised:} \textbf{Bridge bailiff or guildmaster} becomes aware and begins responding.
\item \textbf{Position Lost:} You lose advantageous ground/cover/stealth due to \textbf{crowd movement or flood}.
\item \textbf{Foe Appears:} A \textbf{rival merchant or hired bravo} arrives on scene.
\item \textbf{Gear Trouble:} A piece of equipment becomes \textbf{Compromised/Neglected}.
\item \textbf{Lock/Barrier:} A simple obstacle now requires a test to overcome.
\end{itemize}

\subsection*{Serious Trouble (3 SB)}
\begin{itemize}
\item \textbf{Reinforcements:} Additional \textbf{condottieri, smugglers, or guild enforcers} arrive.
\item \textbf{Key Gear Breaks:} A crucial tool/weapon becomes temporarily unusable.
\item \textbf{Major Twist:} The situation fundamentally changes - \textbf{flood hits/contract voided/duel declared}.
\item \textbf{Rail Tick:} Advance a relevant campaign/front clock by 1 segment.
\item \textbf{Condition Applied:} Mark \textbf{Fatigue 1/Harm 1/Condition} appropriate to fiction.
\end{itemize}

\subsection*{Major Turns (4+ SB)}
\begin{itemize}
\item \textbf{Trap Springs:} A prepared danger activates with full effect.
\item \textbf{Authority Arrival:} \textbf{The Matron, Watch Captain, or Guildmistress} intervenes.
\item \textbf{Scene Shift:} The environment changes dramatically - \textbf{flood rises/bridge collapses/procession blocks streets}.
\item \textbf{Patron Omen:} Divine/arcane forces take notice - \textbf{omen appears/blessing lost/curse manifests}.
\item \textbf{Narrative Pivot:} The story takes an unexpected turn that reframes objectives.
\end{itemize}

\subsection*{Region-Specific SB Options}
\begin{itemize}
\item \textbf{Silkstrand (Bridges \& Trade):} Bridge planks shift underfoot, bargains restructure mid‑speech, coins change hands without touch.
\item \textbf{Silkstrand (Dyewater):} Canal colors shift ominously, dyes stain more than cloth, water carries whispered names.
\item \textbf{Silkstrand (Duel Culture):} Challenges are issued mid‑conversation, seconds appear from the crowd, honor prices rise with each breath.
\end{itemize}

\subsection*{Boss Seeds (three‑phase foes)}
\begin{itemize}
\item \textbf{The Corseted Ledger (cursed accounting bound in silk)} --- \emph{Phases}: Friendly Audit $\rightarrow$ Retroactive Debts $\rightarrow$ Names as Collateral. \emph{Cracks}: burn one forged seal in public; settle a rival's debt at cost; wash the book on Flood‑Stairs at high bell.
\item \textbf{Archivolt Ragman (collector of "missing recitals")} --- \emph{Phases}: Quiet Inserts $\rightarrow$ Paper Hunts $\rightarrow$ Bridge Vetoes. \emph{Cracks}: win a duel for precedent; expose a double‑entry altar; break his pen with the Matron's Seal.
\item \textbf{Color That Walks (dyewater revenant)} --- \emph{Phases}: Stain $\rightarrow$ Spread $\rightarrow$ Drown. \emph{Cracks}: confess on Three‑Queens; bind scent‑rights properly; let the Saint of Warps unpick a life's thread (costly).
\end{itemize}
\newpage

\chapter{Midh-Akhaz}
\emph{The Violet Gate.}

\noindent
{\Large
\textit{Coin talks. Grass walks. Blood settles the rest.}

\vspace{0.5em}

I am a Horse-Braid Apprentice. I tie a new knot for every paying passenger. My hair is a map of debts --- and grudges. The wind carries names, not ledgers, but here, in the shadow of the Oasis Compact, both trade and oaths change hands. Do not mistake our patience for weakness. The grass remembers every hoof that crossed it.

\vspace{0.5em}

Midh Ahkaz is not a city. It is a wound that learned to trade. A day's ride north of the Ecktorian border, it squats on the edge of the Violet Steppe like a patient wolf. No walls. No gates. Just rings of felt tents, latrine trenches, and the smell of horse sweat, cedar smoke, and the faint copper tang of blood money.

\vspace{0.5em}

The grass grows tall enough to hide a body. The locals cut it back at dawn --- but never at dusk, because the things that crawl through it at night have sharp teeth and no patience for property lines. Ecktorian emissaries call this place ``the Brink.'' Caravan guards call it home. Thieves call it paradise.

\vspace{0.5em}

Every morning, the Coin-Weigh Tribunal opens its scales under a canvas dome. Every night, the Dust Market auctions goods that fell off a wagon --- or off a corpse. The Khatun's word is law until someone stronger speaks. And in the Silver Stall Stable, a horse thief can buy a fresh mount with a story and a knife.

\vspace{0.5em}

Never refuse the guest-cup. The first pour is courtesy. The second is negotiation. The third is a blood-price unspoken. And the fourth --- the fourth means you're already dead, you just haven't stopped breathing.

\vspace{1em}

\noindent\textit{--- Horse-Braid Apprentice, who ties a knot for every secret she carries}
}

\vspace{1em}

\begin{almanac}
\textbf{What do people eat for breakfast?} Curd cheese and flatbread --- the cheese is from the Ykrul herds, the bread is baked in the embers of a dung fire. It tastes of smoke and the open sky. The poor dip it in salt; the rich dip it in honey stolen from a border hive.

\textbf{What does a child learn before they can walk?} To count horses. Midh Ahkaz children are taught the eight gaits of a healthy mount: walk, trot, canter, gallop, pace, amble, rack, and the ninth --- the death-sprint that only a stolen horse knows. The ninth gait is never taught; it is discovered when the rider's life depends on it.

\textbf{What is the first thing you notice about a stranger?} Their braids. Ykrul warriors braid their hair with clan colors; Vilikari traders wear oiled braids with silver beads. A stranger whose braids are loose has just fought. A stranger whose braids are cut has lost a blood-price. A stranger with no braids is either a slave or a ghost. Neither stays long.

\textbf{What is the most common crime?} Grass-right forgery --- but forgery is not a crime, it is a negotiation. The Oasis Compact issues branded sticks for grazing rights. A clever carver can duplicate the brand. The penalty is exile into the tall grass at dusk. The grass does not have a court of appeals.

\textbf{What does no one talk about but everyone knows?} The Khatun is not the only ruler. The Silver Stall Stable master knows where every stolen horse sleeps. The Dust Market auctioneer knows the price of every sin. And the Iron Well keeper knows the name of everyone who will die before the next grass moon. The three of them meet once a season. No one records what they say. The grass records it.

\textbf{Where do people go when they die?} To the grass. The dead are not buried. Their bodies are carried beyond the caravan rings and left for the wolves and the wind. The grass grows tall and green where a body falls. A patch of deep violet is a sign that the dead were not properly mourned. The locals avoid those patches.

\textbf{How do you apologize for a serious wrong?} You offer a blood-price in silver, weighed on the Coin-Weigh Tribunal's scales. If the scale tips in your favor, you are forgiven. If it tips against you, you owe a horse. If the scale balances, you owe a story --- and the story must be true, and the true story must make the wronged party weep. If it does not, you owe a finger.

\textbf{What is the most important festival?} The Grass Moon (autumn equinox). The violet seed-heads glow under moonlight. The tribes ride in a great circle, singing the names of the dead. At midnight, the Khatun pours a cup of water onto the earth. If the water sinks, the winter will be mild. If it pools, the winter will be hungry. If it steams, the Hollow is walking. The festival ends at dawn. No one sleeps.

\textbf{What is the secret that could destroy a powerful person?} The Oasis Compact's original seal was not lost. It was stolen by a Vilikari factor who wanted to control water rights. The factor is still alive, living in the Foreign Quarter under a false name. The seal is hidden in a hollow horse skull buried under the Grass Ring. The skull belonged to the Khatun's favorite stallion. The stallion did not die of old age.
\end{almanac}

\newpage
\section{Midh Ahkaz --- ``The Violet Gate''}
\label{chap:midh-ahkaz}

\begin{quote}
  \textbf{Khatun of the Ring (Ykrul Elder):}  
  ``The wind carries names, not ledgers. But here, in the shadow of the Oasis Compact, both trade and oaths change hands. Do not mistake our patience for weakness. The grass remembers every hoof that crossed it.''
\end{quote}

\begin{quote}
  \textbf{Vilikari Factor (Night Market Regular):}  
  ``Ecktoria calls it uncivilized. We call it unbeatable. No tariffs, no seals, no questions. Coin talks. The rest walks.''
\end{quote}

\begin{tcolorbox}[colback=black!3,colframe=green!60!black,title={Theme \& Atmosphere}]
Midh Ahkaz is not a city of stone—it is a city of felt, iron, and shifting grass. A day’s ride north of the Ecktorian border, it squats on the edge of the Violet Steppe like a patient wolf. Here, the Ykrul, Vilikari, Ostrikari (their eastern kin from the grassy fringes of this same prairie), and a dozen other peoples gather to trade, haggle, and settle grudges. The air smells of horse sweat, cedar smoke, and the faint copper tang of blood money.

The city has no walls. Instead, rings of tents and caravanserais spiral outward from the central Counting House. At dawn, the grass—tall enough to hide a horse—is cut back. At dusk, the perimeter fires are lit. The colour of the steppe comes from the violet hue of mature grass seeds, which ripen in waves across the prairie. The Ecktorian emissaries call it “the Brink.” Locals call it home.

\textbf{What you notice first:} The constant low drumming of hooves. The way every merchant has a scale on their table. The braids in every horseman’s mane—each knot a story, a debt, a promise.

\textbf{The one rule that kills newcomers:} Never refuse the guest-cup. The first pour is courtesy. The second is negotiation. The third is a blood‑price unspoken.
\end{tcolorbox}

\begin{tcolorbox}[colback=blue!5!white,colframe=blue!60!black,title={Western Steppe vs. Eastern Steppe (GM Sidebar)}]
The Ykrul divide their world into two great grass seas:

\medskip
\noindent\textbf{Western Steppe (Violet Prairie):}  
Rolling, fertile, forgiving. The grass grows tall enough to hide a crouched rider, but the land is laced with caravanserais, seasonal wells, and the smoke of permanent camps. This is where Midh Ahkaz sits—the soft underbelly of the steppe, where trade happens and clans winter. The violet seed‑heads give the prairie its name.

\medskip
\noindent\textbf{Eastern Steppe (True Steppe):}  
Vast, dry, harsh. Shared with the Kuvani horse‑clans, it stretches east of the Valewood. Grass is shorter, winds are colder, and a rider can see a day’s ride in every direction. The Ostrikari who appear in Midh Ahkaz are from the eastern fringes of the \textit{Western} Steppe—they speak with a harder accent, braid their horses’ tails differently, and carry bows of horn and sinew, not wood.

\medskip
\noindent\textit{Why this matters:} A Ykrul from the Eastern Steppe might find Midh Ahkaz almost claustrophobic. A trader from the Violet Prairie might feel lost in the east. Use the contrast to create tension between clans or to flavour foreign visitors.
\end{tcolorbox}

\subsection*{Regional Mood: Wind, Grass, and Contract}

Midh Ahkaz thrives on movement. Caravans arrive from every direction: Ykrul horse‑tribes from the deep steppe, Ostrikari mercenaries from the eastern fringes, Vilikari traders from the river towns, and the occasional bold Ecktorian factor with a silver tongue and a hidden dagger. The city’s “law” is the Oasis Compact—a shifting web of seasonal rights, grazing permits, and arbitration courts held under open sky. Disputes are settled by the Bowl (fairness) then the Board (geometry). A broken promise can cost a herd. A kept oath can buy a ford for a generation.

\textbf{Starting Location:} The Violet Bazaar, a sprawling market of felt tents and chalk circles, where horse traders shout over the jingle of brass bridles. A Vilikari merchant with oiled braids hands you a sealed letter. “The Coin‑Weigh Tribunal suspects a counterfeit seal. Find the forger before the Grass Moon, or the contract on your head becomes public.”

\paragraph*{(Ring/Well/Cairn)} The Oasis Well where water shares are drawn; counting pebbles placed on stone ledges; the distant smoke of a Steppe Ranger campfire.

\section*{Spades --- Places}
\label{sec:midh-ahkaz-places}
\begin{enumerate}
\setcounter{enumi}{1}
\item \textbf{Violet Bazaar} --- Felt stalls, chalk circles, horse traders and rumor merchants.
  \hfill \textit{The chalk circles are for disputes. Step inside and speak truth—or pay in grass.}
\item \textbf{Coin‑Weigh Tribunal} --- Tariff court under a canvas dome; scales of bronze and iron.
  \hfill \textit{The scales were blessed by a dying Sky‑Speaker. They tip even when the wind is still.}
\item \textbf{Oasis Well} --- Stone‑rimmed pool, token-keeper with a leather bag of water rights.
  \hfill \textit{Each token is a carved pebble. Lose it, and your herd drinks dust.}
\item \textbf{Steppe Ranger Camp} --- Tent line on the eastern edge; trackers, scouts, feud mediators.
  \hfill \textit{The rangers wear wolf tails. They do not smile—they wait.}
\item \textbf{Caravan Rings} --- Concentric circles of wagon shelters; fires lit inward, watch posts on ox carts.
  \hfill \textit{The inner ring is for sworn allies. Outsiders sleep in the second ring, coin in hand.}
\item \textbf{Sky‑Speaker’s Cairn} --- Stone heap atop a low hill; wind‑knot ropes and bones.
  \hfill \textit{The shaman leaves offerings of horsehair. The wind carries his prayers to Kuva.}
\item \textbf{Dust Market} --- No‑questions auction at the city’s edge; stolen goods, cursed relics, bad memories.
  \hfill \textit{The auctioneer wears a hooded mask. No one knows his face. No one has asked twice.}
\item \textbf{Silver Stall Stable} --- Remount station; horses for hire, gossip for free, blood for more.
  \hfill \textit{The stable master has three fingers on his left hand. He lost the others in a bet over a mare.}
\item \textbf{Grass Ring} --- Ritual dueling ground; mourners and bookies share benches.
  \hfill \textit{The grass grows back faster after a duel. The dead fertilize the peace.}
\item[J] \textbf{Foreign Quarter} --- Ecktorian, Vhasian, and Viterran merchants in permanent felt mansions.
  \hfill \textit{The Ecktorian envoy winters here. He has not left in five years.}
\item[Q] \textbf{Khatun’s Ring} --- Largest tent, carpeted with woven battle scenes; the old chieftain’s seat.
  \hfill \textit{The carpet is made of enemy cloaks. Each sash is a surrendered banner.}
\item[K] \textbf{Under‑Arcades} --- Earthen tunnels below the bazaar; root‑walled, bone‑dusted.
  \hfill \textit{The walls are packed with bone dust and old burrows. The dead whisper route tips.}
\item[A] \textbf{Iron Well} --- Ancient Aeler shaft, now used for vengeance contracts and whispered names.
  \hfill \textit{The water tastes of rust and memory. Those who drink never lie again—or never speak.}
\end{enumerate}

\paragraph*{(Horse/Tally/Token)} Sky‑Speaker’s apprentice with a drum and a dream; Ostrikari arms‑master selling short bows; Ykrul elder counting braids.

\section*{Hearts --- People \& Factions}
\label{sec:midh-ahkaz-people}
\begin{enumerate}
\setcounter{enumi}{1}
\item \textbf{Horse‑Braid Apprentice} --- Young Ykrul with a head full of trade routes and a knot diary.
  \hfill \textit{She ties a new braid for every paying passenger. Her hair is a map of debts.}
\item \textbf{Ostrikari Arms‑Master} --- Selling short bows and longer grudges; fights for coin.
  \hfill \textit{He names each bow after a dead rival. The bows never miss.}
\item \textbf{Coin‑Weigh Clerk} --- Ink‑stained, patient, scales in her eyes; never lies, never forgives.
  \hfill \textit{She records weights in blood. A single drop seals a verdict.}
\item \textbf{Vilikari Factor} --- Braided, oiled, quick with a smile and a second ledger.
  \hfill \textit{His caravan has three names. Only the real one buys safe passage.}
\item \textbf{Steppe Ranger} --- Tracker with a wolf‑tail cloak and a debt to the Sky‑Speaker.
  \hfill \textit{He never draws his bow unless watched. Hunger is his real weapon.}
\item \textbf{Tokan (Ykrul Scout)} --- Mute, sharp‑eyed, always smelling the wind; her ears miss nothing.
  \hfill \textit{She signs with her knuckles. Three taps is “danger.” One tap is “run.”}
\item \textbf{Ecktoria Emissary} --- Proud, wax‑sealed, fuming at the “primitive” customs.
  \hfill \textit{He carries a copy of Imperial Code that he has never opened. The flies love his ink.}
\item \textbf{Sky‑Speaker’s Shadow} --- Elder’s quiet apprentice; listens more than the wind.
  \hfill \textit{She can name every clan’s sky‑omen. She cannot name her own origin.}
\item \textbf{Knot‑Keeper (Tulkani)} --- Counts caravan promises with a rope of nine strands.
  \hfill \textit{Each knot is a contract. Untie one and the deal dissolves.}
\item[J] \textbf{Dust Market Auctioneer} --- Hooded, genderless voice; knows the weight of every sin.
  \hfill \textit{The hammer is a thigh bone. The block is a skull altar from the Direwood.}
\item[Q] \textbf{Khatun of the Ring} --- Grey‑haired, one‑eyed, still deadly at sixty; speaks for seven clans.
  \hfill \textit{She lost her eye to a Linn raider. She killed him with his own hook.}
\item[K] \textbf{Kuva’s Voice (Wandering Priest)} --- Weathered, ash‑marked, offers omen readings for grass.
  \hfill \textit{He prays to the sky by throwing bones. The bones never lie—they just mislead.}
\item[A] \textbf{Iron Well Keeper (Vengeance Broker)} --- Mute, scarred, accepts names written in ash.
  \hfill \textit{A drink from the well costs a memory. A death costs a fortune.}
\end{enumerate}

\paragraph*{(Drought/Feud/Blight)} Dry well cuts water rations; blood‑price feud spills into the bazaar; locust swarm from the east grains.

\section*{Clubs --- Complications/Threats}
\label{sec:midh-ahkaz-complications}
\begin{enumerate}
\setcounter{enumi}{1}
\item \textbf{Dry Well Ration} --- Water shares halved; prices spike, tempers fray.
  \hfill \textit{The token‑keeper has a list. The list has a knife mark beside your name.}
\item \textbf{Blood‑Price Feud} --- A killing in the Horse Ring; kinsmen demand wergild in coin or flesh.
  \hfill \textit{The victim’s tribe drums the death rhythm. Every beat is a deadline.}
\item \textbf{Locust Swarm} --- Grass blackens; caravans stranded, fodder worthless.
  \hfill \textit{The swarm leaves a sticky residue. It smells of burnt honey and dread.}
\item \textbf{Counterfeit Coins} --- Light silver floods the Dust Market; Oasis Compact inquiry.
  \hfill \textit{The counterfeit coins are warm. The real ones are cold. Touch quickly.}
\item \textbf{Grass Fire} --- Wind‑whipped flame from the east; evacuations, scorched trade paths.
  \hfill \textit{The smoke spells a name. Yours. Or your debtor’s. Hard to tell.}
\item \textbf{Horse Sickness} --- Contagion at the Silver Stall Stable; quarantines, lost mounts.
  \hfill \textit{The sick horses dream. Their screams sound like your last argument.}
\item \textbf{Khatun’s Illness} --- Old chieftain fading; succession knives out.
  \hfill \textit{Her death will split seven clans. The Sword Oath lasts one year.}
\item \textbf{Census Raid} --- Ecktorian legion patrol “just visiting”; tensions high, bribes higher.
  \hfill \textit{The patrol leader wears a mask of imperial wax. Pull it off and see a dead man’s face.}
\item \textbf{Sky‑Omen (Red Wind)} --- Dust storm from the Ashaani desert; all navigation Desperate.
  \hfill \textit{The dust tastes of iron and old blood. Breathe it and dream of falling.}
\item[J] \textbf{Oath Challenge} --- A Kon’reh master demands geometry arbitration; your bond is the stake.
  \hfill \textit{The board is a painted skull. The pieces are horse teeth.}
\item[Q] \textbf{Tariff Strike} --- Coin‑Weigh clerks refuse to audit; trade freezes, tempers explode.
  \hfill \textit{The clerks chant a counting rhyme. Stop them before they reach nine.}
\item[K] \textbf{Duel for Rights} --- Two clans claim the same well; witnesses needed to prove water share.
  \hfill \textit{The duel is fought with blunt spears. First blood wins the well. Both bleed.}
\item[A] \textbf{Kuva’s Vengeance} --- The sky itself turns; hail, lightning, and a voice from the clouds.
  \hfill \textit{The voice names a sinner. Your name. The hail has teeth.}
\end{enumerate}

\paragraph*{(Writ/Token/Pass)} Grass Right for one season of grazing; Oath‑Bead to compel a Kon’reh master’s arbitration.

\section*{Diamonds --- Rewards/Leverage}
\label{sec:midh-ahkaz-rewards}
\begin{enumerate}
\setcounter{enumi}{1}
\item \textbf{Grass Right} --- One season of grazing on a named meadow; herders pay you tribute.
  \hfill \textit{The right is a branded stick. Burn it to end the lease—or renew.}
\item \textbf{Oath‑Bead} --- A single knot of horsehair; compels a Kon’reh master to arbitrate your dispute.
  \hfill \textit{The bead grows warm when the master speaks truth. Cold when he lies.}
\item \textbf{Water Share} --- Priority access to a well during drought; token from the Oasis Compact.
  \hfill \textit{The token is a polished pebble. Drop it in the well to claim your draw.}
\item \textbf{Steppe Ranger Favor} --- One dangerous tracking job resolved off‑screen.
  \hfill \textit{The ranger marks your name on his tent pole. He never forgets a favor.}
\item \textbf{Vilikari Voucher} --- No‑questions cargo space on three caravans; smuggle as you will.
  \hfill \textit{The voucher is a strip of oiled leather. The factor’s scent guides the mules.}
\item \textbf{Census Waiver} --- Ecktorian patrol leaves you alone for one crossing; paper shield.
  \hfill \textit{The waiver is stamped with a bear. The bear’s eyes follow you.}
\item \textbf{Khatun’s Whistle} --- Summon a small warband once (Cap 3) for a single scene.
  \hfill \textit{The whistle is a hollow horse tooth. Blow it and the grass shivers.}
\item \textbf{Sky‑Speaker’s Omen} --- One weather forecast that cannot fail; plan your escape.
  \hfill \textit{The omen is a bone thrown into fire. Read the cracks—carefully.}
\item \textbf{Silver Stall Chit} --- One free remount exchange; never be horseless.
  \hfill \textit{The chit is a brass hoof. Present it and the stable master nods once.}
\item[J] \textbf{Black Ledger Entry} --- A written promise from a rival, witnessed by three camp elders.
  \hfill \textit{The entry is written in ash on felt. Wash it and the vow dissolves.}
\item[Q] \textbf{Dust Market Cut} --- A share of the auctioneer’s commission for a season; wealth in sin.
  \hfill \textit{The cut is paid in mismatched coins. Spend them fast before they curse you.}
\item[K] \textbf{Blood‑Price Bond} --- A criminal clan owes you a life; they will kill once for you or die trying.
  \hfill \textit{The bond is a braid of enemy hair. Cut it to call the debt.}
\item[A] \textbf{Seal of the Oasis Compact} --- Casting vote on any water right dispute. Once.
  \hfill \textit{The seal is a bronze well with a snake curled at the bottom. Turn it and the water booms.}
\end{enumerate}

\section*{Quick use notes}
\label{sec:midh-ahkaz-quick-use}
\begin{itemize}
\item Draw until you have all four suits: Spade = place, Heart = actor, Club = pressure, Diamond = leverage. Highest rank sets main Clock (2–5 → 4, 6–10 → 6, J/Q/K → 8, A → 10).
\item Diamonds are narrative permits (grazing rights, oaths, witnesses) that shift Position rather than demand a roll.
\item If any A appears, echo \textbf{grass and wind} motifs—shifting alliances, debts that ride with the seasons, the sky as witness.
\item The violet seed‑heads of the steppe are a recurring omen: when they glow under moonlight, a spirit walks; when they shed, a caravan will be late.
\end{itemize}

\subsection*{Additional Features}
\begin{itemize}
\item \textbf{Kon’reh Arbitration:} Disputes are settled on a board painted on hide. Winners gain seasonal rights. Losers pay in grass or labor.
\item \textbf{Grass Omens:} The steppe flickers violet when rain is near; when the seed‑heads droop at noon, a blood‑feud approaches.
\item \textbf{Khatun’s Braid:} The chieftain’s hair holds the memory of every oath she swore. Cutting a braid voids a promise.
\end{itemize}

\subsection*{Lore Echoes (Cross‑Expansion Hooks)}
\begin{itemize}
  \item \textbf{The Bitter Root:} Daughters of the Cord run a “free soup” wagon near the Dust Market. They ask no coin—but they measure your shadow.
  \item \textbf{The Threshold Compendium:} Saikou Ira once tracked a ghoul to the Under‑Arcades. He left a silver mirror in the wall. It still shows the ghoul’s face.
  \item \textbf{Malachai, the Chained Angel:} A broken Kon’reh board at the Iron Well is said to compel truth. Players who touch it gain +1 die to resist lies—but mark 1 Obligation to Malachai.
  \item \textbf{The Grey Wanderer’s Grimoire:} Midh Ahkaz psions are called “Wind‑Readers.” They sense lies by the tremor in grass. The Khatun employs two; they never speak to each other.
  \item \textbf{The Ninth:} In any Midh Ahkaz scene, the ninth horse table in the Violet Bazaar sells nothing but bad advice. Spend a Boon to reroll a failed negotiation—but the advice may haunt you.
\end{itemize}

\subsection*{Midh Ahkaz Superstitions (burn for +1 die once)}
\begin{itemize}
  \item \textbf{Hoofprint Pour:} Pour a cup of water into the first hoofprint of a new horse. Ignore the first \emph{Horse Sickness} penalty.
  \item \textbf{Braid the Knot:} Tie a strand of mane hair around your wrist. Cancel \emph{Counterfeit Coin} or take –1 die on your next Charisma roll (your choice).
  \item \textbf{Name to the Grass:} Whisper a debtor’s name into the wind at noon. Gain +1 Effect when collecting from them this scene, but mark \emph{Obligation} to a Sky‑Speaker.
\end{itemize}

\begin{tcolorbox}[colback=black!3,colframe=black!40!white,title={Patronage \& Power in Midh Ahkaz}]
In Midh Ahkaz, power flows through horseflesh, water rights, and the geometry of oaths. The Khatun rules by consensus, but the true authority rests with the Oasis Compact—a loose council of clan elders, Vilikari factors, and Steppe Rangers who interpret the Wind Law.

\textbf{For the GM – Actionable Examples:}
\begin{itemize}
\item A player who wins a Kon’reh match earns a \textbf{Grass Right} (Diamond). That right could be contested by a rival clan the very next session—forcing a rematch or a chase across the prairie.
\item A \textbf{Water Share} token is stolen. The thief leaves a violet seed‑head as a signature. Tracking them leads to a seasonal camp where debt is repaid not in coin but in stories.
\item The Khatun’s illness (Club) is not natural. A Vilikari factor (Heart) knows the poisoner’s name but demands a \textbf{Black Ledger Entry} (J♠ Diamond) in exchange—a dangerous witness the player must then protect.
\end{itemize}
These hooks turn the generator entries into immediate playable situations.
\end{tcolorbox}

\index{Midh Ahkaz|see{Violet Gate}}
\index{Theme generator}
\index{Steppe generator}
\index{Places!Midh Ahkaz}
\index{People!Midh Ahkaz}
\index{Complications!Midh Ahkaz}
\index{Rewards!Midh Ahkaz}
\index{Khatun of the Ring}
\index{Violet Bazaar}
\index{Coin-Weigh Tribunal}
\index{Oasis Well}
\index{Steppe Ranger}
\index{Grass Ring}
\index{Kon'reh arbitration}
\index{Sky-Speaker}
\index{Iron Well}
\index{Tokan}

\section*{Thematic SB Spend Table (Midh Ahkaz)}
\label{sec:midh-ahkaz-sb}

\subsection*{Minor Complications (1 SB)}
\begin{itemize}
\item \textbf{Exposure:} Your actions draw attention from \textbf{Steppe Rangers or Coin‑Weigh clerks}.
\item \textbf{Noise:} The sound carries across the grass; a \textbf{rival clan’s scout} overhears.
\item \textbf{Trace:} Hoofprints, broken grass, or losing an arrow marks your path.
\item \textbf{Delay:} A water ration takes longer; a horse tires without rest.
\item \textbf{Grass Strain:} Mark +1 segment on a relevant \textbf{resource clock}.
\end{itemize}

\subsection*{Moderate Setbacks (2 SB)}
\begin{itemize}
\item \textbf{Alarm Raised:} \textbf{Border riders} converge on your position; a horn sounds.
\item \textbf{Position Lost:} You lose high ground/cover/the wind due to a \textbf{grass fire or horse stampede}.
\item \textbf{Foe Appears:} A \textbf{Vengeance broker, Ostrikari sell‑sword, or clan champion} arrives.
\item \textbf{Gear Trouble:} A \textbf{bridle breaks, bow string snaps, waterskin leaks}.
\item \textbf{Oath Hearsay:} A false rumor about your oath reaches the wrong ear.
\end{itemize}

\subsection*{Serious Trouble (3 SB)}
\begin{itemize}
\item \textbf{Reinforcements:} Additional \textbf{horsemen, blood‑price hunters, or Dust Market thugs} arrive.
\item \textbf{Key Gear Breaks:} A \textbf{horse goes lame, seal ring cracks, trade ledger burns}.
\item \textbf{Major Twist:} The situation changes – \textbf{Khatun dies, well runs dry, duel interrupted by grass fire}.
\item \textbf{Rail Tick:} Advance a relevant campaign/front clock by 1 segment.
\item \textbf{Condition Applied:} Mark \textbf{Fatigue 1/Harm 1/Condition} (e.g., Blinded by dust, Parched).
\end{itemize}

\subsection*{Major Turns (4+ SB)}
\begin{itemize}
\item \textbf{Trap Springs:} A \textbf{hidden pit, poisoned water share, or false oath} activates.
\item \textbf{Authority Arrival:} \textbf{Khatun’s warband, Ecktorian patrol, Sky‑Speaker} intervenes.
\item \textbf{Scene Shift:} The environment changes – \textbf{grass fire sweeps through, dust storm blinds, night falls without warning}.
\item \textbf{Patron Omen:} A \textbf{bolide, strange animal, or wind that speaks your name} appears.
\item \textbf{Narrative Pivot:} The story takes an unexpected turn – \textbf{the real forger was your ally; the water right was a trap.}
\end{itemize}

\subsection*{Region-Specific SB Options}
\begin{itemize}
\item \textbf{Midh Ahkaz (Grass \& Trade):} Grass catches fire, horses spook at shadows, bargains settled with a knife on the table, violet seed‑heads floating in the wind.
\item \textbf{Midh Ahkaz (Oath \& Sky):} Clouds spell forgotten names, wind carries false testimony, a duel’s outcome shifts with the sun, a violet dusk signals a spirit.
\item \textbf{Midh Ahkaz (Horse \& Bow):} An arrow splits in flight, a horse refuses a rider, a bowstring sings a dying clansman’s tune, a braid comes undone at the wrong moment.
\end{itemize}

\subsection*{Boss Seeds (three‑phase foes)}
\begin{itemize}
\item \textbf{The Counting Circle (cursed scale bound to the Oasis)} — \emph{Phases}: Fair Auditing → Retroactive Drought → Naming of Debtors. \emph{Cracks}: burn a counterfeit token in the Oasis well; settle a seven‑clan water debt; let a Sky‑Speaker rebless the scale.
\item \textbf{Grass Reaver (vengeful horse‑lord, half‑legend)} — \emph{Phases}: Hoof Drum → Ghost Herd → Blood‑Price Ride. \emph{Cracks}: confess a broken oath at his cairn; return a stolen foal; outrun him on a borrowed horse.
\item \textbf{Coin‑Weigh Wraith (greedy merchant who sold his shadow)} — \emph{Phases}: Tally of the Lost → Ledger Chain → Crack the Seal. \emph{Cracks}: pay a forgotten wergild; break his contract with a witnessed error; let the Iron Well drink his name.
\end{itemize}
\newpage

% ------------------------------
% CHAPTER: THE WAYS BETWEEN
% ------------------------------
\chapter{The Ways Between}
\emph{Spiritways \& Veilways.}

\noindent
{\Large
\textit{You are reading this sentence twice. The first time, you understood it. The second time, it understood you.}

\vspace{0.5em}

I am a Lost Pilgrim. Or I was. Or I will be. The paths here don't care about tenses. They care about which version of you arrives at the end. I've been walking for---no. That's a lie. The walking has been walking me. I forget which of us started first.

\vspace{0.5em}

The Ways Between do not measure miles. They measure choices. A wrong turn might lead to yesterday. A right turn might lead to a version of yourself who made a different decision and is still alive to regret it. The mist does not hide landmarks. It hides destinations. And sometimes, when you walk long enough, you realize the destination was never the point. The point was the moment you turned left instead of right. The point was the breath you held at the crossroads. The point was the name you did not speak.

\vspace{0.5em}

Here is the secret the Wayfinder doesn't want you to know: there is no Wayfinder. There is only the road, and the road is a liar. It tells you it wants you to find your way. What it actually wants is to watch you choose. Every reflection in every pool is a version of you that chose differently. They are all still walking. They are all still watching. And they are all still waiting for you to make the choice that finally lets them stop.

\vspace{0.5em}

I met myself on a bridge once. They were coming from the opposite direction. We stopped. We looked at each other. We said nothing. We both knew that if one of us spoke, the other would have to listen. And listening is a kind of debt.

\vspace{0.5em}

Never answer a question asked by a reflection. The answer will bind you to the path you did not choose. But here is the second secret: not answering is also a choice. And the road is very, very patient.

\vspace{1em}

\noindent\textit{--- The Wayfinder, who exists only in the moment a choice is made, and who is not the same person at the end of this sentence as at the beginning}
}

\vspace{1em}

\begin{almanac}
\textbf{What do people eat for breakfast?} Nothing. There is no breakfast in the Ways Between. There is no morning. There is no hunger. There is only the walking. The walking is sustenance.

\textbf{What does a child learn before they can walk?} To forget. Children born in the Ways Between do not learn. They remember. They remember the path they took to be born. The path has eight turns. The ninth turn is the birth. The birth is a choice. The child does not know they chose.

\textbf{What is the first thing you notice about a stranger?} Their direction. A stranger walking toward you may be coming from your future. A stranger walking away may be going to your past. A stranger walking beside you may be yourself from a timeline where you did not turn left. Do not ask them which. They will not answer.

\textbf{What is the most common crime?} Stealing---but stealing is not a crime, it is a paradox. A thief in the Ways Between steals a memory from a version of you that no longer exists. The memory never belonged to you. You cannot miss it. You will miss it anyway.

\textbf{What does no one talk about but everyone knows?} The Ways Between are not between anything. They are the only thing that exists. The solid ground you remember---the cities, the roads, the hearths---is the dream. The Ways Between are the waking. The waking is the walking. The walking never ends.

\textbf{Where do people go when they die?} Deeper. The Ways Between have no surface, no bottom, no end. The dead walk deeper. The deeper they walk, the less they remember. When they remember nothing, they become the path. The path is the walking. The walking is the dead.

\textbf{How do you apologize for a serious wrong?} You cannot. Apologies require memory. Memory requires a self. The Ways Between erase both. The only apology is to stop walking. Stopping is not an option.

\textbf{What is the most important festival?} There are no festivals. There is only the moment between heartbeats when you realize you have been walking for nine years and you do not remember why. That moment is the festival. Celebrate it by taking the next step.

\textbf{What is the secret that could destroy a powerful person?} The Wayfinder is not a person. It is the road itself, dreaming of a traveler who will finally reach the end. The end does not exist. The road knows this. The road dreams anyway.
\end{almanac}

\newpage
\section{The Ways Between −−− Obishaal, the Realm of Dreams and Death}
łabel{chap:ways−between}

\subsection*{Quick Reference}

\textbf{Tagline:} \textit{Every dream is a door. Every door is a judgment. The dead dream forever; the living borrow their sleep.}

\begin{quote}
  \textbf{The Wayfinder (Elite):}  
  ``I exist only in the moment of choice, the breath between one path and another. Obishaal is not a place—it is a \emph{state}. To guide here is to understand that every dream leads somewhere, but only the virtuous know which somewhere is worth the waking.''
\end{quote}

\begin{quote}
  \textbf{Lost Pilgrim (Commoner):}  
  ``The dead don't leave. They just change addresses. Obishaal is their new neighborhood, and the scales are balanced with virtue. If you're lucky, you'll only visit. If you're not... you'll find yourself climbing stairs that never end, and every landing offers a bargain you can't refuse.''
\end{quote}

\noindent\textbf{Genre / Mood:} Dream‑death metaphysical, virtue‑based ascent, temptation and fall, folk‑surrealism.

\noindent\textbf{Starting Location:} A crossroad shrine of black glass stones, where travelers meet their reflections coming the other way. The shrine is silent except for the soft hum of the stones. A Backwards Pilgrim stands at the crossroads, facing the direction you came from. ``The Thirteenth Stone speaks,'' they say. ``But the path chooses who arrives. What did you bury on the road here? The stones can smell it. And in Obishaal, every buried thing dreams of waking.''

\vspace{2mm}
\begin{tcolorbox}[colback=black!3,colframe=blue!60!black,title={Theme & Atmosphere}]
The Ways Between is a transit way, the eternal crossroads through Obishaal, the element of Dreams and Death, the eighth element after Fire, Life, Air, Fortune, Earth, Fate, Water, Dreams/Death. It is the realm where the settled dead go not to oblivion but to an \emph{ideal afterlife}—a dream that reflects the life they lived. A farmer dreams of endless harvest; a warrior dreams of eternal battle; a merchant dreams of scales that never tip. These dreams are real, and they are paradise or prisons.

The living may enter the Ways through thresholds of sleep, near‑death, or deliberate ritual. Once inside, they find a landscape that shifts with their own psyche—but also with the weight of their choices. \textbf{To ascend} through the layers of Obishaal, a traveler must demonstrate \emph{virtue}: honesty, courage, mercy, self‑sacrifice. \textbf{To fall} is to yield to \emph{vice}: greed, cowardice, cruelty, despair. Each layer is a test, and each test is a dream that tempts you to fail.

At the very bottom, pressed into the silt of the deepest dream, lies \textbf{Ikasha}, She Who Sleeps. Her sisters Inaea and Isoka imprisoned her there eons ago, for Ikasha dreamed of summoning and confronting the Clockwork Monad—a confrontation that, in their judgment, would unmake reality. Now Ikasha dreams fitfully, and her nightmares seep upward, infecting the layers above. Sometimes a traveler who falls far enough can hear her voice, offering escape in exchange for service—a service that would free her.

Between the layers drift the dead who have not yet found their final dream. They are the \emph{Unsettled}, and they will trade memories, secrets, and virtues for a chance to ascend—or to pull you down with them.

\vspace{2mm}
\textbf{What you notice first:} The air tastes of copper and old keys. Your footsteps echo half a beat too late. The waystones are warm—and they pulse like hearts.

\textbf{The one rule that kills newcomers:} Never answer a question asked by a reflection. The answer will bind you to the path you did not choose.
\end{tcolorbox}

\subsection*{Regional Mood: The Dream That Remembers}

Here, the dream is not a thing—it is a witness. It remembers every promise made at a deathbed, every lie told to a dying ear, every virtue you failed to keep. Obishaal does not forgive, but it does not forget either. It offers passage in exchange for truth. The longer you stay, the more the dream learns about you. By the time you reach your destination, you may find that the destination knows you better than you know yourself. Some travelers never leave. They become dream‑echoes, their faces pressed into the glass of a forgotten dream, their names still warm on the lips of the next pilgrim.

\begin{tcolorbox}[colback=black!3,colframe=red!60!black,title={The Ninth Taboo Manifestation}]
On any path in the Ways Between, the ninth step you take in a straight line will land on a threshold you did not see. The wise take eight steps, then pause to look around. In Obishaal, the ninth step is always a dream of falling. Those who fall hear a voice ask, ``What did you owe?'' Answering costs a memory; silence costs a name.
\end{tcolorbox}

% −−−−−−−−−−−−−−−−−−−−−−−−−−−−−−−−−−−−−−−−−−−−−−−−−−−−−−−−−−−−−−−−−−−−−−
% Thin Places — Where the Veil Wears Thin
% −−−−−−−−−−−−−−−−−−−−−−−−−−−−−−−−−−−−−−−−−−−−−−−−−−−−−−−−−−−−−−−−−−−−−−

\subsection*{Thin Places — Where the Veil Wears Thin}

Not every threshold is a door. Some are a \emph{crack}—a place where the world forgot to finish its seam. The Pereshi call them \emph{breath‑holes}, where the mountain's exhalation smells of yesterday. The Tulkani knot a thread and leave it untied, marking where the Hollow's attention wanders. The Aeler do not name them at all. They simply count to eight and leave the ninth breath unspoken.

A \textbf{thin place} is a location where the boundary between the waking world and the Ways Between is unusually permeable. Crossing is not a matter of ritual or intent—it is a matter of \emph{being in the wrong place at the wrong angle}. A wrong step. A miscounted breath. A name spoken at the wrong hour.

⫕section*{When Does a Thin Place Open?}
The GM may declare a thin place when:
\begin{itemize}
ℹtem A character rolls a \textbf{critical failure} (all $1$s) on a navigation or survival test in a liminal region (Valewood, Mistlands, high passes, deep mines, ancient ruins).
ℹtem The party camps at a \textbf{crossroads, bridge, or ford} during a new moon, an eclipse, or the ninth hour.
ℹtem A character \textbf{breaks a taboo} (pours the ninth cup, takes the ninth step, speaks the ninth name) in a region where the Hollow is strong.
ℹtem The GM spends \textbf{3+ Story Beats} to introduce a supernatural complication, and the fiction supports a threshold crossing.
ℹtem A player \textbf{willingly seeks} a thin place (rare—most who seek them never return unchanged).
\end{itemize}

⫕section*{Signs of a Thin Place (d6)}
Roll or choose when the party enters an area that might be thin:
\begin{enumerate}
ℹtem \textbf{Echoes answer wrong.} Your voice returns half a heartbeat later, or in a different pitch. Sometimes the echo speaks words you did not say.
ℹtem \textbf{Shadows drift.} Your shadow moves when you stand still. Cast by a light source that isn't there.
ℹtem \textbf{The ninth step appears.} A stair, a bridge plank, a stone—where no ninth should exist. Step on it, and the ground shifts.
ℹtem \textbf{Bells ring without clappers.} A distant chime, a temple bell, a levee alarm. The sound comes from a direction you cannot find.
ℹtem \textbf{The air tastes of copper and old keys.} The same taste Saikou Ira describes before every exorcism. The same taste the Aeler name \emph{breath of the deep}.
ℹtem \textbf{A door stands where no door stood.} Grey wood, warm handle, no frame. It does not open for those who knock. It opens for those who do not.
\end{enumerate}

⫕section*{Accidental Crossing (Mechanics)}
When a character is in a thin place and the conditions align, the GM may call for a \textbf{Spirit + Resolve test} (DV $4$–$6$, depending on the thinness). On a failure, the character crosses into the Ways Between—not by choice, but by \emph{accident}.

\begin{itemize}
ℹtem \textbf{Partial success ($0 < \text{successes} < \text{DV}$):} The character feels the pull. They may choose to step back (mark $1$ Fatigue) or let themselves cross (gain $1$ Boon, but enter the Ways).
ℹtem \textbf{Miss ($0$ successes):} The character crosses involuntarily. They arrive in a random location in the Ways Between (roll on the table in §\ref{sec:ways−between−thin−destinations}).
ℹtem \textbf{Critical failure (all $1$s):} The character crosses—and something crosses \emph{with} them. A Oath‑Keeper, a Counting Wraith, or worse.
\end{itemize}

Once crossed, the character must navigate the Ways Between using the standard travel procedures (Chapter 46). They cannot simply “turn back”—the threshold that opened may close behind them. Finding the way home requires a \emph{balance}: a story, a memory, or a name offered to the road.

⫕section*{Known Thin Places (By Spades Draw)}
The following locations are known thin places, keyed to the \emph{Spades} (place) draws of various regional generators. A GM who draws one of these cards while the party is in that region may declare the area thin, triggering the possibility of an accidental crossing.

\begin{longtable}{p{0.18\textwidth} p{0.32\textwidth} p{0.40\textwidth}}
∩tion{Known Thin Places by Region and Spades Rank}\\
ℎline
\textbf{Region} & \textbf{Spades Entry (Place)} & \textbf{Thin Place Hazard / Hook} \\
ℎline
\endfirsthead
ℎline
\textbf{Region} & \textbf{Spades Entry (Place)} & \textbf{Thin Place Hazard / Hook} \\
ℎline
\endhead
ℎline
μlticolumn{3}{r}{\textit{Continued on next page}} \\
\endfoot
ℎline
\endlastfoot

Aeler & ♠10 (Reliquary Arcade) & \textit{The Bell‑Crypts.} Bells ring without clappers. A ninth tone pulls listeners into a dream‑echo of the mountain's first collapse. \\
Mistlands & ♠Q (Witchlight Bridge) & \textit{The Ninth Plank.} Step where the plank is missing, and you step into the fog of Obishaal. The Hollow Bride sometimes waits on the far side. \\
Ykrul & ♠A (Sky Steppe) & \textit{The Star‑Road Scar.} Where the grass grows in spirals. Crossing at dusk may lead to the Long Mile of Temptation. \\
Vilikari & ♠A (Foedus Stone) & \textit{The Crack of the Ninth Clause.} Touch the stone's newest crack, and you may hear the Oath‑Keeper's question before you answer. \\
Linn & ♠K (High Jarl's Seat) & \textit{The Beacon That Remembers.} When the beacon burns green, the dead can see you. When it burns blue, they can touch you. \\
Thepyrgos & ♠A (The Unfinished Stair) & \textit{The Ninth Flight.} Climb where no stair should be, and you climb through the fossil stairs of Obishaal. \\
Acasia & ♠A (The Pale Causeway) & \textit{The Causeway of Broken Promises.} Every milestone is a broken oath. Walk it at dusk, and the Walking Scale may find you. \\
Vhasia & ♠K (Sun Palace, Lence) & \textit{The Hall of Shattered Mirrors.} Mirrors show not your face, but a reflection of the path you refused. Step through, and you enter the Crossroads of the Ninth Choice. \\
Viterra & ♠A (Highway Stone) & \textit{The Queen's Milestone.} The stone is cracked. The crack is a door. Those who peer inside see the Bone Corridor of Unfinished Business. \\
Valewood & ♠A (Valeheart Spire) & \textit{The Spire's Shadow.} When the spire phases, its shadow becomes a stair. Climb it, and you enter the Stone Circle of Forgotten Villages. \\
Aelinnel & ♠A (The Green Gate) & \textit{The Gate That Expects Change.} If you cannot pay the toll in truths, the gate may open to the Memory Tunnel of the Dying instead of the forest. \\
Aelaerem & ♠A (Hollow Field) & \textit{The Field of Uncounted Steps.} Step nine times in a straight line, and the Hollow Field becomes the Mist‑Ford of First Doubt. \\
Zakov & ♠A (Serpent's Spine) & \textit{The Reef That Sings.} The harmonic hum is the song of the Deep Drake. Follow it, and you may find the Pressure's lair—or the Thirteenth Stone. \\
Tulkani & ♠A (Ikasha's Dreaming Place) & \textit{The Dreamer's Threshold.} A cave that smells of honey and rust. Entering it is already a crossing. The Wayfinder waits inside. \\
Fhara & ♠A (The First Well) & \textit{The Well of Forgotten Names.} Drink from it, and you may wake in the Waystation of the Virtuous, your name already spent. \\
Sidhi & ♠A (The Sealed Coast) & \textit{The Coast of Unmarked Graves.} The drowned walk here. If you answer a drowned sailor's question, you follow them into the Promise Bridge of the Unsettled. \\
Pereshi & ♠A (The First Keystone) & \textit{The Stone That Asks.} Touch it, and the mountain asks a question. Answer poorly, and you climb the Fossil Stairs of Ascension. \\
Kuvani & ♠K (Khagan's Felt Throne) & \textit{The Throne of Tails.} The saddle has no stirrups. Sit in it, and the Long Mile of Temptation begins beneath you. \\
Ashaan & ♠A (Ikasha's Sleeping Chamber) & \textit{The Sealed Door.} The stone is warm. The lock is a whisper. Those who listen too long cross into the deepest layer—and may never wake. \\
Ngomebe & ♠A (The First City) & \textit{The City That Never Stopped.} A walking city seen only on the horizon. Step onto its shadow, and you step into the Spiral Path of Recursion. \\
Taharka & ♠A (The King's Tomb) & \textit{The Tomb of the Unnamed.} The keystone is a door. The door leads to the Threshold Arch of Virtue. The arch judges. \\
Sekogo & ♠A (The Ukwe's Root‑Chamber) & \textit{The Root That Whispers.} Sit on the throne of roots, and the tree asks one question. If you cannot answer, you fall into the Bone Corridor. \\
Lethai‑ar (Crimson Valley) & ♠A (The First Bridge) & \textit{The Invisible Strand.} Only those who have witnessed nine true oaths can see it. Cross it, and you walk the Promise Bridge of the Unsettled. \\
Theona & ♠A (Lookout Cliff) & \textit{The Cliff of Nine Views.} From here you can see the Green Host riding—and sometimes the Fossil Stairs rising beside them. \\
Oshiira & ♠A (The Ledger Vault) & \textit{The Vault of Unpaid Scales.} The records are warm. Read the ninth page, and the Thirteenth Stone speaks to you. \\
Black River (Thepyrgos–Viterra border) & (Special) & \textit{The Ferryman's Slip.} The river is a threshold. Bryn's boat is not the only one. On moonless nights, the dead row themselves. A coin in the water may buy a crossing—or sell your name. \\
\end{longtable}

⫕section*{Denizens of the Ways Between (Expanded)}
The following creatures may appear in the Ways Between, drawn from the bestiaries and expansions. They are not random encounters—they are \emph{judges}. Each has a hunger, a weakness, and a price.

\begin{longtable}{p{0.20\textwidth} p{0.20\textwidth} p{0.20\textwidth} p{0.30\textwidth}}
∩tion{Denizens of the Ways Between}\\
ℎline
\textbf{Creature} & \textbf{Source} & \textbf{Hunger} & \textbf{Price to Banish} \\
ℎline
\endfirsthead
ℎline
\textbf{Creature} & \textbf{Source} & \textbf{Hunger} & \textbf{Price to Banish} \\
ℎline
\endhead
ℎline
μlticolumn{4}{r}{\textit{Continued on next page}} \\
\endfoot
ℎline
\endlastfoot
Counting Wraith & \textit{Witnessed Prey} & Errors (ninth step, ninth breath) & A correct count of eight \\
Oath‑Keeper & \textit{Witnessed Prey} & Unpaid scales & Payment of the original weight \\
Trow & \textit{Witnessed Prey} & Absences, omissions & A memory freely given \\
Neighbor & \textit{Witnessed Prey} & Broken taboos, attention & A never‑told truth (or silence) \\
Poltergeist (Class IV) & \textit{Saikou Ira's Compendium} & Rage, unfinished business & A witnessed apology \\
Deep Drake (Pressure) & \textit{Witnessed Prey / Saikou Ira} & Attention, the ninth breath & A memory of sunlight (or a copper bell) \\
Scale Lich & \textit{Witnessed Prey (appendix)} & Unbalanced judgments & A true name carved in quartz \\
The Hollow Bride & \textit{The Hearth Balance} & Hospitality broken, the ninth cup & An apology witnessed \\
The Walking Scale & \textit{The Salt Balance} & Broken promises & A story never told \\
The Keeper of the Seventh Exhalation & \textit{The Wandering Balance} & The seventh sneeze & A name offered freely \\
\end{longtable}

⫕section*{Accidental Crossing Destination (d6)}
łabel{sec:ways−between−thin−destinations}
When a character crosses into the Ways Between accidentally, roll to determine where they emerge:
\begin{enumerate}
ℹtem \textbf{The Mist−Ford of First Doubt} (Class III Procession site). Water runs uphill. Footsteps echo backwards. The dead march here at midnight.
ℹtem \textbf{The Bone Corridor of Unfinished Business} (Anchor graveyard). Shadows walk the walls but never touch the floor. A Poltergeist may be forming.
ℹtem \textbf{The Crossroads of the Ninth Choice} (Threshold nexus). All four paths lead to the same destination—but different \emph{yesterdays}. A Trow sits at the center, smoking.
ℹtem \textbf{The Promise Bridge of the Unsettled} (Oath‑Keeper territory). Cross with a vow, and the bridge remembers. The toll is a story you have never told.
ℹtem \textbf{The Long Mile of Temptation} (Vice's testing ground). The road stretches differently for each traveler. A Walking Scale appears at the halfway mark.
ℹtem \textbf{The Thirteenth Stone (Ikasha's Seal)}. The path ends at a warm stone. The stone whispers. Beneath it, She Who Sleeps dreams of freedom. A Deep Drake may be coiled nearby.
\end{enumerate}

\begin{tcolorbox}[colback=black!3,colframe=black!40!white,title={GM Guidance: Thin Places as Pressure, Not Punishment}]
Thin places are not traps. They are \emph{complications}. A party that accidentally crosses into the Ways Between should not feel that they have “failed.” They have simply entered a different kind of scene—one where negotiation matters more than steel.

\begin{itemize}
ℹtem \textbf{Offer a way out.} The Ways Between always offer a path home. The path may demand a toll (a memory, a story, a name), but it exists.
ℹtem \textbf{Use the travel rules.} Treat the Ways Between as a leg of a journey. Draw cards. Set timers. The destination is “home” or “the next threshold.”
ℹtem \textbf{Let players use their bonds.} A character who calls on a bond while lost in the Ways may hear a familiar voice—or see a familiar face—guiding them. The road remembers who weighs what.
ℹtem \textbf{The ninth is always watching.} In any scene set in the Ways Between, the ninth card drawn (or the ninth roll made) triggers a manifestation of the Hollow. Not a fight—a \emph{question}. The players must answer truthfully, or the road takes a memory.
\end{itemize}
\end{tcolorbox}

\vspace{2mm}
\noindent\textit{Set down by the Grey Wanderer, after a night spent in a thin place near the Black River. The fire went out three times. The shadows did not. The ninth bell rang once—and then was silent.}

\noindent\textit{The Ways Between are not a dungeon. They are a conversation. Speak carefully. The road listens.}

\section*{Spades −−− Places (Dream Layers)}
łabel{sec:ways−between−places}
\begin{enumerate}
\setcounter{enumi}{1}
ℹtem \textbf{Mist−Ford of First Doubt} −−− Water runs uphill; footsteps echo backwards through time. `[THRESHOLD]`
  ℎfill \textit{The ford shows you wading through your own past. Do not look down at the water.}
ℹtem \textbf{Bone Corridor of Unfinished Business} −−− Shadows walk the walls but never touch the floor; light as memory. `[DEATH]`
  ℎfill \textit{The bones are from those who chose vice and fell. They shift when you blink.}
ℹtem \textbf{Threshold Arch of Virtue} −−− Shows a reflection as someone you might have been; the past as a mirror. `[TRUTH]`
  ℎfill \textit{The reflection has a scar you never got. It asks how you avoided it.}
ℹtem \textbf{Spiral Path of Recursion} −−− Ascends but returns to the same stone marker; the journey as meditation. `[MAGIC]`
  ℎfill \textit{Each loop, the stone moves closer to the center. Or you do.}
ℹtem \textbf{Promise Bridge of the Unsettled} −−− Cross with a vow and it remembers your weight; an oath as a toll. `[OATH]`
  ℎfill \textit{The bridge's planks are old letters. Read one, and you may owe the writer.}
ℹtem \textbf{Crossroads of the Ninth Choice} −−− All four paths lead to the same destination, but different tomorrows. `[THRESHOLD]`
  ℎfill \textit{The center stone is warm. A coin placed there will vanish—or be accepted.}
ℹtem \textbf{Stone Circle of Forgotten Villages} −−− Marks where a village used to be; the houses are now constellations. `[DEATH]`
  ℎfill \textit{The stars are too close. You can touch them. They crumble like ash.}
ℹtem \textbf{Memory Tunnel of the Dying} −−− Hear the thoughts of those who passed here; the past as an echo. `[MEMORY]`
  ℎfill \textit{The thoughts are in a language you almost understand. One word is your name.}
ℹtem \textbf{Fossil Stairs of Ascension} −−− Each step is an eye that watches your ascent; the climb as scrutiny. `[FEAR]`
  ℎfill \textit{The eyes belong to those who fell here. They blink in sequence.}
ℹtem[J] \textbf{Waystation of the Virtuous} −−− The keeper trades in unfinished conversations; words as currency. `[SOCIAL]`
  ℎfill \textit{The keeper has no face. It wears the faces of those who paused too long.}
ℹtem[Q] \textbf{Junction of Vice and Virtue} −−− Paths physically collide; you must choose which reality to follow. `[THRESHOLD]`
  ℎfill \textit{The collision sounds like breaking glass. The glass reforms into a door.}
ℹtem[K] \textbf{Long Mile of Temptation} −−− The road stretches differently for each traveler; distance as perception. `[MAGIC]`
  ℎfill \textit{You are the only one walking. The others are already at the end.}
ℹtem[A] \textbf{The Thirteenth Stone −−− Ikasha's Seal} −−− Where the path reveals its true destination; truth as revelation. `[TRUTH]`
  ℎfill \textit{The stone whispers. The whisper is a promise you made to yourself and broke. Beneath it, the stone is warm—and dreaming.}
\end{enumerate}

\paragraph*{(Wayfarer/Spirit/Dream−walker)} Lost pilgrim who thinks they're going home but forgot where; toll‑taker accepting payment in forgotten memories.

\section*{Hearts −−− Travelers & Guides}
łabel{sec:ways−between−people}
\begin{enumerate}
\setcounter{enumi}{1}
ℹtem \textbf{Lost Pilgrim} −−− Thinks they're going home but forgot where; destination as a memory. `[TRAVEL]`
  ℎfill \textit{Their pack is full of stones. Each stone is a forgotten promise.}
ℹtem \textbf{Toll−Taker (Virtue Assessor)} −−− Accepts payment in virtues you didn't know you had. `[SOCIAL]`
  ℎfill \textit{Their scales weigh truth and echo. A balanced scale means you pass—for now.}
ℹtem \textbf{Wayward Guide} −−− Directions technically correct, morally questionable. `[TRAVEL]`
  ℎfill \textit{They point with a broken compass. The needle spins, but the path is true.}
ℹtem \textbf{Dream−Merchant} −−− Carries a sack of nightmares and one perfect dream for sale. `[MAGIC]`
  ℎfill \textit{The perfect dream is their own. They are selling it because they cannot wake up.}
ℹtem \textbf{Child−Ghost of Unfinished Life} −−− Knows shortcuts but charges in riddles; wisdom as a puzzle. `[DEATH]`
  ℎfill \textit{The riddles have no answer—only a choice. Choose wrong, and you become a child.}
ℹtem \textbf{Wounded Walker (Vice−Bound)} −−− Bleeding metaphor; their pain grows thorns along the path. `[FEAR]`
  ℎfill \textit{The thorns have faces. Each face is someone they failed to save.}
ℹtem \textbf{Map−Merchant of Impossible Dreams} −−− Selling maps to places that don't exist yet; the future as a commodity. `[TRUTH]`
  ℎfill \textit{The maps change as you watch. The destination is always one page ahead.}
ℹtem \textbf{Guide−Dog (Ancestral Spirit)} −−− Made of shadow and starlight; follows those who walk with purpose. `[SPIRIT]`
  ℎfill \textit{It never barks. It whines when you are about to make a mistake.}
ℹtem \textbf{Backwards Pilgrim (Dream−Echo)} −−− Moving through time; their steps erase what just happened. `[MEMORY]`
  ℎfill \textit{They are walking toward their birth. Every step, they forget something.}
ℹtem[J] \textbf{Soul−Ferryman} −−− Carries the luggage of the living; a burden as a service. `[DEATH]`
  ℎfill \textit{The luggage is full of things you never said. The ferryman knows them all.}
ℹtem[Q] \textbf{Road's Child (Dream−Kin)} −−− Born where paths crossed, raised by waymarks; navigation as instinct. `[TRAVEL]`
  ℎfill \textit{They have no navel. They were never born—they simply began.}
ℹtem[K] \textbf{Dead−Road Keeper} −−− Ensures proper passage for those who should not walk. `[DEATH]`
  ℎfill \textit{Their lantern burns a cold blue. The flame does not flicker.}
ℹtem[A] \textbf{The Wayfinder (Dream Avatar)} −−− Exists only in the moment a choice is made; a decision as an entity. `[LEGEND]`
  ℎfill \textit{They are every guide you ever needed. They are also every road you refused.}
\end{enumerate}

\paragraph*{(Veil‑Thin/Dream‑Bleed/Wayward)} Path loops back to show your funeral preparations; reality thins−−−you see the dreams of sleeping travelers; Ikasha's nightmare bleeds into the upper layers.

\section*{Clubs −−− Complications/Threats}
łabel{sec:ways−between−complications}
\begin{enumerate}
\setcounter{enumi}{1}
ℹtem \textbf{Funeral Loop} −−− The path shows your funeral preparations; death as a destination. `[DEATH]` `[FEAR]`
  ℎfill \textit{The mourners are people you have not met yet. They seem to know you well.}
ℹtem \textbf{Reality Thins (Ikasha's Dream Bleed)} −−− You see the dreams of sleeping travelers; the boundary as a veil. `[MAGIC]`
  ℎfill \textit{A dreamer mutters your name. They are dreaming of a choice you will make tomorrow.}
ℹtem \textbf{Waymark Lies (Vice's Test)} −−− Points the wrong way; trust leads to places that never were. `[TRUTH]`
  ℎfill \textit{The waymark has a crack. Through the crack, you see a path you should have taken.}
ℹtem \textbf{Dream−Bleed} −−− Waking memories become someone else's nightmares. `[MEMORY]` `[FEAR]`
  ℎfill \textit{You remember a childhood fear that was not yours. It still terrifies you.}
ℹtem \textbf{Toll Demanded (Virtue Payment)} −−− Currency you did not know you carried; payment as a revelation. `[SOCIAL]` `[WEIGHT]`
  ℎfill \textit{The toll is a secret you told a dead person. They are now standing behind you.}
ℹtem \textbf{Path Split (Virtue vs. Vice)} −−− Each version of you remembers differently; identity as a choice. `[MAGIC]`
  ℎfill \textit{One you is smiling. One you is weeping. One you is already walking away.}
ℹtem \textbf{Gravity Shifts (Guilt's Weight)} −−− Emotional weight grows physically heavy; a burden as physics. `[WEATHER]`
  ℎfill \textit{Your pack feels heavier. Inside is a regret you forgot you carried.}
ℹtem \textbf{Time−Sickness (Dream Chronology)} −−− You arrive before you left; paradox shadows follow. `[MAGIC]`
  ℎfill \textit{Your past self waves at you. Do not wave back.}
ℹtem \textbf{Truth Payment (Ikasha's Price)} −−− The road remembers your lies and demands truth as payment. `[TRUTH]` `[WEIGHT]`
  ℎfill \textit{The lie is carved on a stone. The stone is warm. It knows your temperature.}
ℹtem[J] \textbf{Crossroads Judgment (Virtue Audit)} −−− The path you chose judges the ones you did not; decision as karma. `[TRUTH]` `[SOCIAL]`
  ℎfill \textit{A ghost of the road not taken asks why. You have one sentence to answer.}
ℹtem[Q] \textbf{Memory−Thief (Unsettled Dead)} −−− A rest stop keeper pays in counterfeit recollections. `[MEMORY]` `[SOCIAL]`
  ℎfill \textit{The memory is almost right. The color of the sky is wrong.}
ℹtem[K] \textbf{Forbidden Path (Ikasha's Descent)} −−− Opens only for those already lost; loss as a key. `[DEATH]`
  ℎfill \textit{The path smells of old tears. You taste salt on your lips.}
ℹtem[A] \textbf{Convergence (Ikasha's Waking)} −−− All travelers arrive together, none recall how; arrival as a mystery. `[MAGIC]` `[LEGEND]`
  ℎfill \textit{You recognize everyone. You have never seen any of them before. The air tastes of honey and rust. Ikasha is dreaming of you.}
\end{enumerate}

\paragraph*{(True Name/Safe Passage/Waywisdom)} Waymark that always points toward your next important choice; a token of passage that the road recognizes as belonging.

\section*{Diamonds −−− Rewards/Leverage}
łabel{sec:ways−between−rewards}
\begin{enumerate}
\setcounter{enumi}{1}
ℹtem \textbf{Choice Compass (Virtue's Needle)} −−− Always points toward your next important choice. `[TRAVEL]`
  ℎfill \textit{The needle is a splinter of the Thirteenth Stone. It hums when you are near a decision.}
ℹtem \textbf{Passage Token (Dream−Marked)} −−− The road recognizes you as one who belongs. `[SOCIAL]`
  ℎfill \textit{The token is a braid of your own hair. You do not remember cutting it.}
ℹtem \textbf{Dream−Catcher} −−− Filters nightmares from your rest; sleep as a sanctuary. `[MAGIC]`
  ℎfill \textit{The catcher is woven from old promises. It smells of lavender and regret.}
ℹtem \textbf{Truth−Compass (Ikasha's Fragment)} −−− Points to what you most need to know, not what you want. `[TRUTH]`
  ℎfill \textit{The compass has no needle. The housing spins until you speak a truth.}
ℹtem \textbf{Memory−Anchor} −−− Keeps you from losing yourself on deeper paths. `[MEMORY]`
  ℎfill \textit{The anchor is a river stone. Your name is carved on it in a language you have never seen.}
ℹtem \textbf{Safe Haven (Virtue's Rest)} −−− A guaranteed rest stop without complication. `[TRAVEL]`
  ℎfill \textit{The haven is a dry arch. The arch is a doorframe with no door.}
ℹtem \textbf{Guide−Light (Virtue's Beacon)} −−− Burns the color of your truest intention; purpose as a beacon. `[SOCIAL]`
  ℎfill \textit{The light is blue. Blue means you are honest. Red means you are not.}
ℹtem \textbf{Path−Cutter (Dream Shears)} −−− Cuts distance through metaphor, not space. `[MAGIC]`
  ℎfill \textit{The cutter is a pair of shears. Cut a memory, and the road between here and there vanishes.}
ℹtem \textbf{Scale−Clearing (Virtue's Mercy)} −−− The road forgives one weight you thought you carried. `[WEIGHT]`
  ℎfill \textit{The weight was owed to a dead person. They are no longer dead on this path.}
ℹtem[J] \textbf{Crossroads Boon} −−− Choose among three paths, each exactly where you need to be. `[LEGEND]`
  ℎfill \textit{The paths are labeled: Regret, Forgive, Forget. Choose carefully.}
ℹtem[Q] \textbf{Way−Wisdom (Dream−Logic)} −−− Understand the language of paths and signs for one journey. `[TRAVEL]`
  ℎfill \textit{The waystones whisper. They are complaining about the last traveler. It was you.}
ℹtem[K] \textbf{Grace Passage (Virtue's Shield)} −−− Walk safely through any dangerous crossing for one night. `[TRAVEL]`
  ℎfill \textit{The grace is a bell. Ring it, and the road will pretend not to see you.}
ℹtem[A] \textbf{Road's Name (Ikasha's True Name)} −−− Call the path by its true name and command its nature. `[LEGEND]`
  ℎfill \textit{The name is a single syllable. Speaking it will change the road—and you.}
\end{enumerate}

\section*{Quick use notes}
łabel{sec:ways−between−quick−use}
\begin{itemize}
ℹtem Draw until you have all four suits: Spade = place, Heart = actor, Club = pressure, Diamond = leverage. Highest rank sets the main Timer ($2$–$5 → 4$, $6$–$10 → 6$, J/Q/K $→ 8$, A $→ 10$).
ℹtem Diamonds are codified outcomes (true names/safe passages/waywisdom) that change position rather than call for a roll.
ℹtem If any A appears, echo \textbf{dream‑road} motifs—reflections that lie, paths that judge, destinations that choose the traveler. Also, the presence of an Ace hints that Ikasha's dream is particularly close.
\end{itemize}

\subsection*{Additional Features}
łabel{sec:ways−between−features}
\begin{itemize}
ℹtem \textbf{Dream−Logic Navigation:} Any Ace introduces a metaphysical requirement (a sacrifice, a taboo, an exchange) that must be satisfied to pass. Name it at the table; the road will enforce it.
ℹtem \textbf{Reflection Points:} Face cards (J, Q, K) show alternate selves or unlived choices. Treat them as temporary NPCs or scene tags that can help, hinder, or tempt.
ℹtem \textbf{Memory Currency:} Diamonds may be traded for knowledge or safe passage as if they were memories. Describe the memory paid; the path (or its keeper) pays in kind.
ℹtem \textbf{Virtue and Vice Timers:} Each character in Obishaal tracks a small personal \textbf{Virtue} [3] and \textbf{Vice} [3] timer. Virtue advances when they act with honesty, courage, mercy, or self‑sacrifice. Vice advances when they act with greed, cowardice, cruelty, or despair. When one timer fills, the character \textbf{ascends} (Virtue) to a higher layer of Obishaal or \textbf{falls} (Vice) to a lower one. Ascension offers a boon; falling offers a temptation (a cursed gift from the layer's vice‑spirit). The timers reset after each layer transition. Reaching the bottom (Vice filled in the deepest layer) allows Ikasha to speak directly to the character.
ℹtem \textbf{Ikasha's Whisper (Patron Intervention):} When a character falls to the lowest layer (or deliberately seeks her out), they may hear Ikasha's voice. She offers a \emph{Fragment of Freedom}: $+2$ dice to any roll involving escape, secrecy, or breaking bonds, but each use marks $1$ \textbf{Obligation to Ikasha} and advances a \textbf{Waking Dream} [6] timer. When that timer fills, Ikasha uses the character as a vessel to attempt an escape—a campaign‑level event.
ℹtem \textbf{Local Patrons (Syncretic Names):}
  \begin{itemize}
    ℹtem \textbf{The Unnamed Guide} — The Traveler. \textit{The road that appears only when you are lost.}
    ℹtem \textbf{The Mirror Listener} — The Witness. \textit{The reflection that tells you what you need to hear, never what you want.}
    ℹtem \textbf{The Threshold Child} — The Sealed Gate. \textit{The one who was never born, who guards the places where paths end.}
    ℹtem \textbf{The Echo King} — The Ninth. \textit{The voice that answers when you speak to the empty road.}
    ℹtem \textbf{The Dreamer Below} — Ikasha (imprisoned). \textit{She who sleeps and dreams of waking. Her whispers are bargains; her freedom is a promise.}
  \end{itemize}
\end{itemize}

\subsection*{Lore Echoes (Cross‑Expansion Hooks)}
\begin{itemize}
  ℹtem \textbf{The Bitter Root:} Daughters of the Cord know the Ways Between as the Cord's shadow. A lost Hollowed wanders here, its thread trailing into the mist. Follow it, and you may find a Sister who cut herself free—and who now tends Ikasha's dream.
  ℹtem \textbf{The Threshold Compendium:} Saikou Ira's silver mirror shows the Ways Between as a tangled knot of unfulfilled balances. He has walked three paths here. He will not speak of the fourth, where he glimpsed Ikasha's face.
  ℹtem \textbf{Malachai, the Chained Angel:} A silver coin that never spends lies at every crossroads. Pick it up, and the road will offer you a shortcut—but the coin will reappear in your pocket at the next junction, heavier. In Obishaal, the coin is warm and pulses like a heartbeat.
  ℹtem \textbf{The Grey Wanderer's Grimoire:} Psions call the Ways Between ``the Stair.'' They climb it in their dreams, trading memories for glimpses of paths not yet walked. Some never wake. The Grey Wanderer claims that the Stair is Ikasha's spine.
\end{itemize}

\subsection*{Ways Between Superstitions (burn for $+1$ die once)}
\begin{itemize}
  ℹtem \textbf{Bread to the Stones:} Crumble a crust at the base of a waystone. Ignore the first \emph{Waymark Lies} penalty. `[SUPERSTITION]`
  ℹtem \textbf{Knot the Path:} Tie a single knot in a piece of string and hang it from a milestone. Cancel \emph{Path Split} or take $−−1$ die on your next navigation roll (your choice). `[SUPERSTITION]`
  ℹtem \textbf{Name to the Crossroads:} Whisper a dead traveler's name into the center stone. Gain $+1$ Effect when bargaining with a Toll‑Taker this scene, but mark \emph{Obligation} to the Wayfinder. `[SUPERSTITION]`
\end{itemize}

\subsection*{Thematic SB Spend Table}
łabel{sec:ways−between−sb}

\paragraph{Minor Complications (1 SB)}
\begin{itemize}
ℹtem \textbf{Exposure:} Your actions draw unwanted attention from \textbf{path keepers or wayward spirits}.
ℹtem \textbf{Noise:} Sounds of your actions alert nearby \textbf{travelers or memory echoes}.
ℹtem \textbf{Trace:} Evidence of your passage marks your route for \textbf{trackers or reflection points}.
ℹtem \textbf{Delay:} A brief but meaningful setback costs you \textbf{time or favorable path alignment}.
ℹtem \textbf{Supply Strain:} Mark $+1$ segment on a relevant \textbf{resource timer}.
\end{itemize}

\paragraph{Moderate Setbacks (2 SB)}
\begin{itemize}
ℹtem \textbf{Alarm Raised:} \textbf{Toll‑taker or wayward guide} becomes aware and begins responding.
ℹtem \textbf{Position Lost:} You lose advantageous ground/cover/stealth due to \textbf{path shift or reality thinning}.
ℹtem \textbf{Foe Appears:} A \textbf{memory‑thief or forbidden path entity} arrives on scene.
ℹtem \textbf{Gear Trouble:} A piece of equipment becomes \textbf{Compromised/Neglected}.
ℹtem \textbf{Lock/Barrier:} A simple obstacle now requires a test to overcome.
\end{itemize}

\paragraph{Serious Trouble (3 SB)}
\begin{itemize}
ℹtem \textbf{Reinforcements:} Additional \textbf{wayward spirits, path entities, or alternate selves} arrive.
ℹtem \textbf{Key Gear Breaks:} A crucial tool/weapon becomes temporarily unusable.
ℹtem \textbf{Major Twist:} The situation fundamentally changes − \textbf{path loops/time shifts/reality fractures}.
ℹtem \textbf{Rail Tick:} Advance a relevant campaign/front timer by $1$ segment.
ℹtem \textbf{Condition Applied:} Mark \textbf{Fatigue 1/Harm 1/Condition} appropriate to fiction.
\end{itemize}

\paragraph{Major Turns (4+ SB)}
\begin{itemize}
ℹtem \textbf{Trap Springs:} A prepared danger activates with full effect.
ℹtem \textbf{Authority Arrival:} \textbf{The Wayfinder, Dead‑Road Keeper, or Road's Child} intervenes.
ℹtem \textbf{Scene Shift:} The environment changes dramatically − \textbf{reality thins/paths converge/time loops}.
ℹtem \textbf{Patron Omen:} Divine/arcane forces take notice − \textbf{omen appears/blessing lost/curse manifests}.
ℹtem \textbf{Narrative Pivot:} The story takes an unexpected turn that reframes objectives.
\end{itemize}

\paragraph{Region−Specific SB Options}
\begin{itemize}
ℹtem \textbf{Obishaal (Dream‑Logic):} Paths rearrange without warning, reflections speak unbidden, destinations shift mid‑journey. A character may hear Ikasha's whisper as a free omen (the GM may spend $1$ SB to make it a false omen).
ℹtem \textbf{Obishaal (Memory Currency):} Memories become tangible, forgotten knowledge resurfaces, past and present blur. Trading a memory here may also tick the \emph{Virtue} or \emph{Vice} timer depending on the memory's nature.
ℹtem \textbf{Obishaal (Virtue and Vice):} Decisions echo forward, alternate selves appear, paths judge moral weight. The GM may spend SB to force a character to test a virtue (e.g., ``Resist the temptation to lie'') or to offer a vice that ticks their Vice timer.
\end{itemize}

\subsection*{Path‑Moves & Procedures (Optional Depth)}
łabel{sec:ways−between−procedures}

\paragraph{Weighing of Steps (Travel Move)}
When the company enters a liminal path (Mist‑Ford, Bone Corridor, etc.), choose one:
\begin{itemize}
  ℹtem \textbf{Pay the Toll:} Each PC names a small truth or a fond memory (GM: $1$ SB banked, $+1$ Position on navigation).
  ℹtem \textbf{Refuse the Toll:} $+1$d on Speed, but start a \textbf{Way‑Scale} [4] timer.
  ℹtem \textbf{Bargain the Toll:} Trade a \emph{Diamond} reward as if it were a memory (see Memory Exchange below); cancel one imminent \emph{Club} card in this scene.
\end{itemize}
When \textbf{Way‑Scale} fills, trigger \emph{Crossroads Judgment} at the next junction.

\paragraph{Reflection Duel (Social/Moral Test)}
When confronted by an alternate self or \emph{Reflection Point}, state what that self wants from you and roll your best social action.
\begin{itemize}
  ℹtem \textbf{Success:} Gain a \emph{Guide‑Light} tag for the scene and clear $1$ tick from \textbf{Paths Decide}.
  ℹtem \textbf{Mixed:} Trade places for one exchange or take a \textbf{Condition: Disoriented}.
  ℹtem \textbf{Miss:} The reflection walks away with one of your \emph{Diamonds} (GM picks the most narratively apt).
\end{itemize}

\paragraph{Memory Exchange (Rates & Risks)}
\begin{tabular}{p{0.36\textwidth} p{0.28\textwidth} p{0.30\textwidth}}
\textbf{Memory Offered} & \textbf{Buys You} & \textbf{Side Effect} \\
ℎline
A childhood smell, a forgotten nickname & $+1$ Position on one crossing & A familiar place feels strange for a scene \\
A first failure or small betrayal & Cancel one minor Complication & You gain a \emph{tell} (the Ways recognize you) \\
A true fear you have never spoken & \emph{Guide‑Light} or \emph{Way‑Cord}‑like effect & A dream pursues you for $1$–$3$ nights \\
The details of a kept promise & \emph{Safe Passage} for the party & An NPC tied to that promise misremembers you \\
The meaning behind a scar & Convert one \emph{Club} into \emph{Trace} only & The scar “moves” (cosmetic) until the next dawn \\
A perfect day (all details) & Treat the next \emph{Heart} face as an ally & Start \textbf{Dream‑Bleed} [3] on the map \\
\end{tabular}

\medskip
\textbf{Exchange Rule:} If you pay the same \emph{kind} of memory twice in a journey, start \textbf{Hollowing} [4]. When it fills, you gain the \emph{Backwards Pilgrim} tag until you reclaim what was lost.

\subsection*{Way‑Keys & Taboos (d6)}
\begin{enumerate}
ℹtem \textbf{Speak Last}: Only the final word said at a gate is true; others are “practice.” First speakers take $−−1$d on their next test.
ℹtem \textbf{Step Odd}: Crossings require an odd number of steps; even counts loop you back one scene.
ℹtem \textbf{Name Nothing}: True names wake \emph{Name‑Theft}; use titles or lose a \emph{Diamond}.
ℹtem \textbf{Wet Iron}: A drop of water on iron cancels one glamour once per scene.
ℹtem \textbf{Back Gift}: Leave something behind to go forward; refuse and tick \textbf{Way‑Scale}.
ℹtem \textbf{Answer Bent}: Answer a waystone's question with a metaphor, not a fact, or trigger \emph{Truth Payment}.
\end{enumerate}

\subsection*{Boss Seeds (three‑phase foes)}
\begin{itemize}
ℹtem \textbf{The Archivist of Unmade Roads (cloak stitched from surveyor's cords)} −−− \emph{Phases}: Receipts of Refusal $→$ Stay of Journey $→$ Unfiled Options. \emph{Cracks}: a \textbf{Choice Compass} reveals the true path; a \textbf{Scale‑Clearing} voids his claim.
ℹtem \textbf{The Toll‑Matriarch Beneath the Bridge (many hands, one voice like wet stone)} −−− \emph{Phases}: Truth Payment $→$ Scale‑Clearing $→$ Forbidden Path. \emph{Cracks}: a \textbf{Passage Token} marks you as belonging; a \textbf{Road's Name} commands her to stand aside.
ℹtem \textbf{The Wayfinder's Shadow (your silhouette, two steps ahead)} −−− \emph{Phases}: Crossroads Judgment $→$ Guide‑Light $→$ Haste Mark. \emph{Cracks}: a \textbf{Memory‑Anchor} keeps you grounded; a \textbf{Grace Passage} slips past its gaze.
ℹtem \textbf{Ikasha's Nightmare (Dream−Eater)} −−− Appears only in the deepest layer. \emph{Phases}: Whispers $→$ Temptation $→$ Consumption. \emph{Cracks}: a \textbf{Virtue} filled timer can be spent to reject the nightmare; a memory of true self‑sacrifice (cost $1$ permanent Bond) banishes it entirely.
\end{itemize}

\subsection*{Encounter Seeds (quick hooks)}
\begin{itemize}
ℹtem \textbf{The Promise Bridge's Receipt:} An old vow you never spoke was logged here. To cross, define when you would have sworn it−−−and to whom.
ℹtem \textbf{Thirteenth Stone Auction:} The stone sells three possible endings. Each bidder is you from a different failure.
ℹtem \textbf{Bone Corridor Choir:} Shadows sing a name wrong by one letter. Correct it and gain \emph{Truth‑Compass}; ignore it and start \textbf{Name‑Theft}.
ℹtem \textbf{Ikasha's Finger Bone:} A character finds a bone that pulses with dream‑light. If they hold it while sleeping, they can speak to Ikasha directly—but each time, the \textbf{Waking Dream} timer ticks.
\end{itemize}

% Extensions (Plug−in)
\subsection*{Extensions (Plug−in)}
łabel{sec:ways−between−extensions}

\paragraph{Choice Dial (Resolve $↔$ Drift)}
Track the party's footing between firm intent and wandering.
\begin{itemize}
  ℹtem \emph{Resolve High:} $+1$ Position on \emph{Path Split} tests; the first \emph{Dream‑Bleed} each session hits harder (start at $2$).
  ℹtem \emph{Drift High:} $+1$ Effect when bargaining with way‑entities; the first \emph{Truth Payment} each session costs a cherished \emph{detail}.
  ℹtem Center the Dial by naming aloud the next choice \emph{and} why you might refuse it; clear $1$ tick from \textbf{Way‑Scale}.
\end{itemize}

\paragraph{Counter‑Seal (once/scene, if justified)}
Break a single metaphysical instruction (taboo, key, or omen) by presenting a memory in symbol form (a token, a letter, a sketch). Gain $+1$d on that crossing; the GM banks $1$ SB as a \emph{Marked Consequence} to spend when reflections return.

\paragraph{Start‑of‑Scene Murmurs (1d6)}
\begin{enumerate}
  ℹtem \emph{Second Sunrise:} light arrives twice; your first missed navigation roll becomes a mixed success.
  ℹtem \emph{Name Draft:} the air steals a syllable; your next lie sounds true once.
  ℹtem \emph{Echo Scale:} your footprints owe a step—skip a square or tick \textbf{Way‑Scale} $+1$.
  ℹtem \emph{Pale Bell:} time hiccups; the first \emph{Club} can be delayed until the end of the scene.
  ℹtem \emph{Traveler's Chorus:} distant voices agree with you; $+1$d on appeals to path entities.
  ℹtem \emph{Hollow Weather:} rain that does not wet; conditions look dire but impose no penalties this scene.
\end{enumerate}

\paragraph{Plug‑in SB Conversions (Obishaal Flavor)}
\begin{itemize}
  ℹtem $1$ SB $→$ \emph{Footprint Vote}: your tracks choose a different direction; increase DV by $1$ on any roll that relies on following a trail.
  ℹtem $2$ SB $→$ \emph{Echo of Refusal}: a choice you avoided manifests as a minor obstacle (a locked gate, a forgotten toll). If the choice involved a vice, also tick the Vice timer of the character who refused.
  ℹtem $3+$ SB $→$ \emph{Threshold Cracks}: the path fractures; gain a temporary \textbf{Path‑Cutter} effect but start \textbf{Convergence} [4] (all travelers you have met will arrive together at the next junction). Additionally, Ikasha's dream seeps through the cracks; each character hears a single word of her voice.
\end{itemize}

\paragraph{Cross‑Regional Conversions (Using Diamonds Across the Map)}
\begin{itemize}
  ℹtem \emph{Valewood:} \textbf{Road's Name} may count as a \emph{Valeheart Clause} once if spoken under leaves; doing so creates \textbf{Court Tithe} [2].
  ℹtem \emph{Mistlands:} \textbf{Grace Passage} functions as \emph{Pall Indulgence} over one levee crown at night; ring a bell afterward or start \textbf{Witchlight Count}.
  ℹtem \emph{Vhasia:} \textbf{Way‑Wisdom} can be presented as “pilgrim right” to bypass one \emph{Interdict} barrier; begin a \textbf{Parlement Divided} tick.
  ℹtem \emph{Viterra:} \textbf{Passage Token} can be notarized into a temporary \emph{Ferry Priority}; on use, advance \textbf{Border‑Lace}.
  ℹtem \emph{Ubral:} \textbf{Memory‑Anchor} counts as a \emph{Guest‑Token} at one hearth if you recount a family road‑tale.
  ℹtem \emph{Silkstrand:} \textbf{Path‑Cutter} doubles as a night‑boat route once; spend $1$ favor with a smuggler or trigger \emph{Seal Counterfeit}.
  ℹtem \emph{Linn:} \textbf{Choice Compass} acts as a \emph{Pilot's Token} in a reed‑maze if you pay a name to the water.
  ℹtem \emph{Theona:} \textbf{Crossroads Boon} allows counting “eight‑and‑one” to cross an \emph{Uncounted Bridge} safely once.
\end{itemize}

\begin{tcolorbox}[colback=black!3,colframe=black!40!white,title={Patronage & Power}]
In Obishaal, power flows through understanding the metaphysical nature of dreams, the ability to navigate virtue and vice, and the careful management of one's weight before the scales. The true authorities are those who know the True Names of dream‑layers, can mediate between the living and the Unsettled, and understand the currency of virtue.

\textbf{For the GM:}  
Patronage in Obishaal revolves around guidance, safe passage, and the ability to navigate moral tests. Rewards often take the form of dream‑tokens, fragments of Ikasha's memory, or knowledge of a layer's True Name. To emphasize this:
\begin{itemize}
ℹtem Tie rewards to visible symbols (compasses, anchors, bells) that can be challenged, stolen, or voided by the dream itself.
ℹtem Let rival guides (Virtue vs. Vice) issue conflicting directions, forcing players to choose whose favor matters more.
ℹtem Use the crossroads, waystations, and dream‑layers as arenas for moral contests, where a single act of mercy can ascend you—and a single act of cruelty can damn you.
\end{itemize}
In Obishaal, your virtue is your power, and your power determines whether you wake or sleep forever.
\end{tcolorbox}

% Index entries
∈dex{Ways Between|see{Obishaal}}
∈dex{Obishaal}
∈dex{Ikasha}
∈dex{Dreams and Death}
∈dex{Virtue and Vice}
∈dex{Theme generator}
∈dex{Spiritways and Veilways generator}
∈dex{Places!Obishaal}
∈dex{People!Obishaal}
∈dex{Complications!Obishaal}
∈dex{Rewards!Obishaal}
∈dex{The Wayfinder}
∈dex{Dead−Road Keeper}
∈dex{Road's Child}
∈dex{Soul−Ferryman}
∈dex{Toll−Taker}
∈dex{Thirteenth Stone}
∈dex{Long Mile}
∈dex{Promise Bridge}
∈dex{Crossroads}
∈dex{Memory Tunnel}
∈dex{Fossil Stairs}
∈dex{Choice Compass}
∈dex{Passage Token}
∈dex{Dream−Catcher}
∈dex{Truth−Compass}
∈dex{Memory−Anchor}
∈dex{Road's Name}
∈dex{Grace Passage}
∈dex{Local Patrons!Obishaal}
\newpage

% ------------------------------
% CHAPTER: THE WILDS
% ------------------------------
\chapter{The Wilds}
\emph{Roads, Ruins, and Weather.}

\noindent
{\Large
\textit{Between towns and treaties lies the patience of the land. Old roads remember armies; cairns remember names no book keeps. In the Wilds, law thins to trail-marks and favors, and the weather keeps its own counsel. The Edgewalker calls it ''unclaimed inventory.`` The land calls it patience. The land has more patience than the Edgewalker.}

\vspace{0.5em}

I am a ranger now. Stonehaven trained me as a scholar---mapping ruins, copying old ledgers, believing that ink could preserve what the world forgot. Tam taught me that the world does not forget. It waits. And while it waits, it watches. The wilds don't need your roads or your rules; they've got their own ways of keeping score. Follow the cairns if you want to live, and always leave something behind---a coin, a story, a breath held just long enough to be noticed. Every hill hides a secret, and every hollow breathes.

\vspace{0.5em}

The land has a ledger too. It just doesn't use ink. Between the cartographers' claims, the wilds endure---uninterested in borders, indifferent to tariffs. Here, rangers read spoor and weather; hermits keep shrines at forgotten fords; beasts and older things observe travelers with eyes that have seen empires rise and fall. Travel is a negotiation with the land itself, and the land always remembers a slight. The hollows have long memories. They are not cruel. They are not kind. They are patient.

\vspace{0.5em}

Never sleep in a hollow that has soot on the ceiling without adding your own to the fire. The land knows the difference between a guest and a thief. It also knows the difference between a thief and a corpse. It prefers the corpse. Less arguing.

\vspace{1em}

\noindent\textit{--- Moira of Stonehaven, ranger and reluctant debt-keeper}
}

\vspace{1em}

\begin{almanac}
\textbf{What do people eat for breakfast?} Whatever they found---berries, roots, a rabbit snared at dawn. The wilds provide. The wilds also take. The taking is breakfast for something else.

\textbf{What does a child learn before they can walk?} To read the sky. Wilds children are taught the eight winds by taste: North is dry, South is wet, East is cold, West is warm, and the ninth---the wind that has no direction because it is the land's own breath. The land's breath is a warning.

\textbf{What is the first thing you notice about a stranger?} Their pack. A stranger whose pack is heavy has been traveling long. A stranger whose pack is light has just started. A stranger with no pack is either a fool or a ghost. The wilds eat fools. Ghosts eat the wilds.

\textbf{What is the most common crime?} Forgetting---but forgetting is not a crime, it is a survival strategy. The wilds are full of things that want to be forgotten. A traveler who forgets them may pass safely. A traveler who remembers must pay a toll. The toll is a story. The story is the memory. The memory is the trap.

\textbf{What does no one talk about but everyone knows?} The wilds are not empty. They are full of watchers---spirits, beasts, old things that do not have names because names are debts. The watchers do not interfere. They watch. They wait. They remember every footprint. When the footprints stop, the watchers move on.

\textbf{Where do people go when they die?} Into the land. The dead are not buried. They are left where they fell. The land absorbs them. The land grows a little richer. The land's memory grows a little longer. The dead become the land's memory. The memory is the wilds.

\textbf{How do you apologize for a serious wrong?} You leave an offering---a coin, a button, a lock of hair---at the place where the wrong was done. The land accepts or refuses. If the offering is gone by morning, you are forgiven. If it remains, you are not. If it has been moved, the land is testing you. The test is the walk. The walk is the apology.

\textbf{What is the most important festival?} The Turning of the Leaves (autumn). The wilds change color. The colors are the land's mood. Red is anger. Gold is warning. Brown is patience. Green is forgetfulness. This year, the leaves are red. The land is angry. Do not ask why.

\textbf{What is the secret that could destroy a powerful person?} The wilds have a heart. It is a stone, buried beneath the oldest oak. The stone is warm. The stone is beating. The stone is the land's pulse. If the stone stops, the wilds die. If the wilds die, there is nothing left to remember. Nothing left to wait.
\end{almanac}

\newpage
% ============================================================
% The Wilds — Roads, Ruins, and Weather
% Enhanced with evocative generators, land−responsive mechanics
% ============================================================
\section{The Wilds −−− ``Roads, Ruins, and Weather''}
łabel{chap:wilderness}

\subsection*{Quick Reference}
\textbf{Tagline:} \textit{The land walks when you don't watch it, and the roads remember who bled on them.}

\begin{quote}
  \textbf{Claimant Chief (Elite):}  
  ``These lands answer not to charters but to those who can read their signs and respect their moods. Authority here is earned one campfire at a time, by proving you understand the difference between claiming and belonging.''
\end{quote}

\begin{quote}
  \textbf{Forager Child (Commoner):}  
  ``The wilds don't need your roads or your rules −− they got their own ways of keeping score. Follow the cairns if you want to live, and always leave something behind to thank the land for letting you pass. Every hill hides a secret, and every hollow breathes.''
\end{quote}

\textbf{Genre / Mood:} Hexcrawl wanderlust, ruin‑haunted silence, weather as character.

\textbf{Starting Location:} A shelter hollow carved into a fallen giant's ribcage, moss glowing faintly under starlight. A Forager Child offers a route‑song: ``The Bad Ground shifts, and the weather walks behind it. Sing with me, or walk alone.''

\vspace{2mm}
\begin{tcolorbox}[colback=black!3,colframe=blue!60!black,title={Theme & Atmosphere}]
Between towns and treaties lies the patience of the land. Old roads remember armies; cairns remember names no book keeps. In the Wilds, law thins to trail‑marks and favors, and the weather keeps its own counsel. This is the country of the hexcrawl: a patchwork of hidden valleys, forgotten ruins, beast lairs, and the bones of failed settlements. No single authority rules here—instead, the land itself is the measure, and every crossing writes a new line. Walk soft, read the wind, bargain with whatever still listens, and keep one eye on the horizon. The next hex may hold water, treasure, or teeth.

\vspace{2mm}
\textbf{What you notice first:} The silence after a bird stops singing. The way a cairn seems to have shifted since you passed it last. The smell of rain on dry stone, even when the sky is clear.

\textbf{The one rule that kills newcomers:} Never sleep in a hollow that has soot on the ceiling without adding your own to the fire. The land knows the difference between a guest and a thief.
\end{tcolorbox}

\subsection*{Regional Mood: The Land That Watches}

Here, the wilds are not empty—they are patient. Every ridge has been climbed by someone who never came down; every spring has been drunk from by a hundred generations of the lost. The land remembers every footstep, every campfire, every promise whispered into the wind. It does not judge, but it does not forget. Power in the wilds is not given by crowns or charters—it is earned by reading the weather, by leaving the right offering at the right cairn, by knowing when to move and when to wait. Those who rush are swallowed; those who listen are guided. The wilds are not a destination. They are the conversation between the map and the untamed.

\begin{tcolorbox}[colback=black!3,colframe=red!60!black,title={The Ninth Taboo Manifestation}]
In any Wilds scene, the ninth cairn you pass in a day will be warm to the touch. Leave an offering—a pinch of salt, a copper coin, a single word spoken to the wind—or the path will forget you. Those who forget to pay find that the next crossing leads back to the same cairn, and the cairn has grown a new stone. The stone has their name on it.
\end{tcolorbox}

\subsection*{The Unfinished Map}

The Wilds are not a route—they are a grid of secrets. Each day’s march opens a new page. Some hexes are empty: water, fuel, a safe place to sleep. Others hold lairs, ruins, shrines, or the remnants of old wars. The GM should key a handful of interesting places in advance, then let the cards populate the space between. Draw a card for each new leg; let the suit and rank whisper what you find. The land does not repeat itself, but it echoes.

\paragraph*{(Crossing/Lookout/Old Road)} Crossing point with ford/ice‑span/stepping logs/dune saddle/reef shelf; lookout knoll with wind‑carved marker; old road trace with cairns half‑eaten by terrain.

\section*{Spades −−− Places (Hex Terrain)}
łabel{sec:wilderness−places}
\begin{enumerate}
\setcounter{enumi}{1}
ℹtem \textbf{Crossing Point} −−− Ford, ice‑span, stepping logs, dune saddle, reef shelf.
  ℎfill \textit{The crossing has a price: a stone, a story, or a moment of silence. The water hums when it's satisfied.}
ℹtem \textbf{Lookout Knoll} −−− Tor, cliff, dune, ridge with wind‑carved marker.
  ℎfill \textit{From here, you can see three hexes in clear weather—and one that shouldn't exist, flickering like a heat‑shaped lie.}
ℹtem \textbf{Old Road} −−− Trace with cairns and switchbacks half‑eaten by terrain.
  ℎfill \textit{The cobbles are worn smooth by ghosts. They whisper directions at dusk. The directions are always true, but never complete.}
ℹtem \textbf{Shelter Hollow} −−− Overhang, cave, root cellar; soot says ``used lately.''
  ℎfill \textit{The soot is layered in thin, greasy rings. Someone has been coming here for a very long time. The hollow smells of wet fur and old secrets.}
ℹtem \textbf{Water Source} −−− Spring, seep, fog net, snow pan, guarded by thorns or stone.
  ℎfill \textit{The water tastes of iron and old bargains. Drink deep, but leave a coin—copper for a safe draw, silver for a question answered.}
ℹtem \textbf{Windbreak} −−− Rock ribs, lee of pines, reed‑wall berm.
  ℎfill \textit{The windbreak smells of woodsmoke and wool. Someone camped here last night. Their fire is still warm.}
ℹtem \textbf{Ruined Outpost} −−− Watchstack, wayside shrine, marker pile.
  ℎfill \textit{The ruin's flagstones are warm, as if a great beast sleeps beneath them. The air tastes of rust and thyme.}
ℹtem \textbf{Bad Ground} −−− Sinkhole, bog crust, crevasse, loess slump.
  ℎfill \textit{The ground breathes—a slow, deep inhalation that lifts your heels. Step lightly or it will swallow you.}
ℹtem \textbf{Gate Gully} −−− Pinch point between bluffs, dunes, or ice hummocks.
  ℎfill \textit{The gully has a name carved in the rock. The name changes every season. Carve yours, and you may never leave.}
ℹtem[J] \textbf{Boundary Row} −−− Totems, standing stones, prayer flags bent by weather.
  ℎfill \textit{The flags are faded. The last prayer was for rain—and it still hasn't come. The totems have grown moss in the shape of faces.}
ℹtem[Q] \textbf{Worksite} −−− Abandoned mine, quarry, logging camp, salt pan.
  ℎfill \textit{The tools are scattered, but none are rusted. Whoever left did not plan to return—or has not yet finished.}
ℹtem[K] \textbf{Signal Height} −−− Old fire‑pit or mirror stand; smoke stains linger.
  ℎfill \textit{Light the fire, and someone will answer. Whether they are friend or foe is your risk. The smoke always smells of your childhood.}
ℹtem[A] \textbf{Trail Nexus} −−− Migratory, pack, or contraband paths knot here.
  ℎfill \textit{Three trails meet, but only two are on any map. The third belongs to the dead. Their footprints point the wrong way.}
\end{enumerate}

\paragraph*{(Guide/Warden/Caravan)} Forager child with truer maps than yours; guide with three routes and one price (changes daily).

\section*{Hearts −−− People & Factions (Hex Denizens)}
łabel{sec:wilderness−people}
\begin{enumerate}
\setcounter{enumi}{1}
ℹtem \textbf{Forager Child} −−− Goat‑herd or berry‑picker with truer maps than yours.
  ℎfill \textit{Their map is drawn on a scrap of bark. It shows a path that disappears at a standing stone. The stone is not on any other map.}
ℹtem \textbf{Guide} −−− Three routes and one price (changes daily); navigation as negotiation.
  ℎfill \textit{Their price is never coin. It is a favor, a secret, or a night's watch. They will not tell you which until you arrive.}
ℹtem \textbf{Warden Patrol} −−− Local badges, local laws, local patience; order as familiarity.
  ℎfill \textit{Their badge is a chipped sunburst. They serve a crown that no longer exists. The crown does not pay them. The land does.}
ℹtem \textbf{Caravan Crew} −−− Drovers, porters, yarn‑post hands trading speed for coin.
  ℎfill \textit{They know every shortcut and every ambush. They will share both—for a price. Their silence is purchased in stories.}
ℹtem \textbf{Pilgrims} −−− Bound for a tucked‑away shrine or stone; faith as a destination.
  ℎfill \textit{Their shrine has no name. The stone there is warm, and it answers questions poorly. It prefers to ask them.}
ℹtem \textbf{Poachers} −−− Know every snare and shortcut; survival as expertise.
  ℎfill \textit{They wear the teeth of game wardens as necklaces. The game wardens are still alive. They are saving for a full set.}
ℹtem \textbf{Hermit−Healer} −−− Dogs, geese, and opinions about weather; wisdom as isolation.
  ℎfill \textit{Their poultices smell of grave‑moss. They work anyway. The dogs do not bark at strangers—they judge them.}
ℹtem \textbf{Prospectors} −−− Salt, amber, iron, fungus; chasing rumors and glint.
  ℎfill \textit{Their claim is marked with a knot of red thread. Cross it, and they will shoot. The thread glows when a stranger's foot touches it.}
ℹtem \textbf{War−Band} −−− Reavers or ``escorts,'' depending on your purse; violence as commerce.
  ℎfill \textit{Their banner is a painted skull. The paint is from a dead comrade's blood. They add a finger bone to the standard after every contract.}
ℹtem[J] \textbf{Monster−Hunter} −−− Rite‑keeper wearing yesterday's trophies; death as livelihood.
  ℎfill \textit{Their trophy is still warm. The monster's mate is close behind. They are not running—they are leading.}
ℹtem[Q] \textbf{Quartermaster} −−− Ledgers first, hospitality second; supply as authority.
  ℎfill \textit{Their measure is written in a code only they understand. The code is their life story. The margins are full of regrets.}
ℹtem[K] \textbf{Claimant Chief} −−− Papers and spears to match; legitimacy as force.
  ℎfill \textit{Their papers are stained with water and blood. Both are still wet. The water came from a spring they poisoned. The blood from the spring's keeper.}
ℹtem[A] \textbf{The Stranger} −−− Spirit‑touched nomad or emissary; rules bend near them.
  ℎfill \textit{They have no shadow. The ground does not remember their footprints. Their voice echoes twice—once for them, once for the land.}
\end{enumerate}

\subsection*{Additional Denizens of the Wilds}
łabel{sec:wilderness−extra−denizens}
Not every soul in the wilds is a wanderer. Some have built fragile refuges, while others have been broken by war, exile, or the land itself.

\begin{enumerate}
\setcounter{enumi}{13}
ℹtem \textbf{Isolated Town (Palings)} −−− A handful of families eking out a living behind a thorn hedge.
  ℎfill \textit{The gate is a wagon axle lashed upright. They trade hides for salt and ask no questions. Their children have never seen a paved road.} `[COMMUNITY]'
ℹtem \textbf{Refugee Train} −−− A column of worn carts and quieter people fleeing a war or a curse.
  ℎfill \textit{Their leader carries a charter from a lord who is now dead. They will trade stories for passage and bury their dead without markers.} `[SORROW]'
ℹtem \textbf{Deserter Camp} −−− Soldiers who threw down their colours, now living by stealth.
  ℎfill \textit{They still drill at dawn and keep a sharp watch. Their loyalty is to each other, not to any flag. They may know the location of an abandoned fort.} `[STEALTH]'
ℹtem \textbf{Beastman Hunting Band} −−− Feral‑looking folk with antler charms and pelts stitched to skin.
  ℎfill \textit{They do not speak Utaran. They communicate in growls and hand signs. They are not savages—they are the land's first children, and they remember when the roads were deer paths.} `[SPIRIT]'
ℹtem \textbf{Lost Tribe (Cairn−Folk)} −−− Descendants of a settlement that fell off every map.
  ℎfill \textit{They wear the teeth of their ancestors and tattoo their lineage on their faces. They believe the Hollow is a sleeping god and that the party are its messengers. They are wrong—or are they?} `[MYSTERY]'
ℹtem \textbf{Rootless Mercenaries} −−− A dozen sell‑swords between contracts, nursing old wounds.
  ℎfill \textit{Their captain lost an eye to a mushroom spore. They will fight for coin, but they will also fight for a hot meal. They know the location of a forgotten cache of arms.} `[COMBAT]'
ℹtem \textbf{Anchorite in a Dead Tower} −−− A holy woman who has not spoken in twenty years.
  ℎfill \textit{She writes prophecies on the walls with charcoal. The prophecies are always about the next storm or the next stranger to arrive. She will not look at you if your shadow is wrong.} `[FAITH]'
ℹtem \textbf{Feral Children} −−− Orphaned by a raid or a plague, raised by wolves or worse.
  ℎfill \textit{They know every hidden spring and every game trail. They will not come close, but they will watch. Leave them a piece of bread, and they may leave you a knapped arrowhead—or a warning.} `[SURVIVAL]'
ℹtem \textbf{Cult of the Unburied} −−− Grave‑tenders who believe the dead should not be left to the land.
  ℎfill \textit{They carry iron shovels and speak in hushed tones. They burn the bodies of the fallen they find, lest the Hollow take them. They may be willing to burn your dead as well—for a story.} `[RITUAL]'
ℹtem[J] \textbf{Exiled Scholar} −−− A academic who fled a city after proving the wrong theorem.
  ℎfill \textit{He carries a single book—his own—and will trade a chapter for a night's shelter. The book contains a map to a ruin that does not appear on any official chart.} `[KNOWLEDGE]'
ℹtem[Q] \textbf{The Grey Singer} −−− A blind woman who wanders the wilds humming.
  ℎfill \textit{Her songs calm territorial beasts and call down birds. She is looking for a melody she heard in a dream. If you help her find it, she will teach you a verse that quells a storm.} `[MYSTIC]'
ℹtem[K] \textbf{Bandit King's Lookout} −−− A scout for a large robber band, hiding in a hollow tree.
  ℎfill \textit{He has a whistle that can summon a dozen riders. He is not cruel—just hungry. A bag of salt or a true story will buy his silence and his guidance around the ambush.} `[STEALTH]'
ℹtem[A] \textbf{The Hollow's Own} −−− A figure dressed in grey rags, walking the same path every dusk.
  ℎfill \textit{They have no face—only a smooth expanse where features should be. They do not attack. They point. The direction they point is the direction the Hollow will walk next.} `[DEATH]'
\end{enumerate}

\paragraph*{(Weather/Doubleback/Prowlers)} Weather turn with heat snap/cold snap/fog/dust; doubleback with tracks that loop; prowlers shadow with wolves/jackals/ghouls.

\section*{Clubs −−− Complications/Threats (Hex Hazards)}
łabel{sec:wilderness−complications}
\begin{enumerate}
\setcounter{enumi}{1}
ℹtem \textbf{Weather Turn} −−− Heat snap, cold snap, fog, dust; plans sag like wet leather.
  ℎfill \textit{The wind changes direction three times in an hour. Someone is watching. The air tastes of a season that hasn't come yet.}
ℹtem \textbf{Doubleback} −−− Tracks loop; the navigator swears the land moved underfoot.
  ℎfill \textit{Your own footprints face you. They were made yesterday—or tomorrow. The toes point toward a future you have not chosen.}
ℹtem \textbf{Prowlers Shadow} −−− Wolves, jackals, ghouls, seals, ravens by night; hunt as a company.
  ℎfill \textit{Their eyes reflect no light. Their teeth do. They are patient. They have been following since the last hex. You are only now noticing.}
ℹtem \textbf{Route Blocked} −−− Deadfall, rockfall, dune shift, ice heave; the path as an obstacle.
  ℎfill \textit{The blockage was placed. The toolmarks are fresh. Whoever cleared this path did not want you to follow—or wanted you to take the other way.}
ℹtem \textbf{Quarantine Sign} −−− Camp fever; wardens sniff your packs for death.
  ℎfill \textit{The sign is a ribbon tied in a figure‑eight. Do not cross it. The ribbon is damp. It was tied recently. The wardens are still nearby.}
ℹtem \textbf{Territorial Beast} −−− Charge, stampede, or swarm; your choice of which is wrong.
  ℎfill \textit{The beast's lair is nearby. You can smell it—old blood and older musk. The ground is warm. You are standing on its pantry.}
ℹtem \textbf{Elemental Front} −−− Grassfire, peat‑burn, canopy flare, blowing spindrift.
  ℎfill \textit{The fire moves against the wind. The wind is lying. The smoke writes a word in a language you almost understand. It is a warning.}
ℹtem \textbf{Paper vs Spear} −−− Jurisdiction fight in the wilds stalls your day; law as argument.
  ℎfill \textit{Two officials, one road. Both have seals. Neither will show them. The seals are warm. They have been used recently—on death warrants.}
ℹtem \textbf{Supply Pinch} −−− Water, fuel, or feed low; pick what starves first.
  ℎfill \textit{Your canteen is full. The water is from a source that dried up last year. You do not remember filling it. Someone did.}
ℹtem[J] \textbf{Pursuit} −−− Hunters or avengers follow; signs say ``close.''
  ℎfill \textit{The tracks are one hour old. They are gaining. The hunters do not need to sleep. They have not slept for three days. Neither have you.}
ℹtem[Q] \textbf{Bad Omen} −−− Will‑lights, taboo day, saint bells silent; locals refuse to move.
  ℎfill \textit{The omen is written in bird droppings on a milestone. It is still accurate. The birds are watching. They are waiting for you to ignore it.}
ℹtem[K] \textbf{General Alarm} −−− Levy, muster, evacuation; all tracks become checkpoints.
  ℎfill \textit{The alarm horn is cracked. It sounds like a scream. The scream is a name. The name is yours.}
ℹtem[A] \textbf{Catastrophe} −−− Flood, whiteout, sandstorm, lahar; timers jump like frightened deer.
  ℎfill \textit{The catastrophe has a name. The land remembers it. The land is shouting it now. The sound is the same as a heartbeat. Yours.}
\end{enumerate}

\subsection*{Expanded Travel Hazards (from expansions)}
łabel{sec:wilderness−expanded−hazards}
For GMs who want even more variety, add these entries to the \emph{Clubs} table (either replace six existing entries or expand the table to d18 by renumbering).

\begin{enumerate}
\setcounter{enumi}{13}
ℹtem \textbf{Bell−Line Failure (Mistlands)} −−− A ward‑bell cracks. Wraiths cross the levee.
  ℎfill \textit{Test Spirit + Resolve (DV 4) or gain the \textsc{Marked by Mist} condition (the next spirit encounter starts at Desperate Position).} `[DEATH]'
ℹtem \textbf{Star−Road Phasing (Valewood)} −−− An ancient Lethai road flickers in and out of reality.
  ℎfill \textit{Navigation DV +2. On a failure, the party emerges in a different hex entirely (GM chooses).} `[MAGIC]'
ℹtem \textbf{Spirit−Dust Storm (Ashaan)} −−− Red sand from the Red Desert carries whispering spirits.
  ℎfill \textit{All actions Desperate for the scene. On any 1, a bound spirit breaks free and attacks (Cap 2).} `[SPIRIT]'
ℹtem \textbf{Ninth Night (Tulkani Circle)} −−− The party has camped too long. The Hollow walks.
  ℎfill \textit{An Oath‑Keeper appears at midnight. It asks for a story, a memory, or a promise. Refusal deepens the land's watchfulness.} `[ATTENTION]'
ℹtem \textbf{Well−Spirit Honored (Fhara)} −−− A dry well suddenly flows, but the water carries minerals.
  ℎfill \textit{The spirit asks for respect—an offering of dates or a true story. Granted, the water stays sweet; refused, it turns bitter for the next leg.} `[SPIRIT]'
ℹtem \textbf{Toll Shift (Silkstrand/Brass Coast)} −−− A gatekeeper changes the crossing terms.
  ℎfill \textit{All passage costs +1 Measure until the next landmark. Diplomacy requires a successful Presence + Sway (DV 4).} `[AUTHORITY]'
\end{enumerate}

\paragraph*{(Cache/Pass/Favor)} Cache token for hidden food or fuel stash; right‑of‑way pass recognized by a named trail; warden's favor with an escort letter.

\section*{Diamonds −−− Rewards/Leverage (Hex Discoveries)}
łabel{sec:wilderness−rewards}
\begin{enumerate}
\setcounter{enumi}{1}
ℹtem \textbf{Cache Token} −−− Key to a hidden food or fuel stash (once); preparation as power.
  ℎfill \textit{The token is a knucklebone. The stash is marked with a cairn that has one stone missing. The missing stone is the token.}
ℹtem \textbf{Right−of−Pass} −−− Recognized marker for a named trail or crossing.
  ℎfill \textit{The marker is a notch carved in a milestone. The notch is shaped like a bird. The bird is an owl. The owl watches at dusk.}
ℹtem \textbf{Warden's Favor} −−− Escort letter; ``they're with us'' as a shield.
  ℎfill \textit{The letter is stamped with a bell. The bell's clapper is missing. The absence of sound is the signature.}
ℹtem \textbf{Weather Window} −−− Good forecast and a narrow gate to use it.
  ℎfill \textit{The window is a single hour at dawn. The birds will tell you when. They will stop singing. That is the signal.}
ℹtem \textbf{Water Deed} −−− Lawful draw at a scarce source; necessity as a right.
  ℎfill \textit{The deed is a clay cup. Break it, and the spring will flow—once. The cup's shards can be traded for another deed. There are nine shards.}
ℹtem \textbf{Route Song} −−− A map scrap that actually works; knowledge as navigation.
  ℎfill \textit{The song has eight verses. The ninth is humming. The humming is the sound of the road itself. It changes with the weather.}
ℹtem \textbf{Remount Hire} −−− Fresh legs or hulls waiting at a post; speed as a service.
  ℎfill \textit{The remounts have no names. They answer to whistles and bribes. They remember every rider. They do not forgive cruelty.}
ℹtem \textbf{Truce Cord} −−− Taboo exemption at a site (one scene); peace as a thread.
  ℎfill \textit{The cord is braided from wool and ash. Burn it to invoke the truce. The ash must be from the same fire that warmed the truce.}
ℹtem \textbf{Toll Waiver} −−− A ferry, bridge, or reef gate honors this chit; passage as paper.
  ℎfill \textit{The chit is a dried leaf. The toll‑keeper will not touch it. They will nod and wave you through. They will not speak.}
ℹtem[J] \textbf{Favor Owed} −−− Locals owe you (or you owe them); trade for labor or insight.
  ℎfill \textit{The favor is carved on a tally stick. Break it when the favor is done. The stick will not burn. It will only splinter.}
ℹtem[Q] \textbf{Private Audience} −−− A keeper, shrine‑warden, or spirit of the place hears you alone.
  ℎfill \textit{The audience is held at midnight. The spirit wears your face. It asks only one question. You will answer truthfully.}
ℹtem[K] \textbf{Road Commission} −−− Temporary authority over a stretch of wilds.
  ℎfill \textit{The commission is a whistle. Blow it, and the land will listen—once. The whistle is carved from a bone that was never alive.}
ℹtem[A] \textbf{Earth's Exception} −−− One temporary rule‑bend (cross during a storm, pass uncounted, beasts ignore you).
  ℎfill \textit{The exception is a secret. The land will forget it as soon as you use it. You will remember it every night for a year.}
\end{enumerate}

\section*{Quick use notes}
łabel{sec:wilderness−quick−use}
\begin{itemize}
ℹtem Draw until you have all four suits: Spade = place, Heart = actor, Club = pressure, Diamond = leverage. Highest rank sets the main Timer (2–5 → 4, 6–10 → 6, J/Q/K → 8, A → 10).
ℹtem Diamonds are codified outcomes (cache/pass/favor) that change position rather than call for a roll.
ℹtem If any A appears, add a lingering omen of the land (a smell on the wind, a sound that carries too far, tracks that shift) that you can echo in later scenes.
\end{itemize}

\subsection*{Travel at a Glance (Fast & Loose)}
łabel{sec:wilderness−travel}
A journey through the Wilds is a series of choices, not a checklist. Before each leg (a day, a stretch of road, a single hex), the party picks two priorities from \textbf{Speed}, \textbf{Stealth}, \textbf{Safety}, \textbf{Survey}. Those become your edge; the others become vulnerabilities. The GM then frames a single roll based on the terrain and the weather’s mood (see \emph{Weather Whims}, below). Success means you find what you sought; a mixed result trades progress for a cost; a miss throws a complication from the \emph{Clubs} table onto the path.

\paragraph{Weather Whims (d6, one per leg)}
\begin{enumerate}
ℹtem \emph{Glass‑clear} – long sight, easy navigation. +1 to Survey actions.
ℹtem \emph{Low fog} – stealth advantage, speed suffers. +1 to Stealth, –1 to Speed.
ℹtem \emph{Trickle rain} – tracks wash, cold bites. First navigation or tracking roll at –1 die.
ℹtem \emph{Dry cold / heat} – resources drain faster. Mark +1 supply strain at end of leg.
ℹtem \emph{Shifting wind} – smoke and scent lie. Reroll one die on a weather‑related check, but the GM may bank 1 SB.
ℹtem \emph{Omen sky} (aurora, sun‑dog, three moons) – the land is listening. The GM may ask one player to answer a single question about the party’s future.
\end{enumerate}

\paragraph{Making Camp}
When you stop, choose one \emph{boon} and accept one \emph{risk}:
\begin{itemize}
  ℹtem \textbf{Boon:} heal a piece of gear; clear 1 Fatigue; scout tomorrow’s weather; share a story with a road spirit (gain a \emph{Truce Cord}).
  ℹtem \textbf{Risk:} thin shelter (next weather hits harder); smoky fire (attracts prowlers); exposed approach (ambush more likely); hungry neighbors (locals or beasts appear).
\end{itemize}
A road spirit’s favour requires a small offering—a pinch of salt, a dropped coin, a whispered truth. Be polite.

\subsection*{Lore Echoes (Cross‑Expansion Hooks)}
\begin{itemize}
  ℹtem \textbf{The Bitter Root:} Daughters of the Cord use certain shelter hollows as waystations. A renegade Sister has hidden a Hollowed in a ruined outpost; the child has not moved in weeks.
  ℹtem \textbf{The Threshold Compendium:} Old Roads are Class II Anchor sites—ghosts of imperial surveyors still measure the miles. Saikou Ira would leave a silver coin at every milestone.
  ℹtem \textbf{Malachai, the Chained Angel:} A silver coin that never spends lies at the bottom of every spring. Take it, and the water will never run dry—but the spring will remember your face.
  ℹtem \textbf{The Grey Wanderer's Grimoire:} Wilds psions are called ``Stone‑Listeners.'' They hear the memory of old trails. Claimant chiefs pay well for their services—and dispose of those who hear too much.
\end{itemize}

\subsection*{Wilderness Superstitions (burn for +1 die once)}
\begin{itemize}
  ℹtem \textbf{Bread to the Cairn:} Crumble a crust at a boundary stone. Ignore the first \emph{Bad Ground} penalty.
  ℹtem \textbf{Knot the Trail:} Tie a single knot in a piece of grass at a crossing. Cancel \emph{Doubleback} or take –1 die on your next navigation roll (your choice).
  ℹtem \textbf{Name to the Spring:} Whisper a dead traveler's name into a water source. Gain +1 Effect when foraging this scene, but mark \emph{Obligation} to the land.
\end{itemize}

\subsection*{Weather Omens (d10) −−− New Mini−Generator}
łabel{sec:wilderness−weather−omens}
When the sky speaks in signs, roll to interpret its warning or promise.

\begin{longtable}{|p{1.5cm}|p{10cm}|}
ℎline
\rowcolor{gray!20}
\textbf{d10} & \textbf{Omen Meaning} \\
ℎline
1 & Three crows fly west—the next storm will carry voices from the past. \\
ℎline
2 & Moonring broken at midnight—the weather will shift twice before dawn. \\
ℎline
3 & Wind carries iron scent on clear air—storm approaches from the direction of old battlefields. \\
ℎline
4 & Stars vanish in pattern—the land is holding its breath; the next step decides fate. \\
ℎline
5 & Sun dogs flank the rising sun—blessing on journeys begun before noon. \\
ℎline
6 & Fog forms perfect circles—what is lost will return, but changed. \\
ℎline
7 & Lightning strikes twice in same spot—the Hollow marks this place for attention. \\
ℎline
8 & Rain falls upward for three breaths—a reversal is coming; prepare to give what you take. \\
ℎline
9 & Wind dies completely at noon—the land sleeps; move swiftly or become part of its dream. \\
ℎline
10 & Aurora pulses with your heartbeat—the sky knows your name and your burden. \\
ℎline
\end{longtable}

\subsection*{Ruin Findings (d8) −−− New Mini−Generator}
łabel{sec:wilderness−ruin−findings}
When you explore a forgotten structure, roll to discover what echoes remain.

\begin{longtable}{|p{1.5cm}|p{10cm}|}
ℎline
\rowcolor{gray!20}
\textbf{d8} & \textbf{Discovery} \\
ℎline
1 & A mosaic floor depicting a migration no map remembers—follow the tiles to find a hidden valley. \\
ℎline
2 & Wind whistles through broken arches in a tune that feels like home—hum it to gain +1 Effect on the next social roll. \\
ℎline
3 & Stone faces weep clear liquid that tastes of the first rain of spring—drink to clear 1 Fatigue. \\
ℎline
4 & A single tile is warm to the touch—press it to reveal a compartment holding a perfectly preserved seed. \\
ℎline
5 & Shadows move opposite the light—stand in their intersection to become invisible for one exchange. \\
ℎline
6 & Carvings show hands exchanging not goods but breaths—press your palm to the wall to share a memory. \\
ℎline
7 & The air hums at a frequency that makes teeth ache—sing against it to gain +1 die on the next magical roll. \\
ℎline
8 & A doorway frame is grown over with living stone—walk through it to leave one burden behind. \\
ℎline
\end{longtable}

\subsection*{Trail Signs (d6) −−− New Mini−Generator}
łabel{sec:wilderness−trail−signs}
When you pause to read the land’s subtle messages, roll to interpret what the way is telling you.

\begin{longtable}{|p{1.5cm}|p{10cm}|}
ℎline
\rowcolor{gray!20}
\textbf{d6} & \textbf{Sign Meaning} \\
ℎline
1 & Pebbles point not to the path but to a patch of edible mushrooms—the land offers sustenance if you look aside. \\
ℎline
2 & Grass grows in perfect spirals around a stone—walk the spiral three times to find your bearings if lost. \\
ℎline
3 & A single feather lies across the trail—it belongs to a bird that nests only in places of perfect safety. \\
ℎline
4 & Moss grows thickest on the north face of rocks—but here it thrives on the south; the land is confused or changed. \\
ℎline
5 & Twigs are arranged in a pattern that matches your boot tread—something has been watching your steps. \\
ℎline
6 & Water in a puddle holds a perfect reflection… of a place you’ve never been. \\
ℎline
\end{longtable}

\begin{tcolorbox}[colback=black!3,colframe=black!40!white,title={Patronage & Power}]
In the Wilds, power flows through knowledge of the land, the ability to navigate without formal authority, and the careful cultivation of local relationships. True authority comes from understanding weather patterns, reading tracks, and maintaining the favor of both human inhabitants and the spirits of place. Those who can provide safe passage, find water in dry times, or mediate between conflicting claims hold the most influence.

\textbf{For the GM:}  
Patronage in the Wilds revolves around practical knowledge, local customs, and the ability to provide essential services like navigation, shelter, and supply. Rewards often take the form of tokens, passes, and local knowledge that can be leveraged into greater safety and efficiency. To emphasise the hexcrawl feel:
\begin{itemize}
ℹtem Key 6–10 interesting hexes before the session, then use the generator to fill the rest.
ℹtem Let the party discover new routes, short cuts, and hidden lairs through successful scouting or lucky finds.
ℹtem Use the \emph{Roads Move}, \emph{Weather Wakes}, and \emph{Local Resentment} timers (from the hexcrawl procedures) to create a living, reactive wilderness—or simply let the cards speak.
ℹtem Offer Diamonds as ``exploration rewards'' that unlock new options (a cache token, a safe passage, a local ally).
\end{itemize}
In the Wilds, your skill is your standing, and your standing determines whether you survive or perish.
\end{tcolorbox}

\subsection*{Terrain Stock (d20)}
łabel{sec:wilderness−terrain−stock}
Roll or choose a striking detail to dress a hex.
\begin{enumerate}
  ℹtem Fresh ash on an old fire ring; prints circle \emph{around} it.
  ℹtem Boundary totems toppled but re‑stack themselves if watched politely.
  ℹtem Broken wagon with salt‑vein crystals growing through spokes.
  ℹtem Singing culvert; echo names the boldest listener wrongly.
  ℹtem Weather bell hung in a pine; tolling swaps rain and wind.
  ℹtem Cairn with a copper coin balanced on edge; don't breathe hard.
  ℹtem Fence of antlers pointing uphill; follow and skip one \emph{Bad Ground}.
  ℹtem Mummified pack‑beast wearing a courier's braid; letter hums.
  ℹtem Pebble mosaic of an old levy call; stepping it summons wardens.
  ℹtem Fossil rib arch; shelter under it and dream of tomorrow's route.
  ℹtem Abandoned hunter's cache; choice of food or rumor, not both.
  ℹtem Mirror‑puddle; shows last night's sky whatever the hour.
  ℹtem Rope bridge made of prayer knots; one is freshly cut.
  ℹtem Scare‑stakes dressed in saltcloth; something big respects them.
  ℹtem Stone with bite‑marks; taste iron and remember a shortcut.
  ℹtem Old road mile‑post defaced into a saint; offerings still fresh.
  ℹtem Thistle patch buzzing with coin‑flies; feed them copper for safe passage.
  ℹtem Fallen giant's footprint now a pond; frogs chant a route‑song.
  ℹtem Soot writing on cliff: \emph{Storm at Second Bell}. It's right.
  ℹtem Hollow tree full of carved names; the newest is yours.
\end{enumerate}

\subsection*{Lair Seeds (One‑Page Crawls)}
łabel{sec:wilderness−lair−seeds}
d12 quick lair hooks. Each can be a single session’s focus.
\begin{enumerate}
  ℹtem \textbf{Fen Clockworks} — Bog‑gears turn under peat; each tooth a drowned oath.
  ℹtem \textbf{Dry River Gallery} — Petroglyphs re‑arrange at dusk; wrong reading summons bramble knights.
  ℹtem \textbf{Ice−Well Vault} — Stairs drilled through permafrost; breath blooms frost‑runes that answer questions badly.
  ℹtem \textbf{Sunken Forum} — Pillars knee‑deep in silt; rust moths nest in the rostrum.
  ℹtem \textbf{Pilgrim Kilns} — Charcoal mounds hike on stilt roots at night; a hermit‑healer bargains weather for charcoal.
  ℹtem \textbf{Hanging Salt Chapel} — White stalactites ring like bells; grave‑eels coil in the font.
  ℹtem \textbf{Thorn Maze Mill} — Waterwheel turns without water; phase panthers pace the sluice.
  ℹtem \textbf{Brewer's Hollow} — Barrels of wind; decant a gale or drink a calm.
  ℹtem \textbf{Antler Archive} — Tally carved on horn; theft awakens the cairn‑wights' accountant.
  ℹtem \textbf{Black Lantern Mines} — Lantern‑eyes nest like bats; light buys silence, darkness buys passage.
  ℹtem \textbf{Stormbone Spire} — Vertebrae tower; earth‑sharks circle beneath like sharks in sand.
  ℹtem \textbf{Amber Orchard} — Trees fossilised mid‑sway; trapped thunderflies grant \emph{Weather Window} when freed.
\end{enumerate}

\subsection*{Wilderness Bestiary}
łabel{sec:wilderness−bestiary}
A palette of strange creatures, each with \emph{signs} (what you notice before it sees you) and a \emph{hook} (how to engage or avoid).
\begin{itemize}
  ℹtem \textbf{Earth−Sharks} — Burrow beneath loess and tundra; dorsal ridges like plowshares. \emph{Signs:} furrows that breathe; pebbles rattle. \emph{Hook:} lure with drumbeats or goat bells.
  ℹtem \textbf{Phase Panthers} — Shadow‑striped cats that step between gusts. \emph{Signs:} pawprints offset from themselves. \emph{Hook:} mirrored water or bell‑silver disrupts the blink.
  ℹtem \textbf{Rust Moths} — Swarms that taste iron memory. \emph{Signs:} lacework mail; pitted nails. \emph{Hook:} feed them old keys; they ignore fresh steel.
  ℹtem \textbf{Lichen−Bears} — Moss‑matted ursines with rock‑teeth. \emph{Signs:} clawed bark with green scabs. \emph{Hook:} smoke of juniper calms; salt enrages.
  ℹtem \textbf{Wind Drakes} — Rib‑thin sky‑serpents that ride fronts. \emph{Signs:} shed vanes; feathers like glass. \emph{Hook:} whistle back the storm to ground them.
  ℹtem \textbf{Ooze−Remnants} — Alchemical slough from an old field lab. \emph{Signs:} boot‑holes rounded smooth. \emph{Hook:} wood ash firms them; vinegar makes them sprint.
  ℹtem \textbf{Fell Harts} — Antlered shades that herd the living toward cliffs. \emph{Signs:} hoofprints that don't compress soil. \emph{Hook:} speak a hunter's apology while walking backward.
  ℹtem \textbf{Cairn−Wights} — Stone‑stack custodians; patient as frost. \emph{Signs:} pebbles rearranged around your fire. \emph{Hook:} rebuild the stack correctly; they serve until dawn.
  ℹtem \textbf{Thundertusks} — Boars that store lightning in tusk‑veins. \emph{Signs:} fused sand beads; singed roots. \emph{Hook:} ground your spear; throw shadow to draw the strike.
  ℹtem \textbf{Lantern−Eyes} — Floating orbs of watchlight rumor; jealous of secrets. \emph{Signs:} daylight in ravines at midnight. \emph{Hook:} offer a whispered truth; they part like reeds.
  ℹtem \textbf{Grave−Eels} — Burrow through barrow loam; taste names. \emph{Signs:} sinking turf over old wars. \emph{Hook:} sing lineage or throw false names to mislead.
  ℹtem \textbf{Bramble Knights} — Empty helms grown through with thorn; patrol old borders. \emph{Signs:} briar in boot‑grease. \emph{Hook:} present a hedge‑blessing; they escort instead of impale.
\end{itemize}

\begin{tcolorbox}[colback=black!3,colframe=black!40!white,title={Strange Finds (d10)}]
łabel{sec:wilderness−relic−find}
When you discover a hidden cache or an ancient hoard, roll or pick:
\begin{enumerate}
  ℹtem \textbf{Route−stones} that warm near true north.
  ℹtem \textbf{Weather−bell clapper} that tolls only for lightning.
  ℹtem \textbf{Trail−right knot} dyed in berry law; locals honour it once.
  ℹtem \textbf{Salt−vein pick} that cuts oozes like cloth.
  ℹtem \textbf{Glass feather}; break to call a wind drake for one pass.
  ℹtem \textbf{Tally leaf} waterproofed with eel‑fat; ink never runs, words drift.
  ℹtem \textbf{Thorn sigil} that bramble knights read as parley.
  ℹtem \textbf{Coin‑fly cocoon}; hatch it for one honest guide to the nearest water.
  ℹtem \textbf{Bone flute} that calls earth‑sharks to circle elsewhere.
  ℹtem \textbf{Lantern‑eye husk}; squeeze for daylight the length of a camp prayer.
\end{enumerate}
\end{tcolorbox}

\subsection*{Omens & Signs}
łabel{sec:wilderness−omens}
The land speaks. Use these when the party pauses or when the GM needs a portent.

\paragraph{Sky Signs (d8)}
\begin{enumerate}
  ℹtem sun‑dog
  ℹtem mare's tails
  ℹtem anvil cloud faces the wrong way
  ℹtem silent birds
  ℹtem far thunder without storm
  ℹtem moon with halo
  ℹtem ash scent on clear air
  ℹtem stars ``shiver''
\end{enumerate}

\paragraph{Ground Signs (d8)}
\begin{enumerate}
  ℹtem cairn with new pebble
  ℹtem boot prints that end at rock
  ℹtem butchered game left neat
  ℹtem snapped twig at head height
  ℹtem ash under wet moss
  ℹtem ward‑knot cut
  ℹtem trail beads on shrub
  ℹtem a single coin on a stone
\end{enumerate}

\subsection*{Thematic SB Spend Table}
łabel{sec:wilderness−sb}
\paragraph{Minor Complications (1 SB)}
\begin{itemize}
ℹtem \textbf{Exposure:} Your actions draw unwanted attention from \textbf{wardens or local hunters}.
ℹtem \textbf{Noise:} Sounds of your actions alert nearby \textbf{prowlers or caravan guards}.
ℹtem \textbf{Trace:} Evidence of your passage marks your route for \textbf{trackers or claimants}.
ℹtem \textbf{Delay:} A brief but meaningful setback costs you \textbf{time or favorable weather}.
ℹtem \textbf{Supply Strain:} Mark +1 segment on a relevant \textbf{resource timer}.
\end{itemize}

\paragraph{Moderate Setbacks (2 SB)}
\begin{itemize}
ℹtem \textbf{Alarm Raised:} \textbf{Guide or warden patrol} becomes aware and begins responding.
ℹtem \textbf{Position Lost:} You lose advantageous ground/cover/stealth due to \textbf{weather shift or route blockage}.
ℹtem \textbf{Foe Appears:} A \textbf{war−band or territorial beast} arrives on scene.
ℹtem \textbf{Gear Trouble:} A piece of equipment becomes \textbf{Compromised/Neglected}.
ℹtem \textbf{Lock/Barrier:} A simple obstacle now requires a test to overcome.
\end{itemize}

\paragraph{Serious Trouble (3 SB)}
\begin{itemize}
ℹtem \textbf{Reinforcements:} Additional \textbf{wardens, war−bands, or prowlers} arrive.
ℹtem \textbf{Key Gear Breaks:} A crucial tool/weapon becomes temporarily unusable.
ℹtem \textbf{Major Twist:} The situation fundamentally changes − \textbf{weather turns/catastrophe strikes/quarantine declared}.
ℹtem \textbf{Rail Tick:} Advance a relevant campaign/front timer by 1 segment.
ℹtem \textbf{Condition Applied:} Mark \textbf{Fatigue 1/Harm 1/Condition} appropriate to fiction.
\end{itemize}

\paragraph{Major Turns (4+ SB)}
\begin{itemize}
ℹtem \textbf{Trap Springs:} A prepared danger activates with full effect.
ℹtem \textbf{Authority Arrival:} \textbf{Claimant chief, quartermaster, or The Stranger} intervenes.
ℹtem \textbf{Scene Shift:} The environment changes dramatically − \textbf{weather shifts/blockage clears/catastrophe hits}.
ℹtem \textbf{Patron Omen:} Divine/arcane forces take notice − \textbf{omen appears/blessing lost/curse manifests}.
ℹtem \textbf{Narrative Pivot:} The story takes an unexpected turn that reframes objectives.
\end{itemize}

\paragraph{Region−Specific SB Options}
\begin{itemize}
ℹtem \textbf{Wilds (Weather):} Conditions shift without warning, storms arrive early, calm turns to chaos.
ℹtem \textbf{Wilds (Navigation):} Trails reroute mid‑journey, landmarks disappear, guides demand additional payment.
ℹtem \textbf{Wilds (Local Customs):} Spirits demand tribute, territorial beasts become aggressive, local laws change with the terrain.
\end{itemize}

\subsection*{Fast reskin palette}
łabel{sec:wilderness−reskin−palette}
\begin{description}
ℹtem[Forest] swap dunes$→$deadfall, spindrift$→$canopy flare, prowlers$→$boar/wolves; water = spring/stream.
ℹtem[Desert] swap bog$→$salt pan, fog$→$dust, shelter = overhang/wadi; water = seep/fog net.
ℹtem[Tundra/Ice] swap dune shift$→$ice heave, fire$→$spindrift, boats$→$sleds; prowlers = bears/wolves.
ℹtem[Coast/Isles] crossings = reef shelves, prowlers = seals/raiders, alarms = harbor booms; fuel = driftwood.
ℹtem[Swamp/Fen] crossings = corduroy/log causeways, bad ground = peat crust, prowlers = gators/leeches.
ℹtem[Highlands] crossings = cols and scree traverses, alarms = beacon chains, prowlers = cats/eagles.
\end{description}

% Index entries
∈dex{Wilderness|see{Roads, Ruins, and Weather}}
∈dex{Hexcrawl}
∈dex{Trail rights}
∈dex{Weather whims}
∈dex{Supply and forage}
∈dex{Camps and watches}
∈dex{Road spirits}
∈dex{Claimant chief}
∈dex{Forager child}
∈dex{Cache token}
∈dex{Right−of−pass}
∈dex{Truce cord}
∈dex{Earth's Exception}
∈dex{Biome reskin}
∈dex{Isolated towns}
∈dex{Refugee train}
∈dex{Deserter camp}
∈dex{Beastman hunting band}
∈dex{Lost tribe}
∈dex{Rootless mercenaries}
∈dex{Feral children}
∈dex{Cult of the Unburied}
∈dex{The Grey Singer}
∈dex{The Hollow's Own}
∈dex{Expanded travel hazards}
∈dex{Bell−Line Failure}
∈dex{Star−Road Phasing}
∈dex{Spirit−Dust Storm}
∈dex{Ninth Night}
∈dex{Well−Spirit Honored}
∈dex{Toll Shift}
\newpage

% ------------------------------
% CHAPTER: DUNGEONS
% ------------------------------
\chapter{Dungeons}
\emph{Living Infrastructure.}

\noindent
{\Large
\textit{Beneath the surface, ancient places breathe with mechanical lungs and dream electric dreams. The walls remember every footfall; the air carries whispers of forgotten conversations. The Edgewalker believes dungeons are infrastructure. The prisoners believe they are hungry. Both are correct.}

\vspace{0.5em}

I am a Scholar-Prisoner. The walls here have a memory like old parchment---they remember every scream, every spell, every drop of blood spilled on the stones. And sometimes, when the air's just right, they whisper it back to you. The dungeon listens. The question is whether it learns. I have been here for nine years. I have heard it learning. It started with single words. Then sentences. Now it speaks in paragraphs. Last week, it finished my thoughts before I did.

\vspace{0.5em}

Dungeons are not ruins---they are systems. Ventilation shafts that breathe, water channels that pulse, doors that open and close without any hand to move them. The Aeler built many of them; the Covenant sealed others; the Sidhi mined a few; and some were never built at all. They grew. Like fungus in the dark. Like a story told too many times. Like a debt that compounds.

\vspace{0.5em}

Never ignore a door that opens when you're not looking. The dungeon is offering you a choice. Refuse, and the next door may not open at all. Accept, and the door may lead to a room that was always there, waiting for someone to finally walk through. The room has been waiting for nine thousand years. It is not tired of waiting.

\vspace{1em}

\noindent\textit{--- The Dungeon's Architect, whose plans are written in a language no one speaks and whose signature is a thumbprint that could be anyone's}
}

\vspace{1em}

\begin{almanac}
\textbf{What do people eat for breakfast?} Fungus bread and cave water---the fungus grows on the walls, the water drips from the ceiling. Both taste of stone. Both are all there is.

\textbf{What does a child learn before they can walk?} To count footsteps. Dungeon children are taught the eight sounds of the deep: Drip, Creak, Groan, Whisper, Scrape, Echo, Silence, and the ninth---the sound of the dungeon's own heartbeat. The heartbeat is a door opening. The door is always behind you.

\textbf{What is the first thing you notice about a stranger?} Their shadow. Dungeon shadows are long, even when the light is bright. A stranger whose shadow is too short has been here too long. A stranger whose shadow moves independently has been touched by the deep. A stranger with no shadow is not a stranger. They are the dungeon's guest. The dungeon never lets guests leave.

\textbf{What is the most common crime?} Map-making---but map-making is not a crime, it is a heresy. The dungeon does not want to be mapped. The dungeon changes its corridors when it feels watched. A map that is accurate at dawn is a lie by dusk. The lie is the map-maker's fault. The punishment is to walk the lie until it becomes true.

\textbf{What does no one talk about but everyone knows?} The dungeon is not a place. It is a creature---a sleeping thing that dreamed itself into stone. The corridors are its veins. The chambers are its organs. The prisoners are its thoughts. The thoughts are not its own.

\textbf{Where do people go when they die?} Into the walls. The dead are not buried. Their bones are ground into mortar and used to patch the cracks. The cracks always reappear. The dead are never finished being used.

\textbf{How do you apologize for a serious wrong?} You offer a memory---a true one, witnessed by the dungeon itself. The dungeon accepts the memory. The memory is added to the walls. The walls whisper it back to you when you sleep. The whisper is the apology. The apology never ends.

\textbf{What is the most important festival?} The Turning of the Keys (whenever the dungeon decides). The locks are changed. The prisoners are moved. The corridors are rearranged. The festival has no date, no warning, no end. It is happening now. You did not notice. That is the point.

\textbf{What is the secret that could destroy a powerful person?} The dungeon's architect is still alive. They are sealed in the lowest chamber, kept alive by the dungeon's own will. The dungeon does not want to be alone. The architect is the dungeon's only friend. The architect has not spoken in nine centuries. The dungeon does not need words.
\end{almanac}

\newpage
\section{Dungeons −−− Living Infrastructure}
łabel{chap:dungeons}

\subsection*{Quick Reference}

\textbf{Tagline:} \textit{Beneath the surface, ancient places breathe with mechanical lungs and dream electric dreams. The walls remember every footfall; the air carries whispers of forgotten conversations.}

\begin{quote}
  \textbf{The Dungeon's Architect (Elite):}  
  ``I am a Scholar−Prisoner. The walls here have a memory like old parchment—they remember every scream, every spell, every drop of blood spilled on the stones. And sometimes, when the air's just right, they whisper it back to you. The dungeon listens. The question is whether it learns. I have been here for nine years. I have heard it learning. It started with single words. Then sentences. Now it speaks in paragraphs. Last week, it finished my thoughts before I did.''
\end{quote}

\begin{quote}
  \textbf{Cavern−Crawler (Commoner):}  
  ``Dungeons are not ruins—they are systems. Ventilation shafts that breathe, water channels that pulse, doors that open and close without any hand to move them. The Aeler built many of them; the Covenant sealed others; the Sidhi mined a few; and some were never built at all. They grew. Like fungus in the dark. Like a story told too many times. Like a debt that compounds.''
\end{quote}

\noindent\textbf{Genre / Mood:} Survival horror underground, infrastructure as adversary, eco−horror, forgotten machinery.

\noindent\textbf{System Type:} Portable overlay — apply to any dungeon, ruin, or deep place. Not a fixed geography.

\vspace{2mm}
\begin{tcolorbox}[colback=black!3,colframe=blue!60!black,title={Theme & Atmosphere}]
Dungeons are not empty holes. They are \emph{living infrastructure}—semi‑sentient systems that remember, learn, and hunger. Their corridors shift when you aren't looking. Their doors open to places you did not intend. Their air tastes of copper, old keys, and the breath of something that has been waiting a very long time.

Beneath the surface, ancient places breathe with mechanical lungs. The Aeler built many of them as overflow vaults and under−roads; the Covenant sealed others to contain what should not walk; the Sidhi mined a few and left behind ghost−workings. Some were never built at all. They grew from pressure, from forgotten oaths, from the weight of unresolved debts. A dungeon is not a place. It is a process. Like a fungus in the dark. Like a story told too many times. Like a curse that compounds.

When you enter a dungeon, you are not entering a room. You are entering a \emph{system}—and systems have rules. Air must be counted. Light must be rationed. Noise is the first enemy. Attention is the last. The dungeon learns your patterns and adapts. It offers shortcuts to the desperate and dead ends to the confident. It remembers your name after you leave.

\vspace{2mm}
\textbf{What you notice first:} The air tastes of copper and old keys. Your footsteps echo half a beat too late. The walls are warm—and they pulse like hearts.

\textbf{The one rule that kills newcomers:} Never ignore a door that opens when you're not looking. The dungeon is offering you a choice. Refuse, and the next door may not open at all. Accept, and the door may lead to a room that was always there, waiting for someone to finally walk through. The room has been waiting for nine thousand years. It is not tired of waiting.
\end{tcolorbox}

\begin{tcolorbox}[colback=black!3,colframe=red!60!black,title={The Ninth Taboo Manifestation}]
In any dungeon, the ninth door you open (counting only doors that required conscious choice) will lead to a room that is \emph{not on your map}. The room is always warm. Something inside knows your name. If you enter, you may find treasure—or become part of the dungeon's memory. The wise open eight doors, then backtrack.
\end{tcolorbox}

\subsection*{Quickstart (Two Minutes)}
\begin{enumerate}[leftmargin=*]
    ℹtem \textbf{Choose a Dungeon Type:} Keystone Vault (Aeler), Sealed Catacomb (Covenant), Pressure Grotto (Mazereth), Spore Warren (Mycotheres), or Raw Vein (wild).
    ℹtem \textbf{Track Four Things:} Air [6], Light [6], Noise [4], Attention [6] (choose: Drake Notice, Cult Patrol, or Vein Unrest).
    ℹtem \textbf{Intent Dial (per room/zone):} Stealth / Speed / Survey / Communion (parley with place‑spirits).
    ℹtem \textbf{Diamonds to Carry:} keystone courtesy, bell−wire, toll−chalk, web token, shed‑skin writ, pressure map, reed charm.
    ℹtem \textbf{Etiquette on Entry:} Step on stone; speak under beam; ring once. (Grants Position +1 in first scene if observed.)
\end{enumerate}

\subsection*{Names of Dungeon Denizens}
Dungeon dwellers rarely use surnames. Epithets are earned from the living stone.

\begin{longtable}{|p{0.25\textwidth}|p{0.25\textwidth}|p{0.2\textwidth}|p{0.2\textwidth}|}
ℎline
\textbf{First Names (Male)} & \textbf{First Names (Female)} & \textbf{Clan/Origin} & \textbf{Epithets} \\
ℎline
Korr & Svala & Deepforged & \textit{the Unseeing} \\
ℎline
Thrain & Harken & Ironbound & \textit{the Candle−Eater} \\
ℎline
Grimm & Lira & Root−Touched & \textit{the Silent} \\
ℎline
Veth & Myrin & Spore−Kin & \textit{the Soot−Scarred} \\
ℎline
Rurik & Tamsin & Pressure−Born & \textit{the Echo−Thief} \\
ℎline
Cael & Nix & Hollow−Kin & \textit{the Last Step} \\
ℎline
\end{longtable}

\subsection*{Spades — Places (Dungeon Zones)}
\noindent\textit{Room, Corridor, Threshold, Heart}

\begin{enumerate}
\setcounter{enumi}{1}
ℹtem \textbf{Air−Sealed Chamber} −−− A room with no obvious vent; the air is still, warm, and tastes of iron. `[SURVIVAL]` \textit{Testing Air each scene; on a miss, start Suffocation [4].}
ℹtem \textbf{Whispering Gallery} −−− A long corridor where every sound echoes nine times. `[STEALTH]` \textit{Any roll that produces noise ticks Noise +1; on a miss, the echo speaks a name.}
ℹtem \textbf{Flooded Vault} −−− A submerged antechamber; pressure doors block the way. `[TRAVEL]` \textit{Crossing requires a Body+Athletics test (DV 4) or lose 1 Light (wet lamp).}
ℹtem \textbf{Clockwork Junction} −−− Gears turn in the walls; brass conduits hum. `[MAGIC]` \textit{A successful Tinkering test (DV 3) can shift one segment of any dungeon timer. On a 1, start Anomaly [4].}
ℹtem \textbf{Spore Garden} −−− Bioluminescent fungi carpet the floor; soft light, softer danger. `[HEAL]` \textit{Forage DV 2 (glow‑spores, medicinal moss). On a miss, inhale spores → Weariness +1.}
ℹtem \textbf{Bell−Line Corridor} −−− Ropes of copper bells hang from the ceiling; any movement rings them. `[STEALTH]` \textit{Stealth rolls start Desperate. A successful Tinkering test (DV 4) silences them for one scene.}
ℹtem \textbf{Keystone Arch} −−− A stone doorway carved with ancient Aeler runes. `[AUTHORITY]` \textit{Present a Keystone Courtesy or spend 1 Boon to open it without rolling.}
ℹtem \textbf{Mirror Hall} −−− Walls of polished obsidian show reflections that move independently. `[TRUTH]` \textit{Any lie spoken here is answered by a reflection that speaks the truth—at cost of 1 Fatigue to the liar.}
ℹtem \textbf{Pressure Vent} −−− A narrow shaft with rhythmic blasts of hot air. `[TRAVEL]` \textit{Crossing requires timing (Wits+Survival DV 4). On a miss, take Harm 1 (burn).}
ℹtem[J] \textbf{Hoard Chamber} −−− A cavern of treasure, but every coin is warm and whispers. `[WEIGHT]` \textit{Taking anything starts Hoard's Attention [4]; on fill, a guardian awakens.}
ℹtem[Q] \textbf{Ossuary Forge} −−− Bones and scrap metal are smelted into crude tools. `[CRAFT]` \textit{Can produce one Relic tag (e.g., BLEED, WARD) at cost of 1 Obligation (the dead remember).}
ℹtem[K] \textbf{Observatory of Scars} −−− A domed chamber with sky−windows showing past wars, not current skies. `[DIVINE]` \textit{Once per delve, ask the GM one question about the dungeon's history; the answer is true but cryptic.}
ℹtem[A] \textbf{The Heart} −−− The core of the dungeon: a beating crystal, a sleeping dragon, a coven of cultists, or a ticking machine. `[LEGEND]` \textit{The boss lives here. Entering triggers the final phase.}
\end{enumerate}

\subsection*{Hearts — People & Factions (Dungeon Denizens)}

\begin{enumerate}
\setcounter{enumi}{1}
ℹtem \textbf{Cavern−Crawler} −−− A grimy survivor with a lamp and a longer list of scars. `[SURVIVAL]` \textit{Knows one shortcut (reduce Distance by 1) but demands a share of the first Diamond found.}
ℹtem \textbf{Pressure Guilder (Mazereth)} −−− A deep tunneler who counts breaths like coin. `[TRADE]` \textit{Offers a Pressure Map (Diamond) for 1 Supply or a story about a collapse.}
ℹtem \textbf{Aeler Under−Mason} −−− A dwarf with a tuning fork and a level. `[CRAFT]` \textit{Can stabilize one collapsing zone per delve at cost of 1 Fatigue.}
ℹtem \textbf{Lethai−ar Silk−Warden} −−− A masked elf whose web−token grants passage through disputed corridors. `[SOCIAL]` \textit{Will escort the party past one hazard if they swear an oath witnessed by her.}
ℹtem \textbf{Spore−Speaker} −−− A mycothe priest who communes with fungal networks. `[MAGIC]` \textit{Can clear 1 segment of Air or Light in exchange for a memory (erase a minor fact).}
ℹtem \textbf{Drake−Touched Thrall} −−− A cultist with scales growing from their arms. `[FEAR]` \textit{Offers a Drake's Whistle (summon a minor drake for one scene) at cost of starting Attention +1.}
ℹtem \textbf{Iron Avenger} −−− An Aeler feud‑keeper hunting a fugitive in the deep. `[COMBAT]` \textit{Will fight beside you against a common enemy, but will demand a blood‑price afterward.}
ℹtem \textbf{Fossil Juror} −−− A sentient skeleton bound to judge intruders. `[JUSTICE]` \textit{Will not attack; instead, asks a riddle. Answer correctly to pass; fail and it demands a geas.}
ℹtem \textbf{Hoard−Wyrd} −−− A ghost bonded to a specific treasure. `[DEATH]` \textit{It will let you take the item if you promise to return it after one use—or suffer a curse.}
ℹtem[J] \textbf{Ember Imp} −−− A small, sharp‑toothed creature that trades heat for secrets. `[TRICK]` \textit{Will light your lamp for 1 Fatigue, but will also steal a small item when you aren't looking.}
ℹtem[Q] \textbf{Herald Gargoyle} −−− A stone sentinel that enforces the dungeon's etiquette. `[AUTHORITY]` \textit{Demands a tithe (a story, a breath, a name) before allowing passage.}
ℹtem[K] \textbf{Wyrm−Bound Knight} −−− A fallen knight bound by a geas to a sleeping dragon. `[OATH]` \textit{Will aid in a single combat if you promise to break their oath afterward.}
ℹtem[A] \textbf{The Dungeon's Voice} −−− A whisper from the walls, speaking in your own voice. `[MYSTERY]` \textit{It offers a shortcut or a treasure in exchange for a promise. Breaking it costs a memory.}
\end{enumerate}

\subsection*{Clubs — Complications/Threats (Dungeon Hazards)}
\noindent\textit{Environment, Attention, Betrayal, Unstable}

\begin{enumerate}
\setcounter{enumi}{1}
ℹtem \textbf{Micro−Collapse} −−− A ceiling cracks; stone dust fills the air. `[DISASTER]` \textit{Each PC tests Body+Athletics (DV 3) or takes Harm 1. Tick Noise +1.}
ℹtem \textbf{Lantern Sputter} −−− Your light source gasses or gutters. `[LIGHT]` \textit{Reduce Light by 1. On a miss, the lamp goes out entirely—start Stumble in Dark [4] timer.}
ℹtem \textbf{Pocket Gas} −−− A burst of foul air from a crack. `[AIR]` \textit{Tick Air +1. Anyone in the scene who fails a Endurance test (DV 3) marks 1 Fatigue.}
ℹtem \textbf{Echo−Lure} −−− The dungeon mimics a voice you know, calling from a side passage. `[TRAP]` \textit{Whoever follows must test Wits+Resolve (DV 4) or be separated from the party for 1 scene.}
ℹtem \textbf{Leech Bloom} −−− A patch of carnivorous fungi awakens. `[COMBAT]` \textit{Cap 2 adversary; if ignored, ticks Supplies –1 (they eat rations).}
ℹtem \textbf{Silk Strand Snaps} −−− A Lethai−ar web bridge collapses. `[TRAVEL]` \textit{Distance –1; anyone on the bridge tests Body+Athletics or takes Harm 1.}
ℹtem \textbf{Oath−Magnet} −−− A spoken promise in this zone becomes magically binding. `[OATH]` \textit{Any vow made here ticks a Promise Timer [4] that tightens on a lie.}
ℹtem \textbf{Drake Whisper} −−− A distant roar echoes; the dungeon's attention sharpens. `[ATTENTION]` \textit{Tick Attention +1; on a miss, also mark Fatigue (fear).}
ℹtem \textbf{Patrol Net Tightens} −−− Sentinels or cultists adjust their routes. `[AUTHORITY]` \textit{Next Stealth roll starts at Desperate Position.}
ℹtem[J] \textbf{Lose the Bell−Wire} −−− Your guide‑rope is cut; you are lost. `[TRAVEL]` \textit{Distance +2 segments; all navigation rolls suffer –1 die until you find a landmark.}
ℹtem[Q] \textbf{Task Loop} −−− The party repeats the same corridor; time loops. `[MAGIC]` \textit{Tick Weariness +1; on a miss, a character forgets why they came (lose one Clue).}
ℹtem[K] \textbf{Autonomous Spasm} −−− A clockwork mechanism runs wild. `[DISASTER]` \textit{A random piece of gear becomes Compromised; on a critical, it explodes (Harm 1 to wielder).}
ℹtem[A] \textbf{Convergence} −−− The dungeon's Attention fills. The boss manifests immediately. `[BOSS]` \textit{Treat as a Major Turn (4+ SB); roll on the Boss Seed table.}
\end{enumerate}

\subsection*{Diamonds — Treasures & Discoveries}
\noindent\textit{Each Diamond is a physical object with a cost and a story.}

\begin{enumerate}
\setcounter{enumi}{1}
ℹtem \textbf{Pressure Map} −−− A leather chart marked with safe routes. `[TRAVEL]` \textit{DV –1 to navigate collapses or gas pockets for one delve.}
ℹtem \textbf{Bell−Wire Spool} −−− A coil of copper wire that rings when taut. `[STEALTH]` \textit{Lay a witness line; convert one False Order/Echo Lure into a Confusion only.}
ℹtem \textbf{Keystone Courtesy} −−− A carved stone token. `[AUTHORITY]` \textit{Once per leg, raise or lower a portcullis ``by the book'' without a roll—if you can name the builder.}
ℹtem \textbf{Web−Token} −−− A silk charm from a Lethai−ar warden. `[SOCIAL]` \textit{Enter parley with any faction in the dungeon from a position of safety once.}
ℹtem \textbf{Shed−Skin Writ} −−− A parchment signed in Isoka's rite. `[LAW]` \textit{Voids one coercive clause (contract, geas, oath) at the cost of 1 Obligation (a future favour).}
ℹtem \textbf{Reed−Charm (under‑river)} −−− A water‑resistant talisman. `[TRAVEL]` \textit{Treat one blackwater reach as neutral terrain for a crossing.}
ℹtem \textbf{Cinder Pearl} −−− A drake's tooth still warm. `[MAGIC]` \textit{Add IGNITE to one small action (e.g., light a torch, spark a fuse) once.}
ℹtem \textbf{Gleam Reed} −−− A phosphorescent stem. `[LIGHT]` \textit{Remove one Root/Slow effect; also restores 1 Light segment.}
ℹtem \textbf{Cold Scale} −−− A frost‑wyrm scale. `[WARD]` \textit{Ignore one chill or terror effect for a scene; then crumbles.}
ℹtem[J] \textbf{Proof−Mote} −−− A floating mote of amber light. `[TRUTH]` \textit{Ask the GM: ``What here is misfiled?'' The answer is a clue about a hidden treasure or trap.}
ℹtem[Q] \textbf{Scale Draught} −−− A vial of diluted wyrm blood. `[FEAR]` \textit{Resist fear in a tribunal or horror scene; after use, tick Attention +1 (the wyrm notices).}
ℹtem[K] \textbf{Rime Crown} −−− A circlet of frost‑shaped bone. `[WARD] [COLD]` \textit{Once per session, gain COLD and WARD for a scene; then mark 1 Fatigue.}
ℹtem[A] \textbf{The Egg That Will Not Hatch} −−− A cold, glass‑smooth egg containing an idea. `[LEGEND]` \textit{Whisper a question; the egg answers in dreams. It will only hatch when someone dies willingly for it.}
\end{enumerate}

\subsection*{Dungeon Mechanics (Visible Tracks)}

Every delve tracks four resources. Keep them visible.

\begin{itemize}
    ℹtem \textbf{Air [6]} – Stale pockets, smoke, gas leaks. When Air $≥$ 4, all physical actions start –1 Position unless you have a bellows or mask.
    ℹtem \textbf{Light [6]} – Lamps, glow‑silk, bioluminescence. At Light = 0, all actions are Desperate; wandering monsters find you automatically.
    ℹtem \textbf{Noise [4]} – Clatter, echoes, spells. When Noise = 4, a patrol or predator arrives (GM chooses).
    ℹtem \textbf{Attention [6]} – The dungeon's focus. When Attention = 6, the boss manifests or a catastrophic trap triggers. Reset to 0 after.
\end{itemize}

\paragraph{Intent Dial (per zone)}
Choose one before entering a new zone:
\begin{itemize}
    ℹtem \textbf{Stealth} – Low Noise, slower. +1 Position to Hide Sign, but Distance advances slower.
    ℹtem \textbf{Speed} – Tick Noise to push distance; invites mishaps.
    ℹtem \textbf{Survey} – Read the stone/web/flow. On a hit, reduce Air or Light by 1.
    ℹtem \textbf{Communion} – Parley with place‑spirits. Risky boons (a clue, a shortcut) but may tick Attention.
\end{itemize}

\paragraph{Three−Timer Rule}
At most three active timers in a scene. The GM selects which of Air/Light/Noise/Attention are active. Long‑term timers (Hoard's Attention, Promise Timer ) are tracked between sessions.

\subsection*{Boss Seeds (Three‑Phase Foes)}
Roll d6 or choose a boss that fits the dungeon's heart.

\begin{enumerate}
    ℹtem \textbf{The Hoard−Wyrm (Crystalline Drake)} – \emph{Phases}: Greed's Gleam → Scale Horde → Awaken the Hoard. \emph{Cracks}: feed it a cursed treasure, speak a true regret, or shatter its heart‑crystal with a memory of generosity.
    ℹtem \textbf{The Unburnt Cult (Isoka Fanatics)} – \emph{Phases}: Shedding Rite → Trial of Venom → The Fiery Dawn. \emph{Cracks}: refuse the cup three times, expose a false priest, or force a Bowl (fairness) before a Fang (verdict).
    ℹtem \textbf{The Deep Drake (Pressure Below)} – \emph{Phases}: Whisper → Knock → Invitation → Collection. \emph{Cracks}: offer a memory of sunlight, speak a true regret, or allow a willing sacrifice to walk into the deep in your place.
    ℹtem \textbf{The Iron Avenger (Aeler Feud‑Keeper)} – \emph{Phases}: Blood Ledger → Toll Demanded → Reckoning. \emph{Cracks}: pay the wergild in public works, expose a false claimant, or complete the feud with a duel of witnesses (not blades).
    ℹtem \textbf{The Mycothe Sovereign (Spore Lord)} – \emph{Phases}: Bloom → Court of Mold → The Great Reseeding. \emph{Cracks}: burn the central mycelium, negotiate a spore‑truce (sacrifice 1 Supply of ash), or sing a counting song it has never heard.
    ℹtem \textbf{The Hollow Architect (Ghost of a Forgotten Builder)} – \emph{Phases}: Redraw the Map → Collapse the Approach → The Unfinished Stair. \emph{Cracks}: finish its blueprint (Craft test DV 5), return a stolen keystone, or convince it that the empire is truly dead.
\end{enumerate}

\subsection*{Loot Generator (Hoard Strata)}
When the party defeats the boss or cracks a major vault, roll on these tables to build a hoard. Use the same strata as the Dragon's Lair supplement.

\paragraph{Stratum A: Coin & Commodity (d12)}
\begin{enumerate}
    ℹtem Scales of electrum stamped with extinct dynasties.
    ℹtem River‑pearls that hum in rain. \texttt{[WATER]}
    ℹtem Salt‑bricks worth a city's winter.
    ℹtem Demon‑minted coppers—warm to touch. \texttt{[BLIGHT]}
    ℹtem Glass coins: spend with a lie, shatter on truth.
    ℹtem Star‑iron beads (forge‑grade). \texttt{[RELIC]}
    ℹtem Amber lumps with insects that whisper routes.
    ℹtem Scrip of a fallen bank; redeemable with the right story.
    ℹtem Temple tithes bound with blue twine (oath‑sealed).
    ℹtem Trade bars etched with river‑depth marks (pilot's tools).
    ℹtem Coral crowns; brittle but potent at sea. \texttt{[COMMAND]}
    ℹtem Black spice; anesthetic or poison by dose.
\end{enumerate}

\paragraph{Stratum B: Art & Relic (d12)}
\begin{enumerate}
    ℹtem Masks of the Seven Winds (sing for \texttt{WIND} once/scene).
    ℹtem Tapestry that updates current wars nightly.
    ℹtem Bone flutes that call extinct birds.
    ℹtem Chalice that refuses poisoned liquid.
    ℹtem Crown of Oaths: swearing while wearing binds magically. \texttt{[BIND]}
    ℹtem Twin mirrors: speak into one, echo from the other at midnight.
    ℹtem Zodiac astrolabe; points to eclipses. \texttt{[DIVINE]}
    ℹtem Ember harp; strings ignite when lies are told nearby.
    ℹtem Saint's gauntlet; unburnt by any fire.
    ℹtem Library sigil‑stamp; seals are obeyed by lesser courts.
    ℹtem Fossil codex; pages of shale can be read once each.
    ℹtem Lantern with bottled dawn (3 uses). \texttt{[LIGHT]}
\end{enumerate}

\paragraph{Stratum C: Names, Oaths, Intangibles (d12)}
\begin{enumerate}
    ℹtem A king's true name in a lead vial.
    ℹtem A bridge's right‑of‑way (you collect the toll).
    ℹtem Seven unspent apologies; each removes one curse.
    ℹtem A ship's luck bound in twine; cut to claim it.
    ℹtem The deed to a moonlit crossroads.
    ℹtem A treaty's missing clause; add it to rewrite borders.
    ℹtem The last lullaby of a nation; sing to calm armies.
    ℹtem A festival's first toast; begin it to unite feuds.
    ℹtem Weather's favor over one valley for a year.
    ℹtem The memory of a siege ladder; place it to open a gate.
    ℹtem A duel's outcome, never fought; declare to bind fate.
    ℹtem The debt of a cathedral to its mason's bloodline.
\end{enumerate}

\paragraph{Stratum D: Cursed Complications (d10)}
\begin{enumerate}
    ℹtem Taking coins awakens a Hoard−Wyrd.
    ℹtem The art objects are witnesses; they may testify in court.
    ℹtem Removing any crown asserts a claim; rivals appear.
    ℹtem Intangibles are tracked by Herald Gargoyles.
    ℹtem A rival dragon has a lien on half the hoard.
    ℹtem Curse of Counting: must tally treasure nightly or suffer Fatigue +1.
    ℹtem Hoard Scent: predators pursue the bearer.
    ℹtem Echo of Theft: the original owners dream of you.
    ℹtem Oath−Magnet: you attract sworn duels.
    ℹtem Tide−Tithe: waters reclaim 10% during each full moon.
\end{enumerate}

\subsection*{Quick Hoard Build}
\begin{enumerate}
    ℹtem Pick 1–2 from A, 1–2 from B, 1 from C, 1 from D.
    ℹtem Assign 1–2 \texttt{TAGS} as properties (e.g., \texttt{WARD}, \texttt{COMMAND}, \texttt{LIGHT}).
    ℹtem State a \emph{Hoard Law}: who may touch what (breaking it triggers a Lair Move or Attention tick).
\end{enumerate}

\subsection*{Dungeon Superstitions (burn for +1 die once)}
\begin{itemize}
    ℹtem \textbf{Bread to the Door:} Crumble a crust at a sealed threshold. Ignore the first \emph{Micro−Collapse} penalty. `[SUPERSTITION]`
    ℹtem \textbf{Knot the Bell−Wire:} Tie a single knot in a guide rope. Cancel \emph{Echo−Lure} or take –1 die on your next navigation roll (your choice). `[SUPERSTITION]`
    ℹtem \textbf{Name to the Wall:} Whisper a dead explorer's name into a crack. Gain +1 Effect when negotiating with the dungeon's Voice this scene, but mark \emph{Obligation} to the Architect. `[SUPERSTITION]`
    ℹtem \textbf{Salt to the Threshold:} Leave a pinch of salt at a door you do not wish to open again. The door will not open from either side for one day. `[SUPERSTITION]`
\end{itemize}

\subsection*{SB Menu (Dungeons)}
\paragraph{Minor (1 SB)}
\begin{itemize}
    ℹtem \textbf{Exposure:} A patrol or warden notices you.
    ℹtem \textbf{Noise:} Your footsteps echo; a nearby denizen stirs.
    ℹtem \textbf{Trace:} Scratched stone or spent lamp oil marks your route.
    ℹtem \textbf{Delay:} A door jams, a ladder breaks—lose time.
    ℹtem \textbf{Air Strain:} Tick Air +1 (stale pocket).
\end{itemize}

\paragraph{Moderate (2 SB)}
\begin{itemize}
    ℹtem \textbf{Alarm Raised:} A sentinel sounds a chime; Attention +1.
    ℹtem \textbf{Position Lost:} A floor collapses or a wall seals; you are separated.
    ℹtem \textbf{Foe Appears:} A minor drake, a fossil juror, or a cultist squad arrives.
    ℹtem \textbf{Gear Trouble:} A lamp sputters (Light –1) or a weapon snags.
    ℹtem \textbf{Lock/Barrier:} A keystone door now requires a test to open.
\end{itemize}

\paragraph{Serious (3 SB)}
\begin{itemize}
    ℹtem \textbf{Reinforcements:} A warband of Iron Avengers or a Hoard−Wyrd appears.
    ℹtem \textbf{Key Gear Breaks:} A bell‑wire snaps, a pressure map tears.
    ℹtem \textbf{Major Twist:} The dungeon shifts—a corridor loops, a new path opens.
    ℹtem \textbf{Rail Tick:} Advance a campaign timer (e.g., Hoard's Attention, Manifestation).
    ℹtem \textbf{Condition Applied:} Mark \textbf{Fatigue 1} or \textbf{Harm 1} from gas or fall.
\end{itemize}

\paragraph{Major (4+ SB)}
\begin{itemize}
    ℹtem \textbf{Trap Springs:} A prepared seal activates (Harm 2, Area).
    ℹtem \textbf{Authority Arrival:} The Dungeon's Voice, a Herald Gargoyle, or the boss interrupts.
    ℹtem \textbf{Scene Shift:} A collapse seals the exit; the Heart chamber opens prematurely.
    ℹtem \textbf{Patron Omen:} A draconic echo or a Lethai−ar omen appears.
    ℹtem \textbf{Narrative Pivot:} The dungeon reveals a forgotten truth—a secret door, a hidden geas, a broken oath.
\end{itemize}

\subsection*{Cross−Regional Use}
\begin{itemize}
    ℹtem \textbf{Aeler Holds:} A \textbf{Keystone Courtesy} can open sealed vaults; using it ticks \textbf{Ancestral Grudge} if the builder's lineage is insulted.
    ℹtem \textbf{Mistlands:} A \textbf{Bell−Wire Spool} can be used to rehang a broken bell−line; start \textbf{Witchlight Count} after.
    ℹtem \textbf{Ykrul Steppe:} A \textbf{Cold Scale} counts as a \emph{Salt Allotment} in a kurgan negotiation; the dead will demand a story.
    ℹtem \textbf{Sihai:} A \textbf{Proof−Mote} may be accepted as a \emph{Silence Token} in a Censorate audit; the auditor will ask where you found it.
    ℹtem \textbf{Valewood:} A \textbf{Rime Crown} can be offered to the Alder King as a winter tithe; he will ask for a memory of fire in return.
\end{itemize}

\begin{tcolorbox}[colback=black!3,colframe=black!40!white,title={Patronage & Power}]
In a dungeon, power flows through breath, light, and the careful management of the dungeon's attention. The true authorities are those who understand its systems—the Aeler under‑masons who know keystone law, the Lethai‑ar silk‑wardens who read web‑law, the Mazereth pressure guilders who count every exhalation. But the dungeon itself is the final patron. It offers shortcuts, treasures, and secrets—always at a price. Accept a shortcut, and it may close behind you. Take a treasure, and the Hoard's Attention ticks. Speak a true name, and the dungeon learns your voice.

\textbf{For the GM:}
\begin{itemize}
    ℹtem Use the four tracks (Air, Light, Noise, Attention) as visible pressure. Let players buy relief with Strings or etiquette, not just dice.
    ℹtem Enforce the \textbf{Intent Dial} each zone; switching intent costs a scene or a tick of Noise.
    ℹtem Offer one \textbf{Diamond} that could solve the current scene—but only if the party accepts a future Price (a geas, a curse, an obligation to a faction).
    ℹtem Spend SB as environment first (collapse, sputter, echo) before monsters. The dark should be the first antagonist, the denizens the second, the boss the third.
    ℹtem Use the \textbf{Boss Seeds} as capstones for major delves. A boss is not a random encounter; it is the dungeon's answer to the party's intrusion.
\end{itemize}
The under realms are not a pit to be cleared but a commons to be negotiated. Where lights are counted and breaths are tallied, procedure is mercy—and every crossing is a promise that must hold under weight.
\end{tcolorbox}

% Index entries
∈dex{Dungeons|see{Living Infrastructure}}
∈dex{Theme generator}
∈dex{Living Infrastructure generator}
∈dex{Places!Dungeons}
∈dex{People!Dungeons}
∈dex{Complications!Dungeons}
∈dex{Rewards!Dungeons}
∈dex{Air track}
∈dex{Light track}
∈dex{Noise track}
∈dex{Attention track}
∈dex{Boss Seed}
∈dex{Hoard generator}
∈dex{Keystone Courtesy}
∈dex{Bell−Wire}
∈dex{Pressure Map}
∈dex{The Dungeon's Voice}

\newpage

% ============================================================
% EPILOGUE: THE LEDGER CLOSES
% ============================================================
\chapter*{Epilogue: The Ledger Closes}
\addcontentsline{toc}{chapter}{Epilogue: The Ledger Closes}

\noindent
{\Large
I closed the book. The ink was dry.

\vspace{0.5em}

I had written the final entry months ago, in a damp counting-house on the Dhaharan coast, by candlelight. My hands did not shake---they were too good for that---but I set down my abacus very carefully. I had found something in the margins. A footnote I did not remember writing. A clause I did not remember adding. The ink was the same color as my own blood. The handwriting was my own. The date was tomorrow.

\vspace{0.5em}

I have learned to trust the things I do not remember. They are usually the truest things I know.

\vspace{0.5em}

The hat---the one the Trow had given me, the one that was still warm---sat on the corner of the desk. I picked it up, turned it over in my hands, and for the first time in years, I could not think of a single angle. That should have frightened me. Instead, it felt like a door opening. Not the kind of door you walk through. The kind of door that walks through you.

\vspace{0.5em}

The road was south. I did not know who would be waiting. I went anyway.
}

\vspace{1em}

\noindent\textit{--- The Edgewalker}

\vspace{1em}

% After the Edgewalker's final line, add this:

\vspace{2em}
\begin{quote}
\small
\textit{``The road is a ledger. Every oasis is a credit, every broken wheel a debt. The Tulkani tie knots to remember who owes whom. The Fhara weigh water before they weigh silver. The Pereshi recite verses to keep the bridges open. And the Sidhi keep lighthouse lamps burning for the dead who still have not reached the shore.''}

\noindent --- Dusana of the Raven Road, at a caravanserai whose name she refused to give
\end{quote}